%-----------------------------------------------------------------------------
% Master Prefix Router for Randomized Algorithms Class Notes & Book
%
% Dispatches to:
%   - prefix_lwarp.tex when warping/HTML mode is active (lwarpmk html or \warpMode)
%   - prefix_4ht.tex when TeX4ht/make4ht defines \HCode
%   - prefix_latex.tex for standard LaTeX / XeLaTeX / LuaLaTeX PDF builds
%-----------------------------------------------------------------------------

% Include only once - second time... igmore
\ifdefined\PrefixLoadedAlready
  \endinput
\fi

\def\PrefixLoadedAlready{}


\makeatletter
\def\input@path{{styles/}}
\makeatother


\makeatletter
\def\LWR@isWarp{0}%
\ifx\warpMode\undefined\else\def\LWR@isWarp{1}\fi
\expandafter\ifx\csname opt@lwarp.sty\endcsname\relax\else\def\LWR@isWarp{1}\fi

\if1\LWR@isWarp
  \def\warpMode{1}%
  \IfFileExists{styles/prefix_lwarp.tex}{%-----------------------------------------------------------------------------
% Master Prefix File for Lwarp (HTML Web Generation)
% Compatible with LuaLaTeX and Lwarp
%-----------------------------------------------------------------------------

\RequirePackage{iftex}
\ifLuaTeX
  % LuaLaTeX verified
\else
  \PackageError{lwarp-gate}{%
    prefix_lwarp requires LuaLaTeX! Please run with lualatex.%
  }{HTML generation via lwarp must be run with LuaLaTeX.}
  \stop
\fi

% Always use memoir
\documentclass[12pt,oneside,fleqn]{memoir}

% Load lwarp immediately after documentclass
\usepackage[mathjax]{lwarp}

% Lwarp settings: split at chapter level (0) with section names
\setcounter{FileDepth}{0}
\booltrue{FileSectionNames}
\setcounter{SideTOCDepth}{1}

\makeatletter
% Multi-page SVG resolving hook for Ipe drawings with page=N
\def\LWR@current@figpage{}
\define@key{igraph}{page}{\def\LWR@current@figpage{#1}}
% Perpage compatibility for auxiliary files
\providecommand{\pp@pagectr}[4]{}
\AtBeginDocument{
  \ifdef{\LWR@nullfonts}{%
    \appto\LWR@nullfonts{%
      \def\xspace{ }%
      \def\Turan{Turán}%
      \def\Lovasz{Lovász}%
      \def\Erdos{Erdős}%
      \def\Holder{Hölder}%
      \def\Matousek{Matoušek}%
    }%
  }{}%
  \ifdef{\LWR@replaceEPSSVG}{
    \appto\LWR@replaceEPSSVG{%
      \ifx\LWR@current@figpage\@empty\else
        \edef\LWR@tempone{\LWR@tempone_p\LWR@current@figpage}%
        \def\LWR@current@figpage{}%
      \fi
    }%
  }{}%
}
\makeatother

%-----------------------------------------------------------------------------
% Typography, Fonts & Math Packages
%-----------------------------------------------------------------------------
\usepackage{mathtools}
\PassOptionsToPackage{warnings-off={mathtools-overbracket,mathtools-colon}}{unicode-math}
\usepackage{fontspec}
\usepackage{unicode-math}

% New Computer Modern OpenType text and math fonts
\setmainfont{NewCM10-Book.otf}[
  BoldFont          = NewCM10-Bold.otf,
  ItalicFont        = NewCM10-BookItalic.otf,
  BoldItalicFont    = NewCM10-BoldItalic.otf,
  SlantedFont       = NewCM10-Book.otf,
  SlantedFeatures   = {FakeSlant = 0.25},
  BoldSlantedFont   = NewCM10-Bold.otf,
  BoldSlantedFeatures = {FakeSlant = 0.25},
  SmallCapsFeatures = {Letters = SmallCaps}
]
\setsansfont{NewCMSans10-Book.otf}[
  BoldFont          = NewCMSans10-Bold.otf,
  ItalicFont        = NewCMSans10-BookOblique.otf,
  BoldItalicFont    = NewCMSans10-BoldOblique.otf
]
\setmonofont{NewCMMono10-Book.otf}[
  BoldFont          = NewCMMono10-Bold.otf,
  ItalicFont        = NewCMMono10-BookItalic.otf,
  BoldItalicFont    = NewCMMono10-BoldOblique.otf
]

\setmathfont{NewCMMath-Book.otf}
\setmathfont{NewCMMath-Bold.otf}[version=bold]

% OpenType math aliases and fixes
\providecommand{\square}{\mdlgwhtsquare}
\def\Box{\mdlgwhtsquare}
\def\Diamond{\mdlgwhtdiamond}
\def\blacksquare{\mdlgblksquare}
\def\leadsto{\rightsquigarrow}
\providecommand{\llbracket}{\lBrack}
\providecommand{\rrbracket}{\rBrack}
\DeclareTextCommand{\textcircled}{TU}[1]{%
  \mbox{\ooalign{\hfil\raise0.1ex\hbox{#1}\hfil\crcr\char"25EF}}%
}

\usepackage{alphabeta}
\usepackage[normalem]{ulem} % Provides \sout
\usepackage{mleftright}
\usepackage{microtype}
\usepackage{xspace}
\usepackage{xurl}
\usepackage[dvipsnames,table]{xcolor}
\usepackage{graphicx}
\usepackage{caption}
\usepackage{booktabs,array,colortbl,hhline,multirow,tabularx}
\usepackage[inline]{enumitem}
\usepackage[thicklines]{cancel}
\renewcommand\CancelColor{\color{purple}}
\usepackage{euscript}
\usepackage{marvosym}
\let\Cross\undefined
\let\Square\undefined
\usepackage{bbding}
\usepackage{epigraph}
\usepackage{makeidx}
\usepackage{perpage}
\usepackage{tikz}

%-----------------------------------------------------------------------------
% Boondox / mathcalb font declaration
%-----------------------------------------------------------------------------
\DeclareFontFamily{U}{BOONDOX-calo}{\skewchar\font=45 }
\DeclareFontShape{U}{BOONDOX-calo}{m}{n}{
  <-> s*[1.05] BOONDOX-r-calo}{}
\DeclareFontShape{U}{BOONDOX-calo}{b}{n}{
  <-> s*[1.05] BOONDOX-b-calo}{}
\DeclareMathAlphabet{\mathcalb}{U}{BOONDOX-calo}{m}{n}
\SetMathAlphabet{\mathcalb}{bold}{U}{BOONDOX-calo}{b}{n}
\DeclareMathAlphabet{\mathbcalb}{U}{BOONDOX-calo}{b}{n}

%-----------------------------------------------------------------------------
% Algorithm Environment (formerly salgorithm.sty)
%-----------------------------------------------------------------------------
\newcommand{\ProgWordX}[1]{{\small\textbf{#1}}{\ }}
\newcommand{\Do}{\ProgWordX{do}\xspace}
\newcommand{\End}{\ProgWordX{end}\xspace}
\newcommand{\Return}{\ProgWordX{return}\xspace}%
\newcommand{\Procbegin}{\ProgWordX{begin}}
\newcommand{\If}{\ProgWordX{if}\xspace}
\newcommand{\Then}{\ProgWordX{then}\xspace}
\newcommand{\Else}{\ProgWordX{else}\xspace}
\newcommand{\For}{\ProgWordX{for}\xspace}
\newcommand{\While}{{\small\bfseries while}\ }

\newcommand{\xbeginlgox}{\begin{minipage}{1in}\begin{tabbing}
           \quad\=\qquad\=\qquad\=\qquad\=\qquad\=\qquad\=\qquad\=\kill}
\newcommand{\xendlgox}{\end{tabbing}\end{minipage}}

\newenvironment{algorithmEnv}{%
   \begin{tabular}{|l|}\hline\xbeginlgox\bigskip}%
   {\xendlgox\\\hline\end{tabular}}

\newenvironment{program}{%
   \begin{minipage}{4.0in}%
   \begin{tabbing}%
       \ \ \ \ \=\ \ \ \=\ \ \ \ \=\ \ \ \ \=\ \ \ \ \=%
      \ \ \ \ \=\ \ \ \ \=\ \ \ \ \=\ \ \ \ \=%
      \ \ \ \ \=\ \ \ \ \=\ \ \ \ \=\ \ \ \ \=\kill%
}{%
   \end{tabbing}%
   \end{minipage}%
}

\newcommand{\AlgorithmBox}[1]{%
   \fbox{\begin{program} #1 \end{program}}%
}

\definecolor{AlgorithmColor}{cmyk}{0.07,0.90,0,0.34}
\providecommand{\Algorithm}[1]{{\AlgorithmI{#1}\index{algorithm!#1@{\AlgorithmI{#1}}}}}
\newcommand{\CodeComment}[1]{\textcolor{blue}{\texttt{#1}}}

%-----------------------------------------------------------------------------
% Path & Directory Resolution
%-----------------------------------------------------------------------------
\newcommand{\RootDIR}{../}
\IfFileExists{../../../book.tex}{\renewcommand{\RootDIR}{../../../}}{}
\IfFileExists{../../book.tex}{\renewcommand{\RootDIR}{../../}}{}
\IfFileExists{../book.tex}{\renewcommand{\RootDIR}{../}}{}
\IfFileExists{./book.tex}{\renewcommand{\RootDIR}{}}{}

\providecommand{\File}[1]{#1}
\providecommand{\RootFile}[1]{\RootDIR #1}

%-----------------------------------------------------------------------------
% QED Duck Symbol
%-----------------------------------------------------------------------------
\newlength{\CharHeight}
\AtBeginDocument{\setlength{\CharHeight}{\fontcharht\font`X}}

\IfFileExists{styles/qed_duck.svg}{%
   \providecommand{\myqedsymbol}{%
      \reflectbox{\includegraphics[height=1.4\CharHeight]{styles/qed_duck}}%
   }%
}{%
   \providecommand{\myqedsymbol}{$\blacksquare$}%
}

%-----------------------------------------------------------------------------
% BibLaTeX Setup (Pure Biber, No BibTeX)
%-----------------------------------------------------------------------------
\usepackage[bibencoding=utf8,style=alphabetic,refsegment=chapter,
indexing=cite,
backend=biber]{biblatex}

\ExecuteBibliographyOptions{%
   maxnames = 40,%
   maxalphanames = 4,%
   maxcitenames = 3,%
   minalphanames = 3,%
   minnames = 1,%
   giveninits = true,%
   doi = false,%
   url = false,%
   isbn = false%
}

\renewbibmacro{in:}{}

\AtEveryBibitem{%
   \ifentrytype{inproceedings}{%
      \clearname{editor}%
      \clearlist{editor}%
      \clearfield{editor}%
      \clearlist{publisher}%
      \clearfield{series}%
      \clearlist{location}%
      \clearfield{publisher}%
   }{}%
   \ifentrytype{book}{%
      \clearname{publisher}%
   }{}%
}
\AtEveryCitekey{%
   \ifentrytype{inproceedings}{%
      \clearlist{publisher}%
   }{}%
   \clearname{editor}%
   \clearlist{location}%
}
\DeclareFieldFormat[inproceedings]{date}{}
\DeclareFieldFormat[inproceedings]{pages}{#1, \printfield{year}}
\DeclareFieldFormat[inproceedings]{booktitle}{\mkbibemph{#1,\nopunct}}

\DeclareFieldFormat{journaltitle}{\mkbibemph{#1}\isdot}

\DeclareFieldFormat[article]{number}{\mkbibparens{#1}}
\DeclareFieldFormat[article]{date}{}
\DeclareFieldFormat[article]{pages}{#1, \printfield{year}}
\DeclareFieldFormat[article]{issue}{\nopunct#1}
\renewbibmacro*{issue+date}{%
   \iffieldundef{issue}{}%
   {%
      \printtext{\unspace}%
      (\printfield{issue})%
   }%
   \iffieldundef{pages}%
   {\addcomma\printfield{year}}%
   {}%
}

\renewcommand{\bibpagespunct}{%
   \ifentrytype{article}{\addcolon}{\addcomma}\space%
}

\DeclareFieldFormat[article]{titlecase}{\MakeSentenceCase{#1}}
\DeclareFieldFormat[inproceedings]{titlecase}{\MakeSentenceCase{#1}}
\DeclareFieldFormat[incollection]{titlecase}{\MakeSentenceCase{#1}}
\DeclareFieldFormat[book]{titlecase}{\MakeSentenceCase{#1}}

\renewbibmacro*{journal}{%
   \iffieldundef{journaltitle} {} {\printtext[journaltitle]{%
         \printfield[myplain]{journaltitle}%
         \setunit{\subtitlepunct}%
         \printfield[myplain]{journalsubtitle}}}}
\DeclareFieldFormat{myplain}{#1}

\renewbibmacro*{booktitle}{%
   \ifboolexpr{ test {\iffieldundef{booktitle}} and test
      {\iffieldundef{booksubtitle}} } {} {\printtext[booktitle]{%
         \printfield[myplain]{booktitle}%
         \setunit{\subtitlepunct}%
         \printfield[myplain]{booksubtitle}}%
      \newunit}%
   \printfield{booktitleaddon}}

\renewbibmacro*{volume+number+eid}{%
   \setunit*{\addcomma\addspace}%
   \printfield{volume}%
   \printfield{number}%
   \setunit{\addcomma\space}%
   \printfield{eid}%
}

\newbibmacro{string+doi}[1]{%
   \iffieldundef{doi}{%
      \iffieldundef{url}{%
         #1%
      }{%
         \href{\thefield{url}}{#1}%
      }%
   }{%
      \href{http://dx.doi.org/\thefield{doi}}{#1}}%
}
\DeclareFieldFormat*{title}{\usebibmacro{string+doi}{\mkbibemph{#1}}}
\renewcommand*{\multicitedelim}{\addcomma\space}
\DeclareFieldFormat[inproceedings]{title}{\usebibmacro{string+doi}{{#1}}}
\DeclareFieldFormat[article]{title}{\usebibmacro{string+doi}{{#1}}}
\DeclareFieldFormat[incollection]{title}{\usebibmacro{string+doi}{{#1}}}
\DeclareFieldFormat[book]{titlecase}{{#1}}

\DeclareCiteCommand{\tcite}%
{\usebibmacro{prenote}}%
{\usebibmacro{citeindex}\usebibmacro{cite}}%
{\multicitedelim}%
{\usebibmacro{postnote}}

%-----------------------------------------------------------------------------
% Hyperref & Visual Formatting
%-----------------------------------------------------------------------------
\usepackage{hyperref}
\hypersetup{%
   breaklinks,%
   colorlinks=true,%
   linkcolor=[rgb]{0.45,0.0,0.0},%
   citecolor=[rgb]{0,0,0.45}%
}

\definecolor{DeepNavy}{HTML}{002366}
\definecolor{RoyalBlue}{HTML}{004E92}
\definecolor{SteelBlue}{HTML}{4682B4}
\definecolor{TealBlue}{HTML}{008080}
\definecolor{DarkSlate}{HTML}{2F4F4F}
\definecolor{DustyGrey}{HTML}{696969}

\makeatletter
\AtBeginDocument{\let\@elt\relax}
\makeatother

\headstyles{bringhurst}
\chapterstyle{dash}

\renewcommand{\chapnamefont}{\normalfont\huge\bfseries\color{DeepNavy}}
\renewcommand{\chapnumfont}{\normalfont\huge\bfseries\color{DeepNavy}}
\renewcommand{\chaptitlefont}{\normalfont\Huge\bfseries\color{DeepNavy}}

\counterwithin{footnote}{chapter}

\setsubparaheadstyle{\normalfont\normalsize\itshape\color{DustyGrey}}
\setlength{\beforechapskip}{-20pt} 
\setlength{\afterchapskip}{30pt}
\setsecnumformat{\csname the#1\endcsname.\quad}

%-----------------------------------------------------------------------------
% Standalone Chapter vs Master Book Modes
%-----------------------------------------------------------------------------
\newcommand{\Copyright}{This work is licensed under the Creative Commons Attribution-Noncommercial 3.0 License. To view a copy of this license, visit \url{http://creativecommons.org/licenses/by-nc/3.0/} or send a letter to Creative Commons, 171 Second Street, Suite 300, San Francisco, California, 94105, USA.}

\newcommand{\RealChapterEndInner}{%
   {%
      \begingroup%
      \emergencystretch=1em%
      \printbibliography[heading=subbibliography,segment=\therefsegment]{}%
      \endgroup%
   }%
}

\providecommand{\BookMode}[1]{}
\providecommand{\StandAloneMode}[1]{#1}

\ifx\bookMode\undefined
   % STANDALONE CHAPTER MODE
   \renewcommand{\StandAloneMode}[1]{#1}
   \renewcommand{\BookMode}[1]{}

   \newcommand{\ChapterNumPage}[2]{%
      \if\relax\detokenize{#1}\relax
         \setcounter{chapter}{0}%
      \else\ifnum#1>0
         \setcounter{chapter}{\numexpr#1-1\relax}%
      \else
         \setcounter{chapter}{0}%
      \fi\fi
      \if\relax\detokenize{#2}\relax
         \setcounter{page}{1}%
      \else\ifnum#2>0
         \setcounter{page}{#2}%
      \else
         \setcounter{page}{1}%
      \fi\fi
   }

   \DeclareRobustCommand{\setheadingcolor}{\color{RoyalBlue}}
   \setsecheadstyle{\setheadingcolor\normalfont\Large\bfseries}
   \setsubsecheadstyle{\normalfont\large\bfseries\color{SteelBlue}}
   \setsubsubsecheadstyle{\normalfont\normalsize\bfseries\color{TealBlue}}
   \setparaheadstyle{\normalfont\normalsize\bfseries\color{DarkSlate}}

   \DeclareDocumentCommand{\Chapter}{s o m}{%
      \IfBooleanTF{#1}{%
         \IfValueTF{#2}{\chapter*[]{#2}}{\chapter*[]{#3}}%
      }{%
         \IfValueTF{#2}{\chapter*[]{#2}}{\chapter[]{#3}}%
      }%
      \setcounter{footnote}{0}%
   }
   \newcommand{\ChapterPrefix}{}
   \newcommand{\ChapterGen}{\pagestyle{plain}\vspace*{-0.7cm}}
   \newcommand{\ChapterEnd}{%
      \RealChapterEndInner{}%
      \end{document}%
   }

   \renewcommand{\File}[1]{#1}
   \renewcommand{\RootFile}[1]{../#1}
\else
   % MASTER BOOK MODE
   \renewcommand{\StandAloneMode}[1]{}
   \renewcommand{\BookMode}[1]{#1}

   \newcommand{\ChapterNumPage}[2]{}

   \DeclareDocumentCommand{\Chapter}{s o m}{%
      \IfBooleanTF{#1}{%
         \IfValueTF{#2}{\chapter*[#2]{#3}}{\chapter*{#3}}%
         \ifx\Dir\undefined\else
            \typeout{CHAPTER_MAP: \Dir/\CurrentChapFile => chap: 0, page: \thepage}%
         \fi
      }{%
         \IfValueTF{#2}{\chapter[#2]{#3}}{\chapter{#3}}%
         \ifx\Dir\undefined\else
            \typeout{CHAPTER_MAP: \Dir/\CurrentChapFile => chap: \thechapter, page: \thepage}%
         \fi
      }%
   }
   \newcommand{\ChapterPrefix}{%
      \vspace{-1cm}%
      \subsection*{\small
         {598 - Class notes for Randomized Algorithms}\\
         {Sariel Har-Peled}\\
         {\today}}%
   }
   \newcommand{\ChapterGen}{}
   \newcommand{\ChapterEnd}{\RealChapterEndInner{}}

   \renewcommand{\File}[1]{\Dir/#1}
   \renewcommand{\RootFile}[1]{#1}

   \newcommand{\InputChap}[2]{%
      \def\Dir{#1}%
      \def\CurrentChapFile{#2}%
      \typeout{}%
      \typeout{------------------------------------------------------------}%
      \typeout{PROCESSING FILE: #1/#2}%
      \typeout{}%
      \input{#1/#2}%
   }
\fi

%-----------------------------------------------------------------------------
% Fragments & Quotes
%-----------------------------------------------------------------------------
\newcommand{\InputExt}[1]{%
   \IfFileExists{#1}%
   {\input{#1}}%
   {\IfFileExists{../#1}%
      {\input{../#1}}%
      {\IfFileExists{../../#1}%
         {\input{../../#1}}%
         {\typeout{ERROR: Unable to find #1}}%
      }%
   }%
}
\makeatletter
\def\FragNormKey#1{%
  \frag@norm@check#1def_\relax{#1}%
}
\def\frag@norm@check#1def_#2\relax#3{%
  \if\relax\detokenize{#1}\relax
    #3%
  \else
    def_#3%
  \fi
}
\makeatother

\newcommand{\FragmentName}[1]{\RootFile{fragment/\FragNormKey{#1}.tex}}

\newenvironment{fragment}[1]{}{}

\makeatletter
\AtBeginDocument{%
   \let\real@latex@label\label
   \def\label#1{%
      \global\@namedef{all@labels@seen@#1}{1}%
      \real@latex@label{#1}%
   }%
}

\DeclareDocumentCommand{\IncFragment}{m}{%
   \begingroup
      \def\label##1{%
         \@ifundefined{all@labels@seen@##1}{%
            \global\@namedef{all@labels@seen@##1}{1}%
            \real@latex@label{##1}%
         }{}%
      }%
      \IfFileExists{\FragmentName{#1}}%
         {\InputExt{\FragmentName{#1}}}%
         {\InputExt{\RootFile{fragment/#1.tex}}}%
   \endgroup
}
\makeatother

\newcommand{\Quote}[2]{%
   \SaveIndent
   \begin{center}%
      {\footnotesize
         \begin{quote}
             \RestoreIndent
             {#1}\\
             \hspace*{2cm}-- {#2.}
         \end{quote}%
      }%
      \index{quotation!#2}%
   \end{center}%
   \bigskip
}

\newcommand{\QuotePExt}[3]{%
   \bigskip%
   \hfill%
   \begin{minipage}{0.95\linewidth}%
       \footnotesize {#1}%

       \rule{\linewidth}{0.3mm}%
       \hfill #2, #3%
   \end{minipage}%
   \bigskip%
}

\epigraphfontsize{\small\itshape}
\setlength{\epigraphwidth}{0.6\textwidth}
\newcommand{\QuoteExt}[3]{%
   \vspace*{-1.6cm}%
   \epigraph{{\footnotesize #1}}{{#2}, {#3}}%
   \vspace*{-1.2cm}%
}

\providecommand{\delX}[1]{{\color{red}\sout{#1}}}
\providecommand{\newX}[1]{{\color{blue}#1}}
\providecommand{\chgY}[2]{{\color{red}\sout{#1}}{\color{blue}#2}}
\providecommand{\remX}[1]{{\color{purple}\textbf{[Remark: #1]}}}
\newcommand{\SpellIgnore}[1]{#1}
\newcommand{\EQ}{\mathrel{\;\;=\;\;}}
\newcommand{\LEQ}{\mathrel{\;\;\leq\;\;}}
\newcommand{\LT}{\mathrel{\;\;<\;\;}}
\newcommand{\GT}{\mathrel{\;\;<\;\;}}

\newcommand{\GEQ}{\mathrel{\;\;\geq\;\;}}

%-----------------------------------------------------------------------------
% Save a symbol that we know is going to get redefined.
\def\savesymbol#1{%
  \expandafter\let\expandafter\origsym\expandafter=\csname#1\endcsname
  \expandafter\let\csname orig#1\endcsname=\origsym
  \expandafter\let\csname#1\endcsname=\relax
}

% Restore a previously saved symbol, and rename the current one.
\def\restoresymbol#1#2{%
  \expandafter\let\expandafter\newsym\expandafter=\csname#2\endcsname
  \expandafter\global\expandafter\let\csname#1#2\endcsname=\newsym
  \expandafter\let\expandafter\origsym\expandafter=\csname orig#2\endcsname
  \expandafter\global\expandafter\let\csname#2\endcsname=\origsym
}

\usepackage{pictex}%
\savesymbol{triangle}%
\usepackage{tree}%
\restoresymbol{TreePKG}{triangle}%

% Definitions absorbed from sbook.sty
\providecommand{\MdPrc}{\nu}

\StartDefiningTabulars
\providecommand{\TwoFiguresBase}[4]{%
   \begin{tabular}{cc}
      \begin{minipage}{0.48\linewidth}
          \begin{center}
              {{{#1}}}
          \end{center}
      \end{minipage}
      &
      \begin{minipage}{0.48\linewidth}
          \begin{center}
              {{#3}}
          \end{center}
      \end{minipage}
      \\
      {#2} & {#4}
   \end{tabular}%
}

\providecommand{\TwoFigures}[6]{%
   \begin{tabular}{cc}
      \begin{minipage}{0.48\linewidth}
          \begin{center}
              {\IncludeGraphics[#1]{#2}}
          \end{center}
      \end{minipage}
      &
      \begin{minipage}{0.48\linewidth}
          \centerline{\IncludeGraphics[#4]{#5}}
      \end{minipage}
      \\
      {#3} & {#6}
   \end{tabular}%
}

\providecommand{\ThreeFigures}[6]{%
   \begin{tabular}{ccc}
        \begin{minipage}{0.30\linewidth}
            \IncludeGraphics[width=5.1cm]{#1}\\
        \end{minipage}
        &
        \begin{minipage}{0.30\linewidth}
            \IncludeGraphics[width=5.1cm]{#2}\\
        \end{minipage}
        &
        \begin{minipage}{0.30\linewidth}
            \IncludeGraphics[width=5.1cm]{#3}\\
        \end{minipage}
        \vspace{-1cm}
        \\
        {#4} & {#5}& {#6}
   \end{tabular}%
}
\StopDefiningTabulars

%-----------------------------------------------------------------------------
\providecommand{\remove}[1]{}
\providecommand{\NotAMSMode}[1]{#1}
\providecommand{\Solution}[1]{}
\providecommand{\Version}[1]{\hfill{Version: #1}}
\providecommand{\Distor}{\mathrm{dist}}
\providecommand{\Hardness}[1]{(h:#1)}
\providecommand{\QuoteOptExt}[3]{\QuoteExt{#1}{#2}{#3}}
\providecommand{\Hypercube}{\EuScript{H}}
\providecommand{\emphiC}[1]{\emphic{#1}{#1}}
\providecommand{\CoverTime}{\EuScript{C}}
\providecommand{\Nbr}{\Gamma}
\providecommand{\CommuteTime}[2]{\mathbf{CT}_{#1#2}}
\providecommand{\Resistance}{\mathbf{R}}
\providecommand{\EffectResist}[2]{\mathbf{R}_{#1#2}}
\providecommand{\NetworkE}{\EuScript{N}}
\providecommand{\trinom}[4]{\binom{#1}{#2 \;\; #3 \;\; #4}}
\providecommand{\Ranges}{\EuScript{R}}
\providecommand{\CClass}[1]{\ensuremath{\mathbf{#1}}}
\providecommand{\coCClass}[1]{\ensuremath{\mathrm{co-}}\CClass{#1}}
\providecommand{\CX}{\EuScript{C}}
\providecommand{\packet}{\mathbf{c}}
\providecommand{\packetA}{\mathbf{d}}
\providecommand{\mopt}{m_{\mathrm{opt}}}
\providecommand{\girth}{\mathop{\mathrm{girth}}}
\providecommand{\CR}{\mathop{\mathrm{cr}}}
\providecommand{\TRUE}{{\textcolor[named]{RedViolet}{\textsf{TRUE}}}\xspace}
\providecommand{\FALSE}{{\textcolor[named]{RedViolet}{\textsf{FALSE}}}\xspace}
\providecommand{\Part}[1]{\part{#1}}
\providecommand{\Section}[1]{\section{#1}}
\providecommand{\SubSection}[1]{\subsection{#1}}
\providecommand{\SubSubSection}[1]{\subsubsection{#1}}
\providecommand{\Exercise}[5]{%
\begin{exercise_h}[#1]%
   \exelab{#3}{}%
   #4%
\end{exercise_h}%
}
\providecommand{\si}[1]{#1}
\providecommand{\myParagraph}[1]{\noindent\textbf{#1}}
\providecommand{\DefName}[1]{#1}
\providecommand{\deflabel}[1]{#1}
\providecommand{\Description}[1]{}

%-----------------------------------------------------------------------------
% Mathematical Notations & Environments (from notations.sty)
%-----------------------------------------------------------------------------
%-----------------------------------------------------------------------------
% List environments (questions & compactenum)
%-----------------------------------------------------------------------------

\newlist{compactenumA}{enumerate}{5}%
\setlist[compactenumA]{topsep=0pt,itemsep=-1ex,partopsep=1ex,parsep=1ex,%
   label=(\Alph*)}%

\newlist{compactenumI}{enumerate}{5}%
\setlist[compactenumI]{topsep=0pt,itemsep=-1ex,partopsep=1ex,parsep=1ex,%
   label=(\Roman*)}%

\newlist{compactenuma}{enumerate}{5}%
\setlist[compactenuma]{topsep=0pt,itemsep=-1ex,partopsep=1ex,parsep=1ex,%
   label=(\alph*)}%

\newlist{compactitem}{enumerate}{5}%
\setlist[compactitem]{topsep=0pt,itemsep=-1ex,partopsep=1ex,parsep=1ex,%
   label=\textbullet}%

\newlist{compactenumi}{enumerate}{5}%
\setlist[compactenumi]{topsep=0pt,itemsep=-1ex,partopsep=1ex,parsep=1ex,%
   label=(\roman*)}%

\setcounter{secnumdepth}{5}

\newcommand{\HST}{\Term{HST}\xspace}
\newcommand{\Term}[1]{\textsf{#1}}
\newcommand{\TermI}[1]{\Term{#1}\index{#1@\Term{#1}}}

\DeclareMathAlphabet{\mathpzc}{OT1}{pzc}{m}{it}

\newcommand{\seg}{\mathsf{s}}

\newcommand{\Diam}{\mathrm{diam}}
\newcommand{\diam}{\mathrm{diam}}

\newcommand{\Copt}{C_{\mathrm{opt}}}
\newcommand{\CenterPnt}{\mathrm{cen}}

\newcommand{\parent}{\overline{\mathrm{p}}}
\newcommand{\parentX}[2][\!]{\parent\pth[#1]{#2}}

\newcommand{\FMS}{\EuScript{X}}  
\newcommand{\FMSB}{\EuScript{Y}}  
\newcommand{\FastCut}{\Algorithm{FastCut}\xspace}
\newcommand{\MinCut}{\Algorithm{MinCut}\xspace}
\newcommand{\MinCutRep}{\Algorithm{MinCutRep}\xspace}
\newcommand{\Contract}{\Algorithm{Contract}\xspace}

\newcommand{\Grid}{\mathrm{G}}
\newcommand{\ClosestPair}{\mathcal{CP}}
\newcommand{\SphereChar}{\ensuremath{\mathbb{S}}}
\newcommand{\Sphere}[1]{{\SphereChar^{(#1)}}}

\newcommand{\PntSet}{{\mathsf{P}}}%
\providecommand{\PntSetA}{{\mathsf{Q}}}%

\newcommand{\query}{\mathcalb{q}}

\newcommand{\pnt}{\mathsf{p}}
\newcommand{\pntA}{\mathsf{u}}
\newcommand{\pntB}{\mathsf{v}}
\newcommand{\pntC}{{\mathsf{x}}}
\newcommand{\pntD}{{\mathsf{y}}}
\newcommand{\pntE}{{\mathsf{z}}}

\newcommand{\pntCoord}{p}
\newcommand{\pntCoordA}{q}

%%%%%%%%%%%%%%%%%%%%%%%%%%%%%%%%%%%%%%%%%%%%%%%%%%%%%%%%%%%%%%%%%%%%% 
%%%%%%%%%%%%%%%%%%%%%%%%%%%%%%%%%%%%%%%%%%%%%%%%%%%%%%%%%%%%%%%%%%%%%
% Markov chain notations.

\newcommand{\MCSpc}{\mathtt{S}}
\newcommand{\MCPrb}{\mathtt{P}}
\newcommand{\MCPrbY}[2]{\MCPrb_{#1#2}}
\newcommand{\MCPrbM}{\mathbf{P}}%
\newcommand{\MCPrbT}[3]{\MCPrb_{#1#2}^{(#3)}}
\newcommand{\MCFPrbT}[3]{\mathtt{r}_{#1#2}^{(#3)}}
\newcommand{\MCVstProb}[2]{\mathtt{f}_{#1#2}}
\newcommand{\MCExSteps}[2]{\mathtt{h}_{#1#2}}
\newcommand{\MCPVec}[1]{\mathbf{q}^{(#1)}}
%%%%%%%%%%%%%%%%%%%%%%%%%%%%%%%%%%%%%%%%%%%%%%%%%%%%%%%%%%%%%%%%%%%%%

\newcommand{\Spread}{\Phi}
\newcommand{\VC}{\ensuremath{\mathsf{VC}}\xspace}
\newcommand{\MC}{\ensuremath{\mathsf{MC}}\xspace}
\newcommand{\hyperplane}{h}

\newcommand{\QuickSort}{\Algorithm{QuickSort}\xspace}
\newcommand{\Alg}{\Algorithm{Alg}\xspace}

\newcommand{\Graph}{\mathsf{G}}%
\newcommand{\GraphB}{\mathsf{H}}
\newcommand{\GraphC}{\mathsf{K}}
\newcommand{\GraphD}{\mathsf{F}}

\providecommand{\G}{}
\renewcommand{\G}{\Graph}

\newcommand{\RLP}{\ComplexityClass{RLP}\xspace}
\newcommand{\USTCON}{\ComplexityClass{USTCON}}
\newcommand{\STCON}{\overrightarrow{\ComplexityClass{STCON}}}

\newcommand{\Mtrx}{\mathsf{M}}
\newcommand{\TMtrx}{\mathsf{Q}}
\newcommand{\IMtrx}{\mathcal{I}}

\newcommand{\Expander}[3]{$\left[#1,#2,#3\right]$-expander}

\newcommand{\ChartX}[1]{\ensuremath{\chi_{#1}}}
\newcommand{\Field}{\mathbb{F}}
\newcommand{\LSpace}{\mathbb{G}}
\newcommand{\LD}[2]{\mathsf{L{D}}\pth{#1,#2}}
\newcommand{\SEVN}{\widehat{\lambda_2}}
\newcommand{\NX}[2]{#1\!\left[#2\right]}
\newcommand{\Entry}[2]{#1\!\left[#2\right]}
\newcommand{\NXExt}[3]{#1_{#2}\!\left[#3\right]}
\newcommand{\RangeI}[1]{\llbracket #1 \rrbracket}
\newcommand{\DialX}[1]{\gamma\pth{#1}}
\newcommand{\BiDialX}[1]{{\gamma}_2\pth{#1}}

\newcommand{\EV}[1]{\widehat{\lambda_{#1}}}
\newcommand{\TEV}[1]{\widehat{\lambda_{#1}}}
\newcommand{\EVal}{\widehat{\lambda}}
\newcommand{\Basis}{\mathcal{B}}

\newcommand{\ZProd}[2]{\ensuremath{#1 \, \text{\textcircled{z}} #2}}
\newcommand{\RProd}[2]{\ensuremath{#1 \, \text{\textcircled{r}} #2}}

\newcommand{\OneVec}[1]{\mathbf{1}^{#1}}
\newcommand{\DstrbV}[1]{\mathbf{\mathsf{q}}^{(#1)}}
\newcommand{\DstrbE}[2]{\mathsf{q}^{(#1)}_{#2}}

\newcommand{\DEdges}{\widetilde{\Edges}}

\newcommand{\VV}{\mathsf{V}}%
\newcommand{\EE}{\mathsf{E}}%

\newcommand{\Vertices}{\VV}%
\newcommand{\Edges}{\EE}%

\newcommand{\EdgesX}[1]{\Edges\pth{#1}}
\newcommand{\VerticesX}[1]{V\pth{#1}}

\newcommand{\StringX}[1]{\mathsf{s}\pth{#1}}

\newcommand{\BMtrx}{\mathsf{B}}
\newcommand{\WMtrx}{\mathsf{W}}
\newcommand{\NWMtrx}{\overline{\mathsf{W}}}
\newcommand{\Sequence}{\EuScript{S}}

\newcommand{\med}{\mathrm{med}}
\newcommand{\medX}[1]{\mathrm{med}\pth{#1}}

\newcommand{\medianC}{\nu}%

\newlength{\savedparindent}
\newcommand{\SaveIndent}{\setlength{\savedparindent}{\parindent}}
\newcommand{\RestoreIndent}{\setlength{\parindent}{\savedparindent}}

\newcommand{\NNN}{\mathbb{N}}%
\providecommand{\NN}{\mathbb{N}}%
\renewcommand{\Re}{\mathbb{R}}%
\newcommand{\reals}{\Re}%
\newcommand{\ZZ}{\mathbb{Z}}%

\newcommand{\level}{\mathop{\mathrm{level}}}

\providecommand{\TPDF}[2]{\texorpdfstring{#1}{#2}}

\newcommand{\Hint}[1]{\emph{[#1]}}
\newcommand{\Arr}{\mathop{\mathrm{\EuScript{A}}}}

\newcommand{\Array}{\mathcal{T}}%

\newcommand{\dirEdge}[2]{\left({{#1} \rightarrow {#2}} \right)}

%%%%%%%%%%%%%%%%%%%%%%%%%%%%%%%%%%%%%%%%%%%%%%%%%%%%%%%%%%%%%%%%%%%%%%
% Emphasize stuff...
\definecolor{blue25}{rgb}{0, 0, 11}
\newcommand{\emphic}[2]{%
   \textcolor{blue25}{%
      \bfseries{\emph{#1}}}%
   \index{#2}}
\newcommand{\emphi}[1]{\emphic{#1}{#1}}
\newcommand{\emphOnly}[1]{\textcolor{blue25}{\bfseries{\emph{#1}}}}

\definecolor{almostblack}{rgb}{0, 0, 0.3}
\providecommand{\emphw}[1]{{\textcolor{almostblack}{\emph{#1}}}}%

%
%%%%%%%%%%%%%%%%%%%%%%%%%%%%%%%%%%%%%%%%%%%%%%%%%%%%%%%%%%%%%%%%%%%%%%

\newcommand{\volX}[1]{\mathrm{vol}\pth{#1}}

\newcommand{\Interval}{\EuScript{I}}
\newcommand{\IntervalA}{\EuScript{I}}

\newcommand{\Exp}{\mathrm{Exp}}

\newcommand{\binsY}[2]{\sqcup_{\geq #1}\pth{#2}}
\newcommand{\ballsY}[2]{\mathcal{F}_{\geq #1}\pth{#2}}

\newcommand{\BLX}[1]{{{\sqcupplus{}}_{\geq #1}}}


\let\Ball\undefined
\newcommand{\Ball}{\mathbf{b}}
\newcommand{\hdim}{d}
\newcommand{\Domain}{\mathcal{D}}
\newcommand{\points}[1]{[\bfseries{#1 points}]}
\newcommand{\Event}{\EuScript{E}}%
\newcommand{\EventB}{\EuScript{B}}%
\newcommand{\EventF}{\EuScript{F}}%
\newcommand{\EventG}{\EuScript{G}}%

\newcommand{\distCmd}[1]{\left\| {#1}  \right\|}
\newcommand{\distY}[2]{\distCmd{#1 - #2}}

\newcommand{\NormalC}{\mathbf{N}}
\newcommand{\NormalY}[2]{\NormalC\pth{#1, #2}}%
\newcommand{\NormalDistOneDim}{\NormalC}
\newcommand{\NormalDistX}[1]{\NormalC^{#1}}
\DeclareMathOperator{\lca}{lca}
\newcommand{\RangeSpace}{\mathsf{S}}%
\newcommand{\RangeSpaceA}{\mathsf{T}}

\newcommand{\range}{\mathbf{r}}
\newcommand{\rangeA}{\mathbf{r}'}

\newcommand{\BadProb}{\varphi}

\newcommand{\EntropyChar}{\mathbb{H}}
\newcommand{\Entropy}[1]{\EntropyChar\pth{#1}}
\newcommand{\EntropyX}[1]{\EntropyChar\pth[]{#1}}

\newcommand{\Matrix}{\mathsf{M}}
\newcommand{\MatrixB}{\mathsf{B}}
\newcommand{\MatrixC}{\mathsf{C}}
\newcommand{\MatrixD}{\mathsf{D}}

\renewcommand{\th}{th\xspace}

\newcommand{\Family}{\mathcal{F}}
\newcommand{\FamilyB}{\mathcal{G}}
\newcommand{\FamilyC}{\mathcal{H}}
\newcommand{\FMSC}{\EuScript{Y}\index{metric space}}

\newcommand{\median}{\overline{\mathsf{m}}}
\newcommand{\QuickSelect}{\Algorithm{Quick{}Select}\xspace}

\providecommand{\row}{\mathbf{r}}%
\newcommand{\rowc}{{r}}%

\newcommand{\vecb}{\mathbf{b}}

\newcommand{\vecA}{\mathbf{v}}
\newcommand{\vecB}{\mathbf{u}}
\newcommand{\vecRand}{\mathbf{r}}

\newcommand{\vecH}{\mathbf{h}}
\newcommand{\vecV}{\mathbf{v}}

\newcommand{\vecX}{\mathbf{x}}
\newcommand{\vecY}{\mathbf{y}}

\newcommand{\DistBinom}[2]{\mathrm{Bin}\pth{#1,#2}}

\newcommand{\DistGeom}[1]{\mathrm{Geom}\pth{#1}}

\newcommand{\Sample}{\mathsf{R}}

\newcommand{\lcm}{\mathop{\mathrm{lcm}}}
\newcommand{\bdiv}{\mathop{\mathrm{div}}}%
\newcommand{\divides}{\mid}%
\newcommand{\poly}{\mathrm{poly}}%

\newcommand{\Group}{\mathcal{G}} 
\newcommand{\GroupA}{\mathcal{H}} 

\newcommand{\algGCD}{\Algorithm{EuclidGCD}\xspace}%

\newcommand{\gOp}{\times}%
\newcommand{\gIdent}{\mathsf{i}}%

\newcommand{\ordX}[2][\!]{\mathop{\mathrm{o{r}d}}\pth[#1]{#2}}

\newcommand{\lSym}[2]{%
   \left(
       \ds {#1}\divides{#2}%
   \right) %
}
\newcommand{\jSym}[2]{\left\llbracket%
       \ds {#1}\divides{#2}%
       \right\rrbracket%
}
\newcommand{\Line}{\ell}

\newcommand{\LevelChar}{\mathbb{L}}
\newcommand{\AtMostKLevelChar}{\LevelChar_{\leq k}}
\newcommand{\AtMostKLevel}[1]{\LevelChar_{\leq k} \pth{ #1}}

\newcommand{\LineSet}{\mathsf{L}}%
\newcommand{\ArrX}[1]{\EuScript{A}\pth{#1}}%

\DeclareRobustCommand{\Matousek}{Matoušek\xspace}
\DeclareRobustCommand{\Erdos}{Erdős\xspace}
\DeclareRobustCommand{\Lovasz}{Lovász\xspace}

\newcommand{\ProblemC}[1]{{\textsf{\textcolor{red}{#1}\index{problem!#1}}}}
\DefineNamedColor{named}{OliveGreen} {cmyk}{0.64,0,0.95,0.40}
\providecommand{\ComplexityClass}[1]{{{\textcolor[named]{OliveGreen}{%
            \textsc{#1}}}}}
\providecommand{\NPHard}{{\ComplexityClass{NP-Hard}}%
   \index{NP!hard}\xspace}
\providecommand{\NPComplete}{\ComplexityClass{NP-Complete}%
   \index{NP!complete}\xspace}
\providecommand{\APXHard}{\ComplexityClass{APX-Hard}%
   \index{APX!Hard}\xspace}
\providecommand{\NPC}{\ComplexityClass{NPC}%
   \index{NP!complete}\xspace}
\providecommand{\NP}{\ComplexityClass{NP}%
   \index{NP}\xspace}
\providecommand{\POLYT}{\ComplexityClass{P}\xspace}

\newcommand{\envFrame}[1]{\begin{center}\fbox{\begin{minipage}[c]{0.75\textwidth}{#1}\end{minipage}}\end{center}}

\newcommand{\mysqrt}[1]{\ts\!\!\sqrt{#1}}
\newcommand{\SegSet}{\mathsf{S}}
\newcommand{\ArrVD}[1]{\Arr^{|} \pth{#1}}
\newcommand{\RIC}{\textsf{RIC}\index{random incremental
      construction}\xspace}
\newcommand{\ArrB}{\EuScript{B}}

\newcommand{\TreeOpt}{\EuScript{T}_\mathrm{opt}}
\newcommand{\TreeOptA}{\EuScript{T}_\mathrm{opt}'}

\newcommand{\DepthExt}[2]{\mathrm{depth}_{#1}\pth{#2}}
\newcommand{\depthX}[1]{\mathrm{depth}\pth{#1}}

\newcommand{\TreeA}{\EuScript{U}}
\newcommand{\TreeB}{\EuScript{V}}

\newcommand{\Trap}{\sigma}
\newcommand{\TrapA}{\sigma'}
\newcommand{\TrapB}{\sigma''}

\newcommand{\SampleA}{\mathsf{R}'}

\newcommand{\AllRegions}{\mathcal{T}}
\newcommand{\RegionsSet}[1]{\mathcal{F}\pth{#1}}
\newcommand{\RegionsHeavy}[2]{\mathcal{F}_{\!\geq#1}\pth{#2}}
\newcommand{\NumOfRegionsZero}[1]{\ExChar \!f\pth{#1}}
\newcommand{\NumOfRegions}[2]{\ExChar\! f_{\!\geq #1}\pth[]{#2}}
\newcommand{\ObjSet}{\mathsf{S}}

\newcommand{\WeightX}[2][\!]{\omega \pth[#1]{#2}}
\newcommand{\Edge}{\mathsf{e}}

\newcommand{\Face}{\mathsf{f}}

\newcommand{\DefSet}[1]{D\pth{#1}}
\newcommand{\KillSet}[1]{K\pth{#1}}

\newcommand{\CSProbFull}[4]{
   \,\rho_{\!\!\stackrel{\rule[-.005cm]{0cm}{0.1cm}}{#1,
       #4}}\pth{\MakeSBig #2,#3}}

\newcommand{\disk}{D}
\newcommand{\diskA}{D'}

\newcommand{\IntervalSet}{\widehat{\mathcal{I}}}

\newcommand{\GroundSet}{\textsf{X}}
\newcommand{\FGroundSet}{\mathsf{x}}
\newcommand{\RangeSet}{{\mathcal{R}}}
\newcommand{\RangeSetA}{\mathcal{R}'}
\newcommand{\MeasureChar}{\overline{m}}
\newcommand{\Measure}[1]{\MeasureChar\pth{#1}}
\newcommand{\sMeasure}[1]{\overline{s}\pth{#1}}

\newcommand{\SetA}{S}%
\newcommand{\SetB}{T}%
\newcommand{\SetC}{N}%

\newcommand{\VCProj}[2]{#1_{|#2}}
\newcommand{\VCDim}{\mathrm{dim}_{\VC}}
\newcommand{\Dim}{\delta}
\newcommand{\DimA}{\delta'}
\newcommand{\PntSetB}{{{\mathsf{Q}}}}
\newcommand{\PntSetC}{{{\mathsf{R}}}}
\newcommand{\Growth}[2]{{\EuScript{G}}_{#1}\pth{#2}}
\newcommand{\ShatterFunc}[1]{\pi_{#1}}
\newcommand{\Distribution}{\EuScript{D}}
\newcommand{\Error}{\EuScript{E}}

\newcommand{\bd}{{\partial}}%
\newcommand{\eps}{{\varepsilon}}%

\newcommand{\clmlab}[1]{\label{claim:#1}}
\newcommand{\clmref}[1]{\HLink{Claim}{claim:#1}}

\newcommand{\CHX}[1]{{\mathcal{CH}}\pth{#1}}

\newcommand{\ts}{\hspace{0.6pt}}

\def\qedsymbol{\myqedsymbol}%

\newcommand{\qedsign}{\hfill{\hfill\myqedsymbol}}
\newcommand{\qedsignBig}{\hfill{\hfill\rule{4mm}{4mm}}}
\newcommand{\QED}{{\hfill{\hfill\qedsymbol}}}

\newcommand{\num}{\alpha}

\newcommand{\NEW}[1]{#1}

\StartDefiningTabulars
\newcommand{\ThreeFiguresExt}[9]{%
   \begin{tabular}{ccc}
     \begin{minipage}{0.31\linewidth}
         \begin{center}
             {\IncludeGraphics[#1]{#2}}
         \end{center}
     \end{minipage}
     &
       \begin{minipage}{0.31\linewidth}
           \centerline{\IncludeGraphics[#4]{#5}}
       \end{minipage}
     &
       \begin{minipage}{0.31\linewidth}
           \centerline{\IncludeGraphics[#7]{#8}}
       \end{minipage}
     \\
     ~\\[-0.2cm]
     {#3} & {#6} & {#9}
   \end{tabular}}
\StopDefiningTabulars

\newcommand{\IncludeGraphics}[2][]{%
   \typeout{}%
   \typeout{Graphics: #2}%
   \typeout{\ includegraphics[#1]{#2}}%
   \includegraphics[#1]{#2}
   \typeout{}%
}

%    -----------------------------------------------------------------
%    -----------------------------------------------------------------
%    -----------------------------------------------------------------
%    Handling references
%    -----------------------------------------------------------------
%    -----------------------------------------------------------------
%    -----------------------------------------------------------------
\newcommand{\HLinkShort}[2]{\hyperref[#2]{#1\ref*{#2}}}
\newcommand{\HLink}[2]{\hyperref[#2]{#1~\ref*{#2}}}
\newcommand{\HLinkPage}[2]{\hyperref[#2]{#1~\ref*{#2}%
      $_\text{p\pageref{#2}}$}}

\newcommand{\HLinkSuffix}[3]{\hyperref[#2]{#1\ref*{#2}{#3}}}
\newcommand{\HLinkPageSuffix}[3]{\hyperref[#2]{#1\ref*{#2}%
      #3$_\text{p\pageref{#2}}$}}

\newcommand{\figlab}[1]{\label{fig:#1}}
\newcommand{\figref}[1]{\HLink{Figure}{fig:#1}}

\newcommand{\exelab}[1]{\label{exer:#1}}
\newcommand{\exeref}[1]{Exercise~\ref{exer:#1}}

\providecommand{\deflab}[1]{\label{def:#1}}
\newcommand{\defref}[1]{\HLink{Definition}{def:#1}}
\newcommand{\defrefY}[2]{\hyperref[def:#2]{#1}}
\newcommand{\defrefpage}[1]{\HLinkPage{Definition}{def:#1}}

\newcommand{\etal}{\textit{et~al.}\xspace}

\newcommand{\obslab}[1]{\label{observation:#1}}
\newcommand{\obsref}[1]{\HLink{Observation}{observation:#1}}
\newcommand{\obsrefY}[2]{\hyperref[observation:#2]{#1}}

\newcommand{\exmlab}[1]{\label{example:#1}}
\newcommand{\exmref}[1]{\HLink{Example}{example:#1}}

\newcommand{\remlab}[1]{\label{rem:#1}}
\newcommand{\remref}[1]{\HLink{Remark}{rem:#1}}%

\newcommand{\chaplab}[1]{\label{chap:#1}}

\newcommand{\itemlab}[1]{\label{item:#1}}
\newcommand{\itemref}[1]{\HLinkSuffix{}{item:#1}{}}

\newcommand{\tedlab}[1]{\label{tedium:#1}}
\newcommand{\tedref}[1]{\HLink{Tedium}{tedium:#1}}%

\newcommand{\factlab}[1]{\label{fact:#1}}%

\newcommand{\alglab}[1]{\label{algorithm:#1}}

\newcommand{\net}{\mathcal{N}}%
      
\newcommand{\lemlab}[1]{\label{lemma:#1}}
\newcommand{\lemref}[1]{\HLink{Lemma}{lemma:#1}}%
\newcommand{\lemrefY}[2]{\hyperref[lemma:#2]{#1}}
\newcommand{\lemrefpage}[1]{\HLinkPage{Lemma}{lemma:#1}}
\newcommand{\lemrefshort}[1]{\HLinkShort{L}{lemma:#1}}

\newcommand{\eqlab}[1]{\label{equation:#1}}
\newcommand{\Eqrefbare}[1]{\HLinkSuffix{(}{equation:#1}{)}}
\newcommand{\Eqref}[1]{\HLinkSuffix{Eq.~(}{equation:#1}{)}}
\newcommand{\Eqrefpage}[1]{\HLinkPageSuffix{Eq.~(}{equation:#1}{)}}

\newcommand{\seclab}[1]{\label{sec:#1}}
\newcommand{\secref}[1]{\HLink{Section}{sec:#1}}
\newcommand{\secrefY}[2]{\hyperref[sec:#2]{#1}}
\newcommand{\secrefpage}[1]{\HLinkPage{Section}{sec:#1}}

\newcommand{\thmlab}[1]{{\label{theo:#1}}}
\newcommand{\thmref}[1]{\HLink{Theorem}{theo:#1}}
\newcommand{\thmrefY}[2]{\hyperref[theo:#2]{#1}}
\newcommand{\thmrefpage}[1]{\HLinkPage{Theorem}{theo:#1}}

\newcommand{\corlab}[1]{\label{cor:#1}}
\newcommand{\corref}[1]{\HLink{Corollary}{cor:#1}}%

%    -----------------------------------------------------------------
%    -----------------------------------------------------------------
%    -----------------------------------------------------------------
%    Handling references
%    -----------------------------------------------------------------
%    -----------------------------------------------------------------
%    -----------------------------------------------------------------

\newcommand{\ceil}[1]{\left\lceil {#1} \right\rceil}
\newcommand{\floor}[1]{\left\lfloor {#1} \right\rfloor}
\newcommand{\brc}[1]{\left\{ {#1} \right\}}
\newcommand{\pth}[2][\!]{\mleft({#2}\mright)}%
\newcommand{\brk}[1]{\mleft[{#1}\mright]}%
\newcommand{\pbrc}[2][\!\!]{#1\left[ {#2} \bigr. \right]}
\newcommand{\pbrcS}[1]{\left[ {#1} \right]}

\newcommand{\dotprod}[2]{{#1}\cdot{#2}}
\newcommand{\DotProd}[2]{\permut{{#1},{#2}}}
\newcommand{\permut}[1]{\left\langle {#1} \right\rangle}

\newcommand{\MakeBig}{}
\newcommand{\MakesBig}{}
\newcommand{\MakeSBig}{}
\newcommand{\MakeVBig}{}

\newcommand{\DistChar}{\mathsf{d}}%
\newcommand{\DistSetY}[2]{\DistChar\pth{#1,#2}}
\newcommand{\distG}[2]{\DistChar_\Graph\pth{#1, #2}}
\newcommand{\DistMtrCharX}[1]{\DistChar_{#1}}

\newcommand{\dY}[2]{\DistChar\pth{#1, #2}}%
\newcommand{\dGY}[2]{\DistChar_\Graph\pth{#1, #2}}%

\newcommand{\cset}{\mathcal{C}}%

\newcommand{\MTR}{\mathcal{M}}
\newcommand{\DistMtr}[2]{\DistChar\pth{#1,#2}}

\newcommand{\dist}[1]{\left\| {#1} \right\|}
\newcommand{\norm}[1]{\left\| {#1} \right\|}

\providecommand{\ds}{\displaystyle}
\newcommand{\sep}[1]{\,\left|\, {#1} \bigr.\right.}%
\newcommand{\sepw}[1]{\left|\, {#1} \right.}%

\newcommand{\ProbExt}[2]{\mathop{\mathbf{Pr}}_{#1}\pbrcx{#2}}
\newcommand{\ProbLTR}{\mathbb{P}}%

\newcommand{\Prob}[1]{\mathop{\ProbLTR} \mleft[ #1 \mright]}%
\newcommand{\ProbCond}[2]{\mathop{\ProbLTR}\!\left[%
       #1 \;\middle\vert\; #2 \right]}
\newcommand{\ProbY}[2]{\ProbLTR^{}_{\!#1}\pbrc{#2}}

\newcommand{\ExChar}{\mathbb{E}}%
\newcommand{\ExSym}{\mathop{\ExChar}}%
\newcommand{\Ex}[2][\!]{\ExSym#1\pbrcx{#2}}
\newcommand{\ExExt}[2]{\ExSym_{#1}\!\pbrcx{#2}}
\newcommand{\ExExtS}[2]{\ExSym_{\substack{#1}}\!\pbrcx{#2}}
\newcommand{\ExCond}[2]{\ExSym\!\left[%
       #1 \;\middle\vert\; #2 \right]}

\newcommand{\VarSymb}{\mathop{\mathbb{V}}}
\newcommand{\Var}[1]{\mathop{\mathbb{V}}\!\pbrcx{#1}}

\newcommand{\pbrcx}[1]{\left[ {#1} \right]}%
\newcommand{\nfrac}[2]{{#1}/{#2}}
\newcommand{\cardin}[1]{\left| {#1} \right|}%

\newcommand\CCancel[2][black]{{\renewcommand\CancelColor{\color{#1}}\cancel{#2}}}

\newcommand{\IncGraphPage}[4][]{%
   \IncludeGraphics[page=#4,#1]{\File{#2/#3}}%
}

\newenvironment{enumEx}{
   \begin{compactenumA}
    }{\end{compactenumA}}

\newcommand{\cl}{{\mathrm{cl}}}

\usepackage[amsmath,thmmarks]{ntheorem}
\theoremseparator{.}%

\newtheorem{theorem:unnumbered}{Theorem}[section]
\newtheorem{theorem}{Theorem}[section]

\newtheorem{observation}[theorem]{Observation}
\newtheorem{conjecture}[theorem]{Conjecture}
\newtheorem{lemma}[theorem]{Lemma}
\newtheorem{corollary}[theorem]{Corollary}
\newtheorem{claim}[theorem]{Claim}
\newtheorem{assumption}[theorem]{Assumption}
\newtheorem{tedium}[theorem]{Tedium}

\theoremstyle{plain}%
\theoremheaderfont{\sffamily}
\theorembodyfont{\upshape}%

\newtheorem{defn}[theorem]{Definition}
\newtheorem{definition}[theorem]{Definition}

\newtheorem*{remarkunnumbered}{Remark}%
\newenvironment{remark_unnumbered}{\begin{remarkunnumbered}}{\end{remarkunnumbered}}
\newtheorem*{remarks}{Remarks}%
\newtheorem{remark}[theorem]{Remark}%
\newtheorem{fact}[theorem]{Fact}%
\newtheorem{example}[theorem]{Example}
\newtheorem{exercise}[theorem]{Exercise}
\newtheorem{problem}[theorem]{Problem}
\newtheorem{algorithm}[theorem]{Algorithm}%
\newtheorem{xca}[theorem]{Exercise}
\newtheorem{openproblem}[theorem]{Open Problem}
\newtheorem{exerciseh}[theorem]{Exercise}
\newenvironment{exercise_h}[1][]{\begin{exerciseh}[#1]}{\end{exerciseh}}

% Proof environments
\theoremheaderfont{\em}%
\theorembodyfont{\upshape}%
\theoremstyle{nonumberplain}
\theoremseparator{}
\theoremsymbol{\myqedsymbol}
\newtheorem{proof}{Proof:}
\newtheorem{proofextin}{Proof}
\newtheorem{altproof}{Alternative proof:}

\theoremsymbol{}
\newtheorem{proofnoqed}{Proof:}

\theoremsymbol{\myqedsymbol}
\theorempreskip{0pt}%
\theorempostskip{0pt}%
\newtheorem{proofof}{Proof of\!}%

\newenvironment{proofExt}[1]{{\emph{#1:}}}{\hfill{\hfill\myqedsymbol}}
\newenvironment{proofnoproof}{{}}{{\hfill{\hfill\myqedsymbol}}}

%%%%%%%%%%%%%%%%%%%%%%%%%%%%%%%%%%%%%%%%%%%%%%%%%%%%%%%%%%%%%%%%%%%%
%%%%%%%%%%%%%%%%%%%%%%%%%%%%%%%%%%%%%%%%%%%%%%%%%%%%%%%%%%%%%%%%%%%%
% Macros to embed figures in text

\newcommand{\EBRY}[2]{#1_{\langle#2\rangle}}%

\newcommand{\NiceFrac}[2]{%
   {\raisebox{0.2em}{\ensuremath{#1}}\!\!\!}%
   {\rotatebox[origin=c]{-20}{\Big/}}% Rotated /
   {\! \raisebox{-0.2em}{\ensuremath{#2}} }}%

% end - Macros to embed figures in text
%%%%%%%%%%%%%%%%%%%%%%%%%%%%%%%%%%%%%%%%%%%%%%%%%%%%%%%%%%%%%%%%%%%% 
%%%%%%%%%%%%%%%%%%%%%%%%%%%%%%%%%%%%%%%%%%%%%%%%%%%%%%%%%%%%%%%%%%%%

%%%%%%%%%%%%%%%%%%%%%%%%%%%%%%%%%%%%%%%%%%%%%%%%%%%%%%%%%%%%%%%%%%
% Color in math
\providecommand{\Mh}[1]{{#1}}%
\providecommand{\IfColorMath}[2]{#2}%

\IfFileExists{.latex_printer_friendly}{\def\GenPrinterVer{1}}{}%

\ifx\GenPrinterVer\undefined
   \IfFileExists{.latex_color}{\def\GenColorMath{1}}{}
\else
   \renewcommand{\IfPrinterVer}[2]{#1}%
\fi

\ifx\GenColorMath\undefined
\else
\renewcommand{\IfColorMath}[2]{#1}%
\renewcommand{\Mh}[1]{{\textcolor{red}{#1}}}%
\fi
% Color in math - end
%%%%%%%%%%%%%%%%%%%%%%%%%%%%%%%%%%%%%%%%%%%%%%%%%%%%%%%%%%%%%%%%%%

\newcommand{\HSetY}[2]{H\pth{#1, #2}}%
\newcommand{\CHSetY}[2]{C\pth{#1, #2}}%
\DeclareRobustCommand{\Holder}{Hölder\xspace}
\newcommand{\OldOm}{{\mathcal{O}}}
\newcommand{\NewOm}{{\mathcal{N}}}

%%%%%%%%%%%%%%%%%%%%%%%%%%%%%%%%%%%%%%%%%%%%%%%%%%%%%%%%%%%%%%%%%%%%%%%
%---------------------------------------------------------------
\newcommand{\IntRange}[1]{\left\llbracket #1 \right\rrbracket}
\newcommand{\IRX}[1]{\left\llbracket #1 \right\rrbracket}%
\newcommand{\IRY}[2]{\left\llbracket #1:#2 \right\rrbracket}

%---------------------------------------------------------------
%%%%%%%%%%%%%%%%%%%%%%%%%%%%%%%%%%%%%%%%%%%%%%%%%%%%%%%%%%%%%%%%%%%%%%%

\newcommand{\obj}{f}%
\newcommand{\Wg}{\omega}%
\newcommand{\TotalW}{W}%

\newcommand{\rankY}[2]{{#1}_{ \permut{#2}}}%%
\newcommand{\areaX}[1]{\mathrm{area}\pth{#1}}%%

\newcommand{\mybinom}[2]%
{\Bigl(\begin{array}{@{}c@{}}#1\\#2\end{array}\Bigr)}

\providecommand{\Set}[2]{\left\{ #1 \;\middle\vert\; #2 \right\}}
\newcommand*{\biga}[1]{\vcenter{\hbox{\scalebox{1.5}{\ensuremath#1}}}}
\newcommand{\CNF}{\TermI{CNF}\xspace}%
\newcommand{\Tree}{T}%

\DeclareSymbolFont{sfletters}{OML}{cmbrm}{m}{it}

\DeclareMathSymbol{\salpha}{\mathord}{sfletters}{"0B}

%%%%%%%%%%%%%%%%%%%%%%%%%%%%%%%%%%%%%%%%%%%%%%%%%%%%%%%%%%%%%%%%%%%%

\providecommand{\ProblemInd}[2]{\textsf{\textcolor{red}{#1}\index{Problem!#2}}}
\newcommand{\NPProblemInd}[4]{
   \begin{center}
      \vspace{0.6cm}
       \fbox{
          \begin{minipage}{0.9\linewidth}
              \vspace{-0.6cm} \hspace{-0.35cm}
              {\large \bfseries \ProblemInd{#1}{#2}}
              \vspace{+0.2cm}%

              \noindent{\textcolor{blue}{\bfseries \textsf{Instance}}:} {#3}
              
              \noindent{\textcolor{blue}{\bfseries \textsf{Question}:}} {#4}
          \end{minipage}
       }
   \end{center}
   \smallskip
   }

\newcommand{\NPProblemOpt}[4]{\NPProblemInd{#1}{#2}{#3}{#4}}

%%%%%%%%%%%%%%%%%%%%%%%%%%%%%%%%%%%%%%%%%%%%%%%%%%%%%%%%%%%%%%%%%%%%

\newcommand{\Hastad}{H{\aa}stad\xspace}

\newcommand{\ANNbr}{\Term{ANN}\index{approximate nearest neighbor}\xspace}
\newcommand{\LSH}{\TermI{L{S}H}\xspace}

\newcommand{\DS}{\EuScript{D}} % DataStructure
\newcommand{\DSC}{\mathsf{D}} % DataStructure

\newcommand{\distH}[2]{\DistChar_{H}\pth{#1, #2}}
\newcommand{\distHY}[2]{\DistChar_{H}\pth{#1, #2}}

\newcommand{\DSAprxNearNbr}{\ensuremath{\DataStructure{{\approx%
         \text{Near{}}}}}%
   \index{near neighbor!data-structure!approximate}%
}
\newcommand{\DataStructure}[1]{\ensuremath{\DS_{#1}}}
\newcommand{\DistNN}[2]{\DistChar\pth{#1,#2}}
\newcommand{\pntc}{{p}}%

\newcommand{\seq}{{M}}%
\newcommand{\cTimes}{\beta}%

\newcommand{\DSTimes}{L}%
\newcommand{\rad}{r}%

\newcommand{\HypercubeX}[1]{\Hypercube^{#1}}
\newcommand{\kFunctions}[2]{\mathrm{combine}\pth{#1,#2}}
\newcommand{\WEqualChar}{\Bumpeq}
\newcommand{\WEqual}[2]{{#1}\WEqualChar{#2}}
\newcommand{\FamilyAmp}{\FamilyC_{\WEqualChar}}
\newcommand{\PSPACE}{\ComplexityClass{PSPACE}\xspace}%

\newcommand{\PS}{P}%

\newcommand{\PSetFar}{\ensuremath{\PntSet^{}_{\geq}\xspace}}%
\newcommand{\DA}{D}%
\newcommand{\uniqX}[1]{\mathrm{uniq}\pth{#1}}%
\newcommand{\diffX}[1]{\mathop{}\!\mathrm{d}#1}

\newcommand{\Frq}{F}%
\newcommand{\frq}{f}%
\newcommand{\STRM}{\mathcal{S}}%
\newcommand{\strm}{{s}}%

\DeclareRobustCommand{\Turan}{Turán\xspace}
\newcommand{\fGrowX}[1]{G^{}_{\Dim}\pth{#1}}
\newcommand{\CompressFunc}{\mathtt{Compress}}
\newcommand{\RSA}{\Term{RSA}\xspace}
\newcommand{\MST}{\Term{MST}\xspace}
\providecommand{\CClassPP}{\CClass{P{}P}\xspace}
\providecommand{\ID}{\mathrm{id}}

\newcommand{\EncodeFunc}{\text{\ttfamily{Enc}}}
\newcommand{\DecodeFunc}{\text{\ttfamily{Dec}}}
\newcommand{\discX}[1]{\mathrm{disc}\pth{#1}}%

\newcommand{\propertyC}{\psi}%
\newcommand{\propertyX}[1]{\propertyC\pth{#1}}%
\newcommand{\numC}{\Gamma}%

\newcommand{\degreeX}[1]{\mathrm{degree}\pth{#1}}

\newcommand{\ind}{m}

\providecommand{\nequiv}{\not\equiv}
\providecommand{\sqcupplus}{\sqcup^+}

\newcommand{\AlgorithmI}[1]{{\protect\color{RedViolet}{#1}}}

\usepackage{marginnote}
\usepackage{ragged2e} 
% Short, subtle snark macro
\setlength{\marginparwidth}{1.5in}
\setlength{\marginparsep}{10pt}
\newcommand{\snark}[1]{%
   \marginnote{\raggedright\scriptsize\sffamily\color{gray}#1}%
}


\MakePerPage{footnote} % Resets footnote counter on every page

\allowdisplaybreaks % Allows page breaks inside multiline math environments

%-----------------------------------------------------------------------------
% Bibliography & Document Begin
%-----------------------------------------------------------------------------
\IfFileExists{\jobname.bib}{%
   \addbibresource{\jobname.bib}%
}{%
   \addbibresource{shortcuts.bib}%
   \addbibresource{geometry.bib}%
}
\makeindex
\begin{document}
}{\input{prefix_lwarp}}%
\else
  \ifdefined\HCode
    \IfFileExists{styles/prefix_4ht.tex}{%-----------------------------------------------------------------------------
% TeX4ht Prefix File for Randomized Algorithms Class Notes
% Kept separate from prefix_latex.tex so HTML-specific simplifications cannot
% affect the standard PDF or lwarp builds.
%-----------------------------------------------------------------------------

\RequirePackage{iftex}
\ifPDFTeX
  \PackageError{engine-guard}{This document requires LuaLaTeX or XeLaTeX}{%
    Please recompile using LuaLaTeX or XeLaTeX instead of pdfLaTeX.%
  }
  \stop
\fi

\ifdefined\HCode\else
  \PackageError{tex4ht-gate}{%
    Attempting to use prefix_4ht.tex outside TeX4ht.%
  }{Compile through make4ht, or use prefix_latex.tex for PDF generation.}
  \stop
\fi

% Always use memoir
\documentclass[12pt,oneside,fleqn]{memoir}

%-----------------------------------------------------------------------------
% Typography, Fonts & Math Packages
%-----------------------------------------------------------------------------
\usepackage{mathtools}
\PassOptionsToPackage{warnings-off={mathtools-overbracket,mathtools-colon}}{unicode-math}
\usepackage{fontspec}
\usepackage{unicode-math}

% New Computer Modern OpenType text and math fonts
\setmainfont{NewCM10-Book.otf}[
  BoldFont          = NewCM10-Bold.otf,
  ItalicFont        = NewCM10-BookItalic.otf,
  BoldItalicFont    = NewCM10-BoldItalic.otf,
  SlantedFont       = NewCM10-Book.otf,
  SlantedFeatures   = {FakeSlant = 0.25},
  BoldSlantedFont   = NewCM10-Bold.otf,
  BoldSlantedFeatures = {FakeSlant = 0.25},
  SmallCapsFeatures = {Letters = SmallCaps}
]
\setsansfont{NewCMSans10-Book.otf}[
  BoldFont          = NewCMSans10-Bold.otf,
  ItalicFont        = NewCMSans10-BookOblique.otf,
  BoldItalicFont    = NewCMSans10-BoldOblique.otf
]
\setmonofont{NewCMMono10-Book.otf}[
  BoldFont          = NewCMMono10-Bold.otf,
  ItalicFont        = NewCMMono10-BookItalic.otf,
  BoldItalicFont    = NewCMMono10-BoldOblique.otf
]

\setmathfont{NewCMMath-Book.otf}
\setmathfont{NewCMMath-Bold.otf}[version=bold]

% OpenType math aliases and fixes
\providecommand{\square}{\mdlgwhtsquare}
\def\Box{\mdlgwhtsquare}
\def\Diamond{\mdlgwhtdiamond}
\def\blacksquare{\mdlgblksquare}
\def\leadsto{\rightsquigarrow}
\providecommand{\llbracket}{\lBrack}
\providecommand{\rrbracket}{\rBrack}
\DeclareTextCommand{\textcircled}{TU}[1]{%
  \mbox{\ooalign{\hfil\raise0.1ex\hbox{#1}\hfil\crcr\char"25EF}}%
}

\usepackage{alphabeta}
\usepackage[normalem]{ulem} % Provides \sout
\usepackage{mleftright}
\usepackage{microtype}
\usepackage{xspace}
\usepackage{xurl}
\usepackage[dvipsnames,table]{xcolor}
\usepackage{graphicx}
\usepackage{caption}
\usepackage{booktabs,array,colortbl,hhline,multirow,tabularx}
\usepackage[inline]{enumitem}
\usepackage[thicklines]{cancel}
\renewcommand\CancelColor{\color{purple}}
\usepackage{euscript}
\usepackage{marvosym}
\let\Cross\undefined
\let\Square\undefined
\usepackage{bbding}
\usepackage{epigraph}
\usepackage{makeidx}
\usepackage{perpage}
\usepackage{tikz}

%-----------------------------------------------------------------------------
% Boondox / mathcalb font declaration
%-----------------------------------------------------------------------------
\DeclareFontFamily{U}{BOONDOX-calo}{\skewchar\font=45 }
\DeclareFontShape{U}{BOONDOX-calo}{m}{n}{
  <-> s*[1.05] BOONDOX-r-calo}{}
\DeclareFontShape{U}{BOONDOX-calo}{b}{n}{
  <-> s*[1.05] BOONDOX-b-calo}{}
\DeclareMathAlphabet{\mathcalb}{U}{BOONDOX-calo}{m}{n}
\SetMathAlphabet{\mathcalb}{bold}{U}{BOONDOX-calo}{b}{n}
\DeclareMathAlphabet{\mathbcalb}{U}{BOONDOX-calo}{b}{n}

%-----------------------------------------------------------------------------
% Algorithm Environment (formerly salgorithm.sty)
%-----------------------------------------------------------------------------
\newcommand{\ProgWordX}[1]{{\small\textbf{#1}}{\ }}
\newcommand{\Do}{\ProgWordX{do}\xspace}
\newcommand{\End}{\ProgWordX{end}\xspace}
\newcommand{\Return}{\ProgWordX{return}\xspace}%
\newcommand{\Procbegin}{\ProgWordX{begin}}
\newcommand{\If}{\ProgWordX{if}\xspace}
\newcommand{\Then}{\ProgWordX{then}\xspace}
\newcommand{\Else}{\ProgWordX{else}\xspace}
\newcommand{\For}{\ProgWordX{for}\xspace}
\newcommand{\While}{{\small\bfseries while}\ }

\newcommand{\xbeginlgox}{\begin{minipage}{1in}\begin{tabbing}
           \quad\=\qquad\=\qquad\=\qquad\=\qquad\=\qquad\=\qquad\=\kill}
\newcommand{\xendlgox}{\end{tabbing}\end{minipage}}

\newenvironment{algorithmEnv}{%
   \begin{tabular}{|l|}\hline\xbeginlgox\bigskip}%
   {\xendlgox\\\hline\end{tabular}}

\newenvironment{program}{%
   \begin{minipage}{4.0in}%
   \begin{tabbing}%
       \ \ \ \ \=\ \ \ \=\ \ \ \ \=\ \ \ \ \=\ \ \ \ \=%
      \ \ \ \ \=\ \ \ \ \=\ \ \ \ \=\ \ \ \ \=%
      \ \ \ \ \=\ \ \ \ \=\ \ \ \ \=\ \ \ \ \=\kill%
}{%
   \end{tabbing}%
   \end{minipage}%
}

\newcommand{\AlgorithmBox}[1]{%
   \fbox{\begin{program} #1 \end{program}}%
}

\definecolor{AlgorithmColor}{cmyk}{0.07,0.90,0,0.34}
\providecommand{\Algorithm}[1]{{\AlgorithmI{#1}\index{algorithm!#1@{\AlgorithmI{#1}}}}}
\newcommand{\CodeComment}[1]{\textcolor{blue}{\texttt{#1}}}

%-----------------------------------------------------------------------------
% Path & Directory Resolution
%-----------------------------------------------------------------------------
\newcommand{\RootDIR}{../}
\IfFileExists{../../../book.tex}{\renewcommand{\RootDIR}{../../../}}{}
\IfFileExists{../../book.tex}{\renewcommand{\RootDIR}{../../}}{}
\IfFileExists{../book.tex}{\renewcommand{\RootDIR}{../}}{}
\IfFileExists{./book.tex}{\renewcommand{\RootDIR}{}}{}

\providecommand{\File}[1]{#1}
\providecommand{\RootFile}[1]{\RootDIR #1}

%-----------------------------------------------------------------------------
% QED Duck Symbol
%-----------------------------------------------------------------------------
\IfFileExists{styles/qed_duck.svg}{%
   \providecommand{\myqedsymbol}{%
      \HCode{<img class="qed-duck" src="styles/qed_duck.svg" alt="QED" />}%
   }%
}{%
   \providecommand{\myqedsymbol}{\HCode{<span class="qed-fallback">&#x25A0;</span>}}%
}

%-----------------------------------------------------------------------------
% BibLaTeX Setup (Pure Biber, No BibTeX)
%-----------------------------------------------------------------------------
\usepackage[bibencoding=utf8,style=alphabetic,refsegment=chapter,
indexing=cite,
backend=biber]{biblatex}

\ExecuteBibliographyOptions{%
   maxnames = 40,%
   maxalphanames = 4,%
   maxcitenames = 3,%
   minalphanames = 3,%
   minnames = 1,%
   giveninits = true,%
   doi = false,%
   url = false,%
   isbn = false%
}

\renewbibmacro{in:}{}

\AtEveryBibitem{%
   \ifentrytype{inproceedings}{%
      \clearname{editor}%
      \clearlist{editor}%
      \clearfield{editor}%
      \clearlist{publisher}%
      \clearfield{series}%
      \clearlist{location}%
      \clearfield{publisher}%
   }{}%
   \ifentrytype{book}{%
      \clearname{publisher}%
   }{}%
}
\AtEveryCitekey{%
   \ifentrytype{inproceedings}{%
      \clearlist{publisher}%
   }{}%
   \clearname{editor}%
   \clearlist{location}%
}
\DeclareFieldFormat[inproceedings]{date}{}
\DeclareFieldFormat[inproceedings]{pages}{#1, \printfield{year}}
\DeclareFieldFormat[inproceedings]{booktitle}{\mkbibemph{#1,\nopunct}}

\DeclareFieldFormat{journaltitle}{\mkbibemph{#1}\isdot}

\DeclareFieldFormat[article]{number}{\mkbibparens{#1}}
\DeclareFieldFormat[article]{date}{}
\DeclareFieldFormat[article]{pages}{#1, \printfield{year}}
\DeclareFieldFormat[article]{issue}{\nopunct#1}
\renewbibmacro*{issue+date}{%
   \iffieldundef{issue}{}%
   {%
      \printtext{\unspace}%
      (\printfield{issue})%
   }%
   \iffieldundef{pages}%
   {\addcomma\printfield{year}}%
   {}%
}

\renewcommand{\bibpagespunct}{%
   \ifentrytype{article}{\addcolon}{\addcomma}\space%
}

\DeclareFieldFormat[article]{titlecase}{\MakeSentenceCase{#1}}
\DeclareFieldFormat[inproceedings]{titlecase}{\MakeSentenceCase{#1}}
\DeclareFieldFormat[incollection]{titlecase}{\MakeSentenceCase{#1}}
\DeclareFieldFormat[book]{titlecase}{\MakeSentenceCase{#1}}

\renewbibmacro*{journal}{%
   \iffieldundef{journaltitle} {} {\printtext[journaltitle]{%
         \printfield[myplain]{journaltitle}%
         \setunit{\subtitlepunct}%
         \printfield[myplain]{journalsubtitle}}}}
\DeclareFieldFormat{myplain}{#1}

\renewbibmacro*{booktitle}{%
   \ifboolexpr{ test {\iffieldundef{booktitle}} and test
      {\iffieldundef{booksubtitle}} } {} {\printtext[booktitle]{%
         \printfield[myplain]{booktitle}%
         \setunit{\subtitlepunct}%
         \printfield[myplain]{booksubtitle}}%
      \newunit}%
   \printfield{booktitleaddon}}

\renewbibmacro*{volume+number+eid}{%
   \setunit*{\addcomma\addspace}%
   \printfield{volume}%
   \printfield{number}%
   \setunit{\addcomma\space}%
   \printfield{eid}%
}

\newbibmacro{string+doi}[1]{%
   \iffieldundef{doi}{%
      \iffieldundef{url}{%
         #1%
      }{%
         \href{\thefield{url}}{#1}%
      }%
   }{%
      \href{http://dx.doi.org/\thefield{doi}}{#1}}%
}
\DeclareFieldFormat*{title}{\usebibmacro{string+doi}{\mkbibemph{#1}}}
\renewcommand*{\multicitedelim}{\addcomma\space}
\DeclareFieldFormat[inproceedings]{title}{\usebibmacro{string+doi}{{#1}}}
\DeclareFieldFormat[article]{title}{\usebibmacro{string+doi}{{#1}}}
\DeclareFieldFormat[incollection]{title}{\usebibmacro{string+doi}{{#1}}}
\DeclareFieldFormat[book]{titlecase}{{#1}}

\DeclareCiteCommand{\tcite}%
{\usebibmacro{prenote}}%
{\usebibmacro{citeindex}\usebibmacro{cite}}%
{\multicitedelim}%
{\usebibmacro{postnote}}

%-----------------------------------------------------------------------------
% Hyperref & Visual Formatting
%-----------------------------------------------------------------------------
\usepackage{hyperref}
\hypersetup{%
   breaklinks,%
   colorlinks=true,%
   linkcolor=[rgb]{0.45,0.0,0.0},%
   citecolor=[rgb]{0,0,0.45}%
}

\definecolor{DeepNavy}{HTML}{002366}
\definecolor{RoyalBlue}{HTML}{004E92}
\definecolor{SteelBlue}{HTML}{4682B4}
\definecolor{TealBlue}{HTML}{008080}
\definecolor{DarkSlate}{HTML}{2F4F4F}
\definecolor{DustyGrey}{HTML}{696969}

\makeatletter
\AtBeginDocument{\let\@elt\relax}
\makeatother

\headstyles{bringhurst}
\chapterstyle{dash}

\renewcommand{\chapnamefont}{\normalfont\huge\bfseries\color{DeepNavy}}
\renewcommand{\chapnumfont}{\normalfont\huge\bfseries\color{DeepNavy}}
\renewcommand{\chaptitlefont}{\normalfont\Huge\bfseries\color{DeepNavy}}

\counterwithin{footnote}{chapter}

\setsubparaheadstyle{\normalfont\normalsize\itshape\color{DustyGrey}}
\setlength{\beforechapskip}{-20pt} 
\setlength{\afterchapskip}{30pt}
\setsecnumformat{\csname the#1\endcsname.\quad}

%-----------------------------------------------------------------------------
% Standalone Chapter vs Master Book Modes
%-----------------------------------------------------------------------------
\newcommand{\Copyright}{This work is licensed under the Creative Commons Attribution-Noncommercial 3.0 License. To view a copy of this license, visit \url{http://creativecommons.org/licenses/by-nc/3.0/} or send a letter to Creative Commons, 171 Second Street, Suite 300, San Francisco, California, 94105, USA.}

\newcommand{\RealChapterEndInner}{%
   {%
      \begingroup%
      \emergencystretch=1em%
      \printbibliography[heading=subbibliography,segment=\therefsegment]{}%
      \endgroup%
   }%
}

\providecommand{\BookMode}[1]{}
\providecommand{\StandAloneMode}[1]{#1}

\ifx\bookMode\undefined
   % STANDALONE CHAPTER MODE
   \renewcommand{\StandAloneMode}[1]{#1}
   \renewcommand{\BookMode}[1]{}

   \newcommand{\ChapterNumPage}[2]{%
      \if\relax\detokenize{#1}\relax
         \setcounter{chapter}{0}%
      \else\ifnum#1>0
         \setcounter{chapter}{\numexpr#1-1\relax}%
      \else
         \setcounter{chapter}{0}%
      \fi\fi
      \if\relax\detokenize{#2}\relax
         \setcounter{page}{1}%
      \else\ifnum#2>0
         \setcounter{page}{#2}%
      \else
         \setcounter{page}{1}%
      \fi\fi
   }

   \DeclareRobustCommand{\setheadingcolor}{\color{RoyalBlue}}
   \setsecheadstyle{\setheadingcolor\normalfont\Large\bfseries}
   \setsubsecheadstyle{\normalfont\large\bfseries\color{SteelBlue}}
   \setsubsubsecheadstyle{\normalfont\normalsize\bfseries\color{TealBlue}}
   \setparaheadstyle{\normalfont\normalsize\bfseries\color{DarkSlate}}

   \DeclareDocumentCommand{\Chapter}{s o m}{%
      \IfBooleanTF{#1}{%
         \IfValueTF{#2}{\chapter*[]{#2}}{\chapter*[]{#3}}%
      }{%
         \IfValueTF{#2}{\chapter*[]{#2}}{\chapter[]{#3}}%
      }%
      \setcounter{footnote}{0}%
   }
   \newcommand{\ChapterPrefix}{}
   \newcommand{\ChapterGen}{\pagestyle{plain}\vspace*{-0.7cm}}
   \newcommand{\ChapterEnd}{%
      \RealChapterEndInner{}%
      \end{document}%
   }

   \renewcommand{\File}[1]{#1}
   \renewcommand{\RootFile}[1]{../#1}
\else
   % MASTER BOOK MODE
   \renewcommand{\StandAloneMode}[1]{}
   \renewcommand{\BookMode}[1]{#1}

   \newcommand{\ChapterNumPage}[2]{}

   \DeclareDocumentCommand{\Chapter}{s o m}{%
      \IfBooleanTF{#1}{%
         \IfValueTF{#2}{\chapter*[#2]{#3}}{\chapter*{#3}}%
         \ifx\Dir\undefined\else
            \typeout{CHAPTER_MAP: \Dir/\CurrentChapFile => chap: 0, page: \thepage}%
         \fi
      }{%
         \IfValueTF{#2}{\chapter[#2]{#3}}{\chapter{#3}}%
         \ifx\Dir\undefined\else
            \typeout{CHAPTER_MAP: \Dir/\CurrentChapFile => chap: \thechapter, page: \thepage}%
         \fi
      }%
   }
   \newcommand{\ChapterPrefix}{%
      \vspace{-1cm}%
      \subsection*{\small
         {CS 574: Randomized Algorithms}\\
         {Sariel Har-Peled}\\
         {\today}}%
   } \newcommand{\ChapterGen}{}
   \newcommand{\ChapterEnd}{\RealChapterEndInner{}}

   \renewcommand{\File}[1]{\Dir/#1}
   \renewcommand{\RootFile}[1]{#1}

   \newcommand{\InputChap}[2]{%
      \def\Dir{#1}%
      \def\CurrentChapFile{#2}%
      \typeout{}%
      \typeout{------------------------------------------------------------}%
      \typeout{PROCESSING FILE: #1/#2}%
      \typeout{}%
      \input{#1/#2}%
   }
\fi

%-----------------------------------------------------------------------------
% Fragments & Quotes
%-----------------------------------------------------------------------------
% The PDF build writes fragment bodies to shared files and later includes them
% in standalone appendices.  A chapter webpage already contains each fragment
% at its original location, so TeX4ht only needs a transparent wrapper and can
% omit the repeated standalone copies.
\DeclareDocumentEnvironment{fragment}{m +b}{%
   #2%
}{}
\DeclareDocumentCommand{\IncFragment}{m}{}

\newcommand{\QuoteOpen}{}
\newcommand{\QuoteClose}{}
\newcommand{\QuoteWrite}[2]{}

\newcommand{\Quote}[2]{%
   \SaveIndent
   \QuoteWrite{#1}{#2}%
   \begin{center}%
      {\footnotesize
         \begin{quote}
             \RestoreIndent
             {#1}\\
             \hspace*{2cm}-- {#2.}
         \end{quote}%
      }%
      \index{quotation!#2}%
   \end{center}%
   \bigskip
}

\newcommand{\WebChapterQuote}[3]{%
   \ifvmode\IgnorePar\fi\EndP%
   \HCode{<figure class="chapter-quote"><blockquote>}%
      #1%
   \HCode{</blockquote><hr /><figcaption><span class="quote-author">}%
      #2%
   \HCode{</span><span class="quote-work">}%
      #3%
   \HCode{</span></figcaption></figure>}%
   \par%
}

\newcommand{\QuotePExt}[3]{\WebChapterQuote{#1}{#2}{#3}}

\epigraphfontsize{\small\itshape}
\setlength{\epigraphwidth}{0.6\textwidth}
\newcommand{\QuoteExt}[3]{\WebChapterQuote{#1}{#2}{#3}}

\providecommand{\delX}[1]{{\color{red}\sout{#1}}}
\providecommand{\newX}[1]{{\color{blue}#1}}
\providecommand{\chgY}[2]{{\color{red}\sout{#1}}{\color{blue}#2}}
\providecommand{\remX}[1]{{\color{purple}\textbf{[Remark: #1]}}}
\newcommand{\SpellIgnore}[1]{#1}
\newcommand{\EQ}{\mathrel{\;\;=\;\;}}
\newcommand{\LEQ}{\mathrel{\;\;\leq\;\;}}
\newcommand{\LT}{\mathrel{\;\;<\;\;}}
\newcommand{\GT}{\mathrel{\;\;<\;\;}}

\newcommand{\GEQ}{\mathrel{\;\;\geq\;\;}}

%-----------------------------------------------------------------------------
% Save a symbol that we know is going to get redefined.
\def\savesymbol#1{%
  \expandafter\let\expandafter\origsym\expandafter=\csname#1\endcsname
  \expandafter\let\csname orig#1\endcsname=\origsym
  \expandafter\let\csname#1\endcsname=\relax
}

% Restore a previously saved symbol, and rename the current one.
\def\restoresymbol#1#2{%
  \expandafter\let\expandafter\newsym\expandafter=\csname#2\endcsname
  \expandafter\global\expandafter\let\csname#1#2\endcsname=\newsym
  \expandafter\let\expandafter\origsym\expandafter=\csname orig#2\endcsname
  \expandafter\global\expandafter\let\csname#2\endcsname=\origsym
}

\usepackage{pictex}%
\savesymbol{triangle}%
\usepackage{tree}%
\restoresymbol{TreePKG}{triangle}%

% PicTeX's tree macros perform box arithmetic that TeX4ht cannot convert and
% can terminate the HTML document at the first tree.  These lightweight web
% forms preserve the two children and optional subtree label as HTML/CSS.
\renewcommand{\tree}[2]{%
   \HCode{<span class="web-binary-tree"><span class="tree-node"></span><span class="tree-children"><span class="tree-child">}%
   #1%
   \HCode{</span><span class="tree-child">}%
   #2%
   \HCode{</span></span></span>}%
}
\renewcommand{\TreePKGtriangle}[3][]{%
   \HCode{<span class="web-subtree-triangle"><span class="subtree-label">}%
   #1%
   \HCode{</span></span>}%
}

% Definitions absorbed from sbook.sty
\providecommand{\MdPrc}{\nu}

\providecommand{\TwoFiguresBase}[4]{%
   \begin{tabular}{cc}
      \begin{minipage}{0.48\linewidth}
          \begin{center}
              {{{#1}}}
          \end{center}
      \end{minipage}
      &
      \begin{minipage}{0.48\linewidth}
          \begin{center}
              {{#3}}
          \end{center}
      \end{minipage}
      \\
      {#2} & {#4}
   \end{tabular}%
}

\providecommand{\TwoFigures}[6]{%
   \begin{tabular}{cc}
      \begin{minipage}{0.48\linewidth}
          \begin{center}
              {\IncludeGraphics[#1]{#2}}
          \end{center}
      \end{minipage}
      &
      \begin{minipage}{0.48\linewidth}
          \centerline{\IncludeGraphics[#4]{#5}}
      \end{minipage}
      \\
      {#3} & {#6}
   \end{tabular}%
}

\providecommand{\ThreeFigures}[6]{%
   \begin{tabular}{ccc}
        \begin{minipage}{0.30\linewidth}
            \IncludeGraphics[width=5.1cm]{#1}\\
        \end{minipage}
        &
        \begin{minipage}{0.30\linewidth}
            \IncludeGraphics[width=5.1cm]{#2}\\
        \end{minipage}
        &
        \begin{minipage}{0.30\linewidth}
            \IncludeGraphics[width=5.1cm]{#3}\\
        \end{minipage}
        \vspace{-1cm}
        \\
        {#4} & {#5}& {#6}
   \end{tabular}%
}

%-----------------------------------------------------------------------------
\providecommand{\remove}[1]{}
\providecommand{\NotAMSMode}[1]{#1}
\providecommand{\Solution}[1]{}
\providecommand{\Version}[1]{\hfill{Version: #1}}
\providecommand{\Distor}{\mathrm{dist}}
\providecommand{\Hardness}[1]{(h:#1)}
\providecommand{\QuoteOptExt}[3]{\QuoteExt{#1}{#2}{#3}}
\providecommand{\Hypercube}{\EuScript{H}}
\providecommand{\CoverTime}{\EuScript{C}}
\providecommand{\Nbr}{\Gamma}
\providecommand{\CommuteTime}[2]{\mathbf{CT}_{#1#2}}
\providecommand{\Resistance}{\mathbf{R}}
\providecommand{\EffectResist}[2]{\mathbf{R}_{#1#2}}
\providecommand{\NetworkE}{\EuScript{N}}
\providecommand{\trinom}[4]{\binom{#1}{#2 \;\; #3 \;\; #4}}
\providecommand{\Ranges}{\EuScript{R}}
\providecommand{\CClass}[1]{\ensuremath{\mathbf{#1}}}
\providecommand{\coCClass}[1]{\ensuremath{\mathrm{co-}}\CClass{#1}}
\providecommand{\CX}{\EuScript{C}}
\providecommand{\packet}{\mathbf{c}}
\providecommand{\packetA}{\mathbf{d}}
\providecommand{\mopt}{m_{\mathrm{opt}}}
\providecommand{\girth}{\mathop{\mathrm{girth}}}
\providecommand{\CR}{\mathop{\mathrm{cr}}}
\providecommand{\TRUE}{{\textcolor[named]{RedViolet}{\textsf{TRUE}}}\xspace}
\providecommand{\FALSE}{{\textcolor[named]{RedViolet}{\textsf{FALSE}}}\xspace}
\providecommand{\Part}[1]{\part{#1}}
\providecommand{\Section}[1]{\section{#1}}
\providecommand{\SubSection}[1]{\subsection{#1}}
\providecommand{\SubSubSection}[1]{\subsubsection{#1}}
\providecommand{\Exercise}[5]{%
\begin{exercise_h}[#1]%
   \exelab{#3}{}%
   #4%
\end{exercise_h}%
}
\providecommand{\si}[1]{#1}
\providecommand{\myParagraph}[1]{\noindent\textbf{#1}}
\providecommand{\DefName}[1]{#1}
\providecommand{\deflabel}[1]{#1}
\providecommand{\Description}[1]{}

%-----------------------------------------------------------------------------
% Mathematical Notations & Environments (from notations.sty)
%-----------------------------------------------------------------------------
%-----------------------------------------------------------------------------
% List environments (questions & compactenum)
%-----------------------------------------------------------------------------

\newlist{compactenumA}{enumerate}{5}%
\setlist[compactenumA]{topsep=0pt,itemsep=-1ex,partopsep=1ex,parsep=1ex,%
   label=(\Alph*)}%

\newlist{compactenumI}{enumerate}{5}%
\setlist[compactenumI]{topsep=0pt,itemsep=-1ex,partopsep=1ex,parsep=1ex,%
   label=(\Roman*)}%

\newlist{compactenuma}{enumerate}{5}%
\setlist[compactenuma]{topsep=0pt,itemsep=-1ex,partopsep=1ex,parsep=1ex,%
   label=(\alph*)}%

\newlist{compactitem}{enumerate}{5}%
\setlist[compactitem]{topsep=0pt,itemsep=-1ex,partopsep=1ex,parsep=1ex,%
   label=\textbullet}%

\newlist{compactenumi}{enumerate}{5}%
\setlist[compactenumi]{topsep=0pt,itemsep=-1ex,partopsep=1ex,parsep=1ex,%
   label=(\roman*)}%

\setcounter{secnumdepth}{5}

\newcommand{\HST}{\TermI{HST}\xspace}
\newcommand{\Term}[1]{\textsf{#1}}
\newcommand{\TermI}[1]{\Term{#1}\index{#1@\Term{#1}}}

\DeclareMathAlphabet{\mathpzc}{OT1}{pzc}{m}{it}

\newcommand{\seg}{\mathsf{s}}

\newcommand{\Diam}{\mathrm{diam}}
\newcommand{\diam}{\mathrm{diam}}

\newcommand{\Copt}{C_{\mathrm{opt}}}
\newcommand{\CenterPnt}{\mathrm{cen}}

\newcommand{\parent}{\overline{\mathrm{p}}}
\newcommand{\parentX}[2][\!]{\parent\pth[#1]{#2}}

\newcommand{\FMS}{\EuScript{X}\index{metric space}}  
\newcommand{\FMSB}{\EuScript{Y}\index{metric space}}  
\newcommand{\FastCut}{\Algorithm{FastCut}\xspace}
\newcommand{\MinCut}{\Algorithm{MinCut}\xspace}
\newcommand{\MinCutRep}{\Algorithm{MinCutRep}\xspace}
\newcommand{\Contract}{\Algorithm{Contract}\xspace}

\newcommand{\Grid}{\mathrm{G}}
\newcommand{\ClosestPair}{\mathcal{CP}}
\newcommand{\SphereChar}{\ensuremath{\mathbb{S}}}
\newcommand{\Sphere}[1]{{\SphereChar^{(#1)}}}

\newcommand{\PntSet}{{\mathsf{P}}}%
\providecommand{\PntSetA}{{\mathsf{Q}}}%

\newcommand{\query}{\mathcalb{q}}

\newcommand{\pnt}{\mathsf{p}}
\newcommand{\pntA}{\mathsf{u}}
\newcommand{\pntB}{\mathsf{v}}
\newcommand{\pntC}{{\mathsf{x}}}
\newcommand{\pntD}{{\mathsf{y}}}
\newcommand{\pntE}{{\mathsf{z}}}

\newcommand{\pntCoord}{p}
\newcommand{\pntCoordA}{q}

%%%%%%%%%%%%%%%%%%%%%%%%%%%%%%%%%%%%%%%%%%%%%%%%%%%%%%%%%%%%%%%%%%%%% 
%%%%%%%%%%%%%%%%%%%%%%%%%%%%%%%%%%%%%%%%%%%%%%%%%%%%%%%%%%%%%%%%%%%%%
% Markov chain notations.

\newcommand{\MCSpc}{\mathtt{S}}
\newcommand{\MCPrb}{\mathtt{P}}
\newcommand{\MCPrbY}[2]{\MCPrb_{#1#2}}
\newcommand{\MCPrbM}{\mathbf{P}}%
\newcommand{\MCPrbT}[3]{\MCPrb_{#1#2}^{(#3)}}
\newcommand{\MCFPrbT}[3]{\mathtt{r}_{#1#2}^{(#3)}}
\newcommand{\MCVstProb}[2]{\mathtt{f}_{#1#2}}
\newcommand{\MCExSteps}[2]{\mathtt{h}_{#1#2}}
\newcommand{\MCPVec}[1]{\mathbf{q}^{(#1)}}
%%%%%%%%%%%%%%%%%%%%%%%%%%%%%%%%%%%%%%%%%%%%%%%%%%%%%%%%%%%%%%%%%%%%%

\newcommand{\Spread}{\Phi}
\newcommand{\VC}{\ensuremath{\mathsf{VC}}\xspace}
\newcommand{\MC}{\ensuremath{\mathsf{MC}}\xspace}
\newcommand{\hyperplane}{h}

\newcommand{\QuickSort}{\Algorithm{QuickSort}\xspace}
\newcommand{\Alg}{\Algorithm{Alg}\xspace}

\newcommand{\Graph}{\mathsf{G}}%
\newcommand{\GraphB}{\mathsf{H}}
\newcommand{\GraphC}{\mathsf{K}}
\newcommand{\GraphD}{\mathsf{F}}

\providecommand{\G}{}
\renewcommand{\G}{\Graph}

\newcommand{\RLP}{\ComplexityClass{RLP}\xspace}
\newcommand{\USTCON}{\ComplexityClass{USTCON}}
\newcommand{\STCON}{\overrightarrow{\ComplexityClass{STCON}}}

\newcommand{\Mtrx}{\mathsf{M}}
\newcommand{\TMtrx}{\mathsf{Q}}
\newcommand{\IMtrx}{\mathcal{I}}

\newcommand{\Expander}[3]{$\left[#1,#2,#3\right]$-expander}

\newcommand{\ChartX}[1]{\ensuremath{\chi_{#1}}}
\newcommand{\Field}{\mathbb{F}}
\newcommand{\LSpace}{\mathbb{G}}
\newcommand{\LD}[2]{\mathsf{L{D}}\pth{#1,#2}}
\newcommand{\SEVN}{\widehat{\lambda_2}}
\newcommand{\NX}[2]{#1\!\left[#2\right]}
\newcommand{\Entry}[2]{#1\!\left[#2\right]}
\newcommand{\NXExt}[3]{#1_{#2}\!\left[#3\right]}
\newcommand{\RangeI}[1]{\llbracket #1 \rrbracket}
\newcommand{\DialX}[1]{\gamma\pth{#1}}
\newcommand{\BiDialX}[1]{{\gamma}_2\pth{#1}}

\newcommand{\EV}[1]{\widehat{\lambda_{#1}}}
\newcommand{\TEV}[1]{\widehat{\lambda_{#1}}}
\newcommand{\EVal}{\widehat{\lambda}}
\newcommand{\Basis}{\mathcal{B}}

\newcommand{\ZProd}[2]{\ensuremath{#1 \, \text{\textcircled{z}} #2}}
\newcommand{\RProd}[2]{\ensuremath{#1 \, \text{\textcircled{r}} #2}}

\newcommand{\OneVec}[1]{\mathbf{1}^{#1}}
\newcommand{\DstrbV}[1]{\mathbf{\mathsf{q}}^{(#1)}}
\newcommand{\DstrbE}[2]{\mathsf{q}^{(#1)}_{#2}}

\newcommand{\DEdges}{\widetilde{\Edges}}

\newcommand{\VV}{\mathsf{V}}%
\newcommand{\EE}{\mathsf{E}}%

\newcommand{\Vertices}{\VV}%
\newcommand{\Edges}{\EE}%

\newcommand{\EdgesX}[1]{\Edges\pth{#1}}
\newcommand{\VerticesX}[1]{V\pth{#1}}

\newcommand{\StringX}[1]{\mathsf{s}\pth{#1}}

\newcommand{\BMtrx}{\mathsf{B}}
\newcommand{\WMtrx}{\mathsf{W}}
\newcommand{\NWMtrx}{\overline{\mathsf{W}}}
\newcommand{\Sequence}{\EuScript{S}}

\newcommand{\med}{\mathrm{med}}
\newcommand{\medX}[1]{\mathrm{med}\pth{#1}}

\newcommand{\medianC}{\nu}%

\newlength{\savedparindent}
\newcommand{\SaveIndent}{\setlength{\savedparindent}{\parindent}}
\newcommand{\RestoreIndent}{\setlength{\parindent}{\savedparindent}}

\newcommand{\NNN}{\mathbb{N}}%
\providecommand{\NN}{\mathbb{N}}%
\renewcommand{\Re}{\mathbb{R}}%
\newcommand{\reals}{\Re}%
\newcommand{\ZZ}{\mathbb{Z}}%

\newcommand{\level}{\mathop{\mathrm{level}}}

\providecommand{\TPDF}[2]{\texorpdfstring{#1}{#2}}

\newcommand{\Hint}[1]{\emph{[#1]}}
\newcommand{\Arr}{\mathop{\mathrm{\EuScript{A}}}}

\newcommand{\Array}{\mathcal{T}}%

\newcommand{\dirEdge}[2]{\left({{#1} \rightarrow {#2}} \right)}

%%%%%%%%%%%%%%%%%%%%%%%%%%%%%%%%%%%%%%%%%%%%%%%%%%%%%%%%%%%%%%%%%%%%%%
% Emphasize stuff...
\definecolor{blue25}{rgb}{0, 0, 11}
\newcommand{\emphic}[2]{%
%   \BADVERYBAD%
   \textcolor{blue25}{%
      \bfseries{\emph{#1}}}%
   \index{#2}}
\newcommand{\emphi}[1]{\emphic{#1}{#1}}
\newcommand{\emphOnly}[1]{\textcolor{blue25}{\bfseries{\emph{#1}}}}
\providecommand{\emphiC}[1]{\emphic{#1}{#1}}


\definecolor{almostblack}{rgb}{0, 0, 0.3}
\providecommand{\emphw}[1]{{\textcolor{almostblack}{\emph{#1}}}}%

%
%%%%%%%%%%%%%%%%%%%%%%%%%%%%%%%%%%%%%%%%%%%%%%%%%%%%%%%%%%%%%%%%%%%%%%

\newcommand{\volX}[1]{\mathrm{vol}\pth{#1}}

\newcommand{\Interval}{\EuScript{I}}
\newcommand{\IntervalA}{\EuScript{I}}

\newcommand{\Exp}{\mathrm{Exp}}

\let\Ball\undefined
\newcommand{\Ball}{\mathbf{b}}
\newcommand{\hdim}{d}
\newcommand{\Domain}{\mathcal{D}}
\newcommand{\points}[1]{[\bfseries{#1 points}]}
\newcommand{\Event}{\EuScript{E}}%
\newcommand{\EventB}{\EuScript{B}}%
\newcommand{\EventF}{\EuScript{F}}%
\newcommand{\EventG}{\EuScript{G}}%

\newcommand{\distCmd}[1]{\left\| {#1}  \right\|}
\newcommand{\distY}[2]{\distCmd{#1 - #2}}

\newcommand{\NormalC}{\mathbf{N}}
\newcommand{\NormalY}[2]{\NormalC\pth{#1, #2}}%
\newcommand{\NormalDistOneDim}{\NormalC}
\newcommand{\NormalDistX}[1]{\NormalC^{#1}}
\DeclareMathOperator{\lca}{lca}
\newcommand{\RangeSpace}{\mathsf{S}}%
\newcommand{\RangeSpaceA}{\mathsf{T}}

\newcommand{\range}{\mathbf{r}}
\newcommand{\rangeA}{\mathbf{r}'}

\newcommand{\BadProb}{\varphi}

\newcommand{\EntropyChar}{\mathbb{H}}
\newcommand{\Entropy}[1]{\EntropyChar\pth{#1}}
\newcommand{\EntropyX}[1]{\EntropyChar\pth[]{#1}}

\newcommand{\Matrix}{\mathsf{M}}
\newcommand{\MatrixB}{\mathsf{B}}
\newcommand{\MatrixC}{\mathsf{C}}
\newcommand{\MatrixD}{\mathsf{D}}

\renewcommand{\th}{th\xspace}

\newcommand{\Family}{\mathcal{F}}
\newcommand{\FamilyB}{\mathcal{G}}
\newcommand{\FamilyC}{\mathcal{H}}
\newcommand{\FMSC}{\EuScript{Y}\index{metric space}}

\newcommand{\median}{\overline{\mathsf{m}}}
\newcommand{\QuickSelect}{\Algorithm{Quick{}Select}\xspace}

\providecommand{\row}{\mathbf{r}}%
\newcommand{\rowc}{{r}}%

\newcommand{\vecb}{\mathbf{b}}

\newcommand{\vecA}{\mathbf{v}}
\newcommand{\vecB}{\mathbf{u}}
\newcommand{\vecRand}{\mathbf{r}}

\newcommand{\vecH}{\mathbf{h}}
\newcommand{\vecV}{\mathbf{v}}

\newcommand{\vecX}{\mathbf{x}}
\newcommand{\vecY}{\mathbf{y}}

\newcommand{\DistBinom}[2]{\mathrm{Bin}\pth{#1,#2}}

\newcommand{\DistGeom}[1]{\mathrm{Geom}\pth{#1}}

\newcommand{\Sample}{\mathsf{R}}

\newcommand{\lcm}{\mathop{\mathrm{lcm}}}
\newcommand{\bdiv}{\mathop{\mathrm{div}}}%
\newcommand{\divides}{\mid}%
\newcommand{\poly}{\mathrm{poly}}%

\newcommand{\Group}{\mathcal{G}} 
\newcommand{\GroupA}{\mathcal{H}} 

\newcommand{\algGCD}{\Algorithm{EuclidGCD}\xspace}%

\newcommand{\gOp}{\times}%
\newcommand{\gIdent}{\mathsf{i}}%

\newcommand{\ordX}[2][\!]{\mathop{\mathrm{o{r}d}}\pth[#1]{#2}}

\newcommand{\lSym}[2]{%
   \left(
       \ds {#1}\divides{#2}%
   \right) %
}
\newcommand{\jSym}[2]{\left\llbracket%
       \ds {#1}\divides{#2}%
       \right\rrbracket%
}
\newcommand{\Line}{\ell}

\newcommand{\LevelChar}{\mathbb{L}}
\newcommand{\AtMostKLevelChar}{\LevelChar_{\leq k}}
\newcommand{\AtMostKLevel}[1]{\LevelChar_{\leq k} \pth{ #1}}

\newcommand{\LineSet}{\mathsf{L}}%
\newcommand{\ArrX}[1]{\EuScript{A}\pth{#1}}%

\providecommand{\Matousek}{Matou{\v s}ek\xspace}
\providecommand{\Erdos}{Erd{\H o}s\xspace}
\providecommand{\Lovasz}{Lov\'asz\xspace}

\newcommand{\ProblemC}[1]{{\textsf{\textcolor{red}{#1}\index{problem!#1}}}}
\DefineNamedColor{named}{OliveGreen} {cmyk}{0.64,0,0.95,0.40}
\providecommand{\ComplexityClass}[1]{{{\textcolor[named]{OliveGreen}{%
            \textsc{#1}}}}}
\providecommand{\NPHard}{{\ComplexityClass{NP-Hard}}%
   \index{NP!hard}\xspace}
\providecommand{\NPComplete}{\ComplexityClass{NP-Complete}%
   \index{NP!complete}\xspace}
\providecommand{\APXHard}{\ComplexityClass{APX-Hard}%
   \index{APX!Hard}\xspace}
\providecommand{\NPC}{\ComplexityClass{NPC}%
   \index{NP!complete}\xspace}
\providecommand{\NP}{\ComplexityClass{NP}%
   \index{NP}\xspace}
\providecommand{\POLYT}{\ComplexityClass{P}\xspace}

\newcommand{\envFrame}[1]{\begin{center}\fbox{\begin{minipage}[c]{0.75\textwidth}{#1}\end{minipage}}\end{center}}

\newcommand{\mysqrt}[1]{\ts\!\!\sqrt{#1}}
\newcommand{\SegSet}{\mathsf{S}}
\newcommand{\ArrVD}[1]{\Arr^{|} \pth{#1}}
\newcommand{\RIC}{\textsf{RIC}\index{random incremental
      construction}\xspace}
\newcommand{\ArrB}{\EuScript{B}}

\newcommand{\TreeOpt}{\EuScript{T}_\mathrm{opt}}
\newcommand{\TreeOptA}{\EuScript{T}_\mathrm{opt}'}

\newcommand{\DepthExt}[2]{\mathrm{depth}_{#1}\pth{#2}}
\newcommand{\depthX}[1]{\mathrm{depth}\pth{#1}}

\newcommand{\TreeA}{\EuScript{U}}
\newcommand{\TreeB}{\EuScript{V}}

\newcommand{\Trap}{\sigma}
\newcommand{\TrapA}{\sigma'}
\newcommand{\TrapB}{\sigma''}

\newcommand{\SampleA}{\mathsf{R}'}

\newcommand{\AllRegions}{\mathcal{T}}
\newcommand{\RegionsSet}[1]{\mathcal{F}\pth{#1}}
\newcommand{\RegionsHeavy}[2]{\mathcal{F}_{\!\geq#1}\pth{#2}}
\newcommand{\NumOfRegionsZero}[1]{\ExChar \!f\pth{#1}}
\newcommand{\NumOfRegions}[2]{\ExChar\! f_{\!\geq #1}\pth[]{#2}}
\newcommand{\ObjSet}{\mathsf{S}}

\newcommand{\WeightX}[2][\!]{\omega \pth[#1]{#2}}
\newcommand{\Edge}{\mathsf{e}}

\newcommand{\Face}{\mathsf{f}}

\newcommand{\DefSet}[1]{D\pth{#1}}
\newcommand{\KillSet}[1]{K\pth{#1}}

\newcommand{\CSProbFull}[4]{
   \,\rho_{\!\!\stackrel{\rule[-.005cm]{0cm}{0.1cm}}{#1,
       #4}}\pth{\MakeSBig #2,#3}}

\newcommand{\disk}{D}
\newcommand{\diskA}{D'}

\newcommand{\IntervalSet}{\widehat{\mathcal{I}}}

\newcommand{\GroundSet}{\textsf{X}}
\newcommand{\FGroundSet}{\mathsf{x}}
\newcommand{\RangeSet}{{\mathcal{R}}}
\newcommand{\RangeSetA}{\mathcal{R}'}
\newcommand{\MeasureChar}{\overline{m}}
\newcommand{\Measure}[1]{\MeasureChar\pth{#1}}
\newcommand{\sMeasure}[1]{\overline{s}\pth{#1}}

\newcommand{\SetA}{S}%
\newcommand{\SetB}{T}%
\newcommand{\SetC}{N}%

\newcommand{\VCProj}[2]{#1_{|#2}}
\newcommand{\VCDim}{\mathrm{dim}_{\VC}}
\newcommand{\Dim}{\delta}
\newcommand{\DimA}{\delta'}
\newcommand{\PntSetB}{{{\mathsf{Q}}}}
\newcommand{\PntSetC}{{{\mathsf{R}}}}
\newcommand{\Growth}[2]{{\EuScript{G}}_{#1}\pth{#2}}
\newcommand{\ShatterFunc}[1]{\pi_{#1}}
\newcommand{\Distribution}{\EuScript{D}}
\newcommand{\Error}{\EuScript{E}}

\newcommand{\bd}{{\partial}}%
\newcommand{\eps}{{\varepsilon}}%

\newcommand{\clmlab}[1]{\label{claim:#1}}
\newcommand{\clmref}[1]{\HLink{Claim}{claim:#1}}

\newcommand{\CHX}[1]{{\mathcal{CH}}\pth{#1}}

\newcommand{\ts}{\hspace{0.6pt}}

\def\qedsymbol{\myqedsymbol}%

\newcommand{\qedsign}{\hfill{\hfill\myqedsymbol}}
\newcommand{\qedsignBig}{\hfill{\hfill\rule{4mm}{4mm}}}
\newcommand{\QED}{{\hfill{\hfill\qedsymbol}}}

\newcommand{\num}{\alpha}

\newcommand{\NEW}[1]{#1}

\newcommand{\ThreeFiguresExt}[9]{%
   \begin{tabular}{ccc}
     \begin{minipage}{0.31\linewidth}
         \begin{center}
             {\IncludeGraphics[#1]{#2}}
         \end{center}
     \end{minipage}
     &
       \begin{minipage}{0.31\linewidth}
           \centerline{\IncludeGraphics[#4]{#5}}
       \end{minipage}
     &
       \begin{minipage}{0.31\linewidth}
           \centerline{\IncludeGraphics[#7]{#8}}
       \end{minipage}
     \\
     ~\\[-0.2cm]
     {#3} & {#6} & {#9}
   \end{tabular}}

\newcommand{\IncludeGraphics}[2][]{%
   \typeout{}%
   \typeout{Graphics: #2}%
   \typeout{\ includegraphics[#1]{#2}}%
   \includegraphics[#1]{#2}
   \typeout{}%
}

%    -----------------------------------------------------------------
%    -----------------------------------------------------------------
%    -----------------------------------------------------------------
%    Handling references
%    -----------------------------------------------------------------
%    -----------------------------------------------------------------
%    -----------------------------------------------------------------
\newcommand{\HLinkShort}[2]{\hyperref[#2]{#1\ref*{#2}}}
\newcommand{\HLink}[2]{\hyperref[#2]{#1~\ref*{#2}}}
\newcommand{\HLinkPage}[2]{\hyperref[#2]{#1~\ref*{#2}%
      $_\text{p\pageref{#2}}$}}

\newcommand{\HLinkSuffix}[3]{\hyperref[#2]{#1\ref*{#2}{#3}}}
\newcommand{\HLinkPageSuffix}[3]{\hyperref[#2]{#1\ref*{#2}%
      #3$_\text{p\pageref{#2}}$}}

\newcommand{\figlab}[1]{\label{fig:#1}}
\newcommand{\figref}[1]{\HLink{Figure}{fig:#1}}

\newcommand{\exelab}[1]{\label{exer:#1}}
\newcommand{\exeref}[1]{Exercise~\ref{exer:#1}}

\providecommand{\deflab}[1]{\label{def:#1}}
\newcommand{\defref}[1]{\HLink{Definition}{def:#1}}
\newcommand{\defrefY}[2]{\hyperref[def:#2]{#1}}
\newcommand{\defrefpage}[1]{\HLinkPage{Definition}{def:#1}}

\newcommand{\etal}{\textit{et~al.}\xspace}

\newcommand{\obslab}[1]{\label{observation:#1}}
\newcommand{\obsref}[1]{\HLink{Observation}{observation:#1}}
\newcommand{\obsrefY}[2]{\hyperref[observation:#2]{#1}}

\newcommand{\exmlab}[1]{\label{example:#1}}
\newcommand{\exmref}[1]{\HLink{Example}{example:#1}}

\newcommand{\remlab}[1]{\label{rem:#1}}
\newcommand{\remref}[1]{\HLink{Remark}{rem:#1}}%

\newcommand{\chaplab}[1]{\label{chap:#1}}

\newcommand{\itemlab}[1]{\label{item:#1}}
\newcommand{\itemref}[1]{\HLinkSuffix{}{item:#1}{}}

\newcommand{\tedlab}[1]{\label{tedium:#1}}
\newcommand{\tedref}[1]{\HLink{Tedium}{tedium:#1}}%

\newcommand{\factlab}[1]{\label{fact:#1}}%

\newcommand{\alglab}[1]{\label{algorithm:#1}}

\newcommand{\net}{\mathcal{N}}%
      
\newcommand{\problab}[1]{\label{prob:#1}}%
\newcommand{\probref}[1]{\HLink{Problem}{prob:#1}}%
\newcommand{\probrefY}[2]{\hyperref[prob:#1]{#2}}

\newcommand{\lemlab}[1]{\label{lemma:#1}}
\newcommand{\lemref}[1]{\HLink{Lemma}{lemma:#1}}%
\newcommand{\lemrefY}[2]{\hyperref[lemma:#2]{#1}}
\newcommand{\lemrefpage}[1]{\HLinkPage{Lemma}{lemma:#1}}
\newcommand{\lemrefshort}[1]{\HLinkShort{L}{lemma:#1}}

\newcommand{\eqlab}[1]{\label{equation:#1}}
\newcommand{\Eqrefbare}[1]{\HLinkSuffix{(}{equation:#1}{)}}
\newcommand{\Eqref}[1]{\HLinkSuffix{Eq.~(}{equation:#1}{)}}
\newcommand{\Eqrefpage}[1]{\HLinkPageSuffix{Eq.~(}{equation:#1}{)}}

\newcommand{\seclab}[1]{\label{sec:#1}}
\newcommand{\secref}[1]{\HLink{Section}{sec:#1}}
\newcommand{\secrefY}[2]{\hyperref[sec:#2]{#1}}
\newcommand{\secrefpage}[1]{\HLinkPage{Section}{sec:#1}}

\newcommand{\thmlab}[1]{{\label{theo:#1}}}
\newcommand{\thmref}[1]{\HLink{Theorem}{theo:#1}}
\newcommand{\thmrefY}[2]{\hyperref[theo:#2]{#1}}
\newcommand{\thmrefpage}[1]{\HLinkPage{Theorem}{theo:#1}}

\newcommand{\corlab}[1]{\label{cor:#1}}
\newcommand{\corref}[1]{\HLink{Corollary}{cor:#1}}%

%    -----------------------------------------------------------------
%    -----------------------------------------------------------------
%    -----------------------------------------------------------------
%    Handling references
%    -----------------------------------------------------------------
%    -----------------------------------------------------------------
%    -----------------------------------------------------------------

\newcommand{\ceil}[1]{\left\lceil {#1} \right\rceil}
\newcommand{\floor}[1]{\left\lfloor {#1} \right\rfloor}
\newcommand{\brc}[1]{\left\{ {#1} \right\}}
\newcommand{\pth}[2][\!]{\mleft({#2}\mright)}%
\newcommand{\brk}[1]{\mleft[{#1}\mright]}%
\newcommand{\pbrc}[2][\!\!]{#1\left[ {#2} \bigr. \right]}
\newcommand{\pbrcS}[1]{\left[ {#1} \right]}

\newcommand{\dotprod}[2]{{#1}\cdot{#2}}
\newcommand{\DotProd}[2]{\permut{{#1},{#2}}}
\newcommand{\permut}[1]{\left\langle {#1} \right\rangle}

\newcommand{\MakeBig}{}
\newcommand{\MakesBig}{}
\newcommand{\MakeSBig}{}
\newcommand{\MakeVBig}{}

\newcommand{\DistChar}{\mathsf{d}}%
\newcommand{\DistSetY}[2]{\DistChar\pth{#1,#2}}
\newcommand{\distG}[2]{\DistChar_\Graph\pth{#1, #2}}
\newcommand{\DistMtrCharX}[1]{\DistChar_{#1}}

\newcommand{\dY}[2]{\DistChar\pth{#1, #2}}%
\newcommand{\dGY}[2]{\DistChar_\Graph\pth{#1, #2}}%

\newcommand{\cset}{\mathcal{C}}%

\newcommand{\MTR}{\mathcal{M}}
\newcommand{\DistMtr}[2]{\DistChar\pth{#1,#2}}

\newcommand{\dist}[1]{\left\| {#1} \right\|}
\newcommand{\norm}[1]{\left\| {#1} \right\|}

\providecommand{\ds}{\displaystyle}
\newcommand{\sep}[1]{\,\left|\, {#1} \bigr.\right.}%
\newcommand{\sepw}[1]{\left|\, {#1} \right.}%

\newcommand{\ProbExt}[2]{\mathop{\mathbf{Pr}}_{#1}\pbrcx{#2}}
\newcommand{\ProbLTR}{\mathbb{P}}%

\newcommand{\Prob}[1]{\mathop{\ProbLTR} \mleft[ #1 \mright]}%
\newcommand{\ProbCond}[2]{\mathop{\ProbLTR}\!\left[%
       #1 \;\middle\vert\; #2 \right]}
\newcommand{\ProbY}[2]{\ProbLTR^{}_{\!#1}\pbrc{#2}}

\newcommand{\ExChar}{\mathbb{E}}%
\newcommand{\ExSym}{\mathop{\ExChar}}%
\newcommand{\Ex}[2][\!]{\ExSym#1\pbrcx{#2}}
\newcommand{\ExExt}[2]{\ExSym_{#1}\!\pbrcx{#2}}
\newcommand{\ExExtS}[2]{\ExSym_{\substack{#1}}\!\pbrcx{#2}}
\newcommand{\ExCond}[2]{\ExSym\!\left[%
       #1 \;\middle\vert\; #2 \right]}

\newcommand{\VarSymb}{\mathop{\mathbb{V}}}
\newcommand{\Var}[1]{\mathop{\mathbb{V}}\!\pbrcx{#1}}

\newcommand{\pbrcx}[1]{\left[ {#1} \right]}%
\newcommand{\nfrac}[2]{{#1}/{#2}}
\newcommand{\cardin}[1]{\left| {#1} \right|}%

\newcommand\CCancel[2][black]{{\renewcommand\CancelColor{\color{#1}}\cancel{#2}}}

\newcommand{\IncGraphPage}[4][]{%
   \IncludeGraphics[page=#4,#1]{\File{#2/#3}}%
}

\newenvironment{enumEx}{
   \begin{compactenumA}
    }{\end{compactenumA}}

\newcommand{\cl}{{\mathrm{cl}}}

\usepackage[amsmath,thmmarks]{ntheorem}
\theoremseparator{.}%

\newtheorem{theorem:unnumbered}{Theorem}[section]
\newtheorem{theorem}{Theorem}[section]

\newtheorem{observation}[theorem]{Observation}
\newtheorem{conjecture}[theorem]{Conjecture}
\newtheorem{lemma}[theorem]{Lemma}
\newtheorem{corollary}[theorem]{Corollary}
\newtheorem{claim}[theorem]{Claim}
\newtheorem{assumption}[theorem]{Assumption}
\newtheorem{tedium}[theorem]{Tedium}

\theoremstyle{plain}%
\theoremheaderfont{\sffamily}
\theorembodyfont{\upshape}%

\newtheorem{defn}[theorem]{Definition}
\newtheorem{definition}[theorem]{Definition}

\newtheorem*{remark_unnumbered}[theorem:unnumbered]{Remark}%
\newtheorem*{remarks}[theorem]{Remarks}%
\newtheorem{remark}[theorem]{Remark}%
\newtheorem{fact}[theorem]{Fact}%
\newtheorem{example}[theorem]{Example}
\newtheorem{exercise}[theorem]{Exercise}
\newtheorem{problem}[theorem]{Problem}
\newtheorem{algorithm}[theorem]{Algorithm}%
\newtheorem{xca}[theorem]{Exercise}
\newtheorem{openproblem}[theorem]{Open Problem}
\newtheorem{exercise_h}[theorem]{Exercise}

% Proof environments
\theoremheaderfont{\em}%
\theorembodyfont{\upshape}%
\theoremstyle{nonumberplain}
\theoremseparator{}
% TeX4ht's ntheorem support consumes theorem end symbols.  The web-only
% wrappers below place the duck before ntheorem closes its HTML block.
\theoremsymbol{}
\newtheorem{proof}{Proof:}
\newtheorem{proofextin}{Proof}
\newtheorem{altproof}{Alternative proof:}

\theoremsymbol{}
\newtheorem{proofnoqed}{Proof:}

\theoremsymbol{}
\theorempreskip{0pt}%
\theorempostskip{0pt}%
\newtheorem{proofof}{Proof of\!}%

\makeatletter
\let\proof@theorem\proof
\let\endproof@theorem\endproof
\renewenvironment{proof}{\proof@theorem}{\unskip\myqedsymbol\endproof@theorem}
\let\proofextin@theorem\proofextin
\let\endproofextin@theorem\endproofextin
\renewenvironment{proofextin}{\proofextin@theorem}{\unskip\myqedsymbol\endproofextin@theorem}
\let\altproof@theorem\altproof
\let\endaltproof@theorem\endaltproof
\renewenvironment{altproof}{\altproof@theorem}{\unskip\myqedsymbol\endaltproof@theorem}
\let\proofof@theorem\proofof
\let\endproofof@theorem\endproofof
\renewenvironment{proofof}{\proofof@theorem}{\unskip\myqedsymbol\endproofof@theorem}
\makeatother

\newenvironment{proofExt}[1]{{\emph{#1:}}}{\hfill{\hfill\myqedsymbol}}
\newenvironment{proofnoproof}{{}}{{\hfill{\hfill\myqedsymbol}}}

%%%%%%%%%%%%%%%%%%%%%%%%%%%%%%%%%%%%%%%%%%%%%%%%%%%%%%%%%%%%%%%%%%%%
%%%%%%%%%%%%%%%%%%%%%%%%%%%%%%%%%%%%%%%%%%%%%%%%%%%%%%%%%%%%%%%%%%%%
% Macros to embed figures in text

\newlength{\RecordedWidth}
\newcommand{\RecordWidth}[1]{\settowidth{\RecordedWidth}{#1}}%

\newlength{\msize}%

%-------------------------------------------------------------------
% \minipageWFigure{ TEXT }{ Figure }{ Figure caption/label }
%-------------------------------------------------------------------
\newcommand{\minipageWFigure}[3]{%
   \RecordWidth{#2}%
   \setlength{\msize}{\dimexpr(\textwidth-1.1\RecordedWidth)\relax}
   \noindent%
   \SaveIndent%
   \begin{minipage}{\msize}%
       \RestoreIndent%
       #1
   \end{minipage}%
   \null\hfill%
   \begin{minipage}{\RecordedWidth}
       \null\hfill%
       #2%       

       #3%
   \end{minipage}%
}

\newcommand{\EBRY}[2]{#1_{\langle#2\rangle}}%

\newcommand{\NiceFrac}[2]{%
   {\raisebox{0.2em}{\ensuremath{#1}}\!\!\!}%
   {\rotatebox[origin=c]{-20}{\Big/}}% Rotated /
   {\! \raisebox{-0.2em}{\ensuremath{#2}} }}%

% end - Macros to embed figures in text
%%%%%%%%%%%%%%%%%%%%%%%%%%%%%%%%%%%%%%%%%%%%%%%%%%%%%%%%%%%%%%%%%%%% 
%%%%%%%%%%%%%%%%%%%%%%%%%%%%%%%%%%%%%%%%%%%%%%%%%%%%%%%%%%%%%%%%%%%%

%%%%%%%%%%%%%%%%%%%%%%%%%%%%%%%%%%%%%%%%%%%%%%%%%%%%%%%%%%%%%%%%%%
% Color in math
%\providecommand{\Mh}[1]{{#1}}%
\providecommand{\IfColorMath}[2]{#2}%

\IfFileExists{.latex_printer_friendly}{\def\GenPrinterVer{1}}{}%

\ifx\GenPrinterVer\undefined
   \IfFileExists{.latex_color}{\def\GenColorMath{1}}{}
\else
   \renewcommand{\IfPrinterVer}[2]{#1}%
\fi

\ifx\GenColorMath\undefined
\else
\renewcommand{\IfColorMath}[2]{#1}%
%\renewcommand{\Mh}[1]{{\textcolor{red}{#1}}}%
\fi
% Color in math - end
%%%%%%%%%%%%%%%%%%%%%%%%%%%%%%%%%%%%%%%%%%%%%%%%%%%%%%%%%%%%%%%%%%

\newcommand{\HSetY}[2]{H\pth{#1, #2}}%
\newcommand{\CHSetY}[2]{C\pth{#1, #2}}%
\newcommand{\Holder}{H\"older\xspace}
\newcommand{\OldOm}{{\mathcal{O}}}
\newcommand{\NewOm}{{\mathcal{N}}}

%%%%%%%%%%%%%%%%%%%%%%%%%%%%%%%%%%%%%%%%%%%%%%%%%%%%%%%%%%%%%%%%%%%%%%%
%---------------------------------------------------------------
\newcommand{\IntRange}[1]{\left\llbracket #1 \right\rrbracket}
\newcommand{\IRX}[1]{\left\llbracket #1 \right\rrbracket}%
\newcommand{\IRY}[2]{\left\llbracket #1:#2 \right\rrbracket}

%---------------------------------------------------------------
%%%%%%%%%%%%%%%%%%%%%%%%%%%%%%%%%%%%%%%%%%%%%%%%%%%%%%%%%%%%%%%%%%%%%%%

\newcommand{\obj}{f}%
\newcommand{\Wg}{\omega}%
\newcommand{\TotalW}{W}%

\newcommand{\rankY}[2]{{#1}_{ \permut{#2}}}%%
\newcommand{\areaX}[1]{\mathrm{area}\pth{#1}}%%

\newcommand{\mybinom}[2]%
{\Bigl(\begin{array}{@{}c@{}}#1\\#2\end{array}\Bigr)}

\providecommand{\Set}[2]{\left\{ #1 \;\middle\vert\; #2 \right\}}
\newcommand*{\biga}[1]{\vcenter{\hbox{\scalebox{1.5}{\ensuremath#1}}}}
\newcommand{\CNF}{\TermI{CNF}\xspace}%
\newcommand{\Tree}{T}%

\DeclareSymbolFont{sfletters}{OML}{cmbrm}{m}{it}

\DeclareMathSymbol{\salpha}{\mathord}{sfletters}{"0B}

%%%%%%%%%%%%%%%%%%%%%%%%%%%%%%%%%%%%%%%%%%%%%%%%%%%%%%%%%%%%%%%%%%%%

\providecommand{\ProblemInd}[2]{\textsf{\textcolor{red}{#1}\index{Problem!#2}}}
\newcommand{\NPProblemInd}[4]{
   \begin{center}
      \vspace{0.6cm}
       \fbox{
          \begin{minipage}{0.9\linewidth}
              \vspace{-0.6cm} \hspace{-0.35cm}
              {\large \bfseries \ProblemInd{#1}{#2}}
              \vspace{+0.2cm}%

              \noindent{\textcolor{blue}{\bfseries \textsf{Instance}}:} {#3}
              
              \noindent{\textcolor{blue}{\bfseries \textsf{Question}:}} {#4}
          \end{minipage}
       }
   \end{center}
   \smallskip
   }

\newcommand{\NPProblemOpt}[4]{\NPProblemInd{#1}{#2}{#3}{#4}}

%%%%%%%%%%%%%%%%%%%%%%%%%%%%%%%%%%%%%%%%%%%%%%%%%%%%%%%%%%%%%%%%%%%%

\newcommand{\Hastad}{H{\aa}stad\xspace}

\newcommand{\ANNbr}{\Term{ANN}\index{approximate nearest neighbor}\xspace}
\newcommand{\LSH}{\TermI{L{S}H}\xspace}

\newcommand{\DS}{\EuScript{D}} % DataStructure
\newcommand{\DSC}{\mathsf{D}} % DataStructure

\newcommand{\distH}[2]{\DistChar_{H}\pth{#1, #2}}
\newcommand{\distHY}[2]{\DistChar_{H}\pth{#1, #2}}

\newcommand{\DSAprxNearNbr}{\ensuremath{\DataStructure{{\approx%
         \text{Near{}}}}}%
   \index{near neighbor!data-structure!approximate}%
}
\newcommand{\DataStructure}[1]{\ensuremath{\DS_{#1}}}
\newcommand{\DistNN}[2]{\DistChar\pth{#1,#2}}
\newcommand{\pntc}{{p}}%

\newcommand{\seq}{{M}}%
\newcommand{\cTimes}{\beta}%

\newcommand{\DSTimes}{L}%
\newcommand{\rad}{r}%

\newcommand{\HypercubeX}[1]{\Hypercube^{#1}}
\newcommand{\kFunctions}[2]{\mathrm{combine}\pth{#1,#2}}
\newcommand{\WEqualChar}{\Bumpeq}
\newcommand{\WEqual}[2]{{#1}\WEqualChar{#2}}
\newcommand{\FamilyAmp}{\FamilyC_{\WEqualChar}}
\newcommand{\PSPACE}{\ComplexityClass{PSPACE}\xspace}%

\newcommand{\PS}{P}%

\newcommand{\PSetFar}{\ensuremath{\PntSet^{}_{\geq}\xspace}}%
\newcommand{\DA}{D}%
\newcommand{\uniqX}[1]{\mathrm{uniq}\pth{#1}}%
\newcommand{\diffX}[1]{\mathop{}\!\mathrm{d}#1}

\newcommand{\Frq}{F}%
\newcommand{\frq}{f}%
\newcommand{\STRM}{\mathcal{S}}%
\newcommand{\strm}{{s}}%

\providecommand{\Turan}{Tur\'an\xspace}
\newcommand{\fGrowX}[1]{G^{}_{\Dim}\pth{#1}}
\newcommand{\CompressFunc}{\mathtt{Compress}}
\newcommand{\RSA}{\Term{RSA}\xspace}
\newcommand{\MST}{\Term{MST}\xspace}
\providecommand{\CClassPP}{\CClass{P{}P}\xspace}
\providecommand{\ID}{\mathrm{id}}

\newcommand{\EncodeFunc}{\text{\ttfamily{Enc}}}
\newcommand{\DecodeFunc}{\text{\ttfamily{Dec}}}
\newcommand{\discX}[1]{\mathrm{disc}\pth{#1}}%

\newcommand{\propertyC}{\psi}%
\newcommand{\propertyX}[1]{\propertyC\pth{#1}}%
\newcommand{\numC}{\Gamma}%

\newcommand{\degreeX}[1]{\mathrm{degree}\pth{#1}}

\newcommand{\ind}{m}

\providecommand{\nequiv}{\not\equiv}
\providecommand{\sqcupplus}{\sqcup^+}

\newcommand{\AlgorithmI}[1]{{\protect\color{RedViolet}{#1}}}

\usepackage{marginnote}
\usepackage{ragged2e} 
% Short, subtle snark macro
\setlength{\marginparwidth}{1.5in}
\setlength{\marginparsep}{10pt}
\newcommand{\snark}[1]{%
   \marginnote{\raggedright\scriptsize\sffamily\color{gray}#1}%
}

\newcommand{\BigH}{\rule[-0.4em]{0pt}{1.3em}}%

\newcommand{\binsY}[2]{\sqcup_{\geq #1}\pth{#2}}
\newcommand{\ballsY}[2]{\text{\Football}_{\geq #1}\pth{#2}}

\newcommand{\BLX}[1]{{{\sqcupplus{}}_{\geq #1}}}

% Page-breaking hints have no meaning in HTML.  Display bodies are preserved
% verbatim for MathJax, so the builder also removes occurrences found there.
\renewcommand{\allowdisplaybreaks}[1][]{}

\let\Mh\undefined


% -----------------------------------------------------------------------------
% Bibliography & Document Begin
%-----------------------------------------------------------------------------
\IfFileExists{\jobname.bib}{%
   \addbibresource{\jobname.bib}%
}{%
   \addbibresource{shortcuts.bib}%
   \addbibresource{geometry.bib}%
}
\makeindex
\begin{document}
}{\input{prefix_4ht}}%
  \else
    \IfFileExists{styles/prefix_latex.tex}{%-----------------------------------------------------------------------------
% Master Prefix File for Randomized Algorithms Class Notes & Book
% Unified self-contained configuration for LuaLaTeX and XeLaTeX with Memoir
%-----------------------------------------------------------------------------

\RequirePackage{iftex}
\ifPDFTeX
  \PackageError{engine-guard}{This document requires LuaLaTeX or XeLaTeX}{%
    Please recompile using LuaLaTeX or XeLaTeX instead of pdfLaTeX.%
  }
  \stop
\fi

\ifdefined\LWR@started
  \PackageError{latex-gate}{%
    Attempting to run lwarp with prefix_latex.tex.%
  }{Please use prefix_lwarp.tex (or set \string\warpMode) for HTML generation.}
  \stop
\fi

% Always use memoir
\documentclass[12pt,oneside,fleqn]{memoir}

%-----------------------------------------------------------------------------
% Typography, Fonts & Math Packages
%-----------------------------------------------------------------------------
\usepackage{amssymb}
\usepackage{mathtools}
\PassOptionsToPackage{warnings-off={mathtools-overbracket,mathtools-colon}}{unicode-math}
\usepackage{fontspec}
\usepackage{unicode-math}

% New Computer Modern OpenType text and math fonts
\setmainfont{NewCM10-Book.otf}[
  BoldFont          = NewCM10-Bold.otf,
  ItalicFont        = NewCM10-BookItalic.otf,
  BoldItalicFont    = NewCM10-BoldItalic.otf,
  SlantedFont       = NewCM10-Book.otf,
  SlantedFeatures   = {FakeSlant = 0.25},
  BoldSlantedFont   = NewCM10-Bold.otf,
  BoldSlantedFeatures = {FakeSlant = 0.25},
  SmallCapsFeatures = {Letters = SmallCaps}
]
\setsansfont{NewCMSans10-Book.otf}[
  BoldFont          = NewCMSans10-Bold.otf,
  ItalicFont        = NewCMSans10-BookOblique.otf,
  BoldItalicFont    = NewCMSans10-BoldOblique.otf
]
\setmonofont{NewCMMono10-Book.otf}[
  BoldFont          = NewCMMono10-Bold.otf,
  ItalicFont        = NewCMMono10-BookItalic.otf,
  BoldItalicFont    = NewCMMono10-BoldOblique.otf
]

\setmathfont{NewCMMath-Book.otf}
\setmathfont{NewCMMath-Bold.otf}[version=bold]

% OpenType math aliases and fixes
\providecommand{\square}{\mdlgwhtsquare}
\def\Box{\mdlgwhtsquare}
\def\Diamond{\mdlgwhtdiamond}
\def\blacksquare{\mdlgblksquare}
\def\leadsto{\rightsquigarrow}
\providecommand{\llbracket}{\lBrack}
\providecommand{\rrbracket}{\rBrack}
\DeclareTextCommand{\textcircled}{TU}[1]{%
  \mbox{\ooalign{\hfil\raise0.1ex\hbox{#1}\hfil\crcr\char"25EF}}%
}

\usepackage{alphabeta}
\usepackage[normalem]{ulem} % Provides \sout
\usepackage{mleftright}
\usepackage{microtype}
\usepackage{xspace}
\usepackage{xurl}
\usepackage[dvipsnames,table]{xcolor}
\usepackage{graphicx}
\usepackage{caption}
\usepackage{booktabs,array,colortbl,hhline,multirow,tabularx}
\usepackage[inline]{enumitem}
\usepackage[thicklines]{cancel}
\renewcommand\CancelColor{\color{purple}}
\usepackage{euscript}
\usepackage{marvosym}
\let\Cross\undefined
\let\Square\undefined
\usepackage{bbding}
\usepackage{epigraph}
\usepackage{makeidx}
\usepackage{perpage}
\usepackage{tikz}

%-----------------------------------------------------------------------------
% Boondox / mathcalb font declaration
%-----------------------------------------------------------------------------
\DeclareFontFamily{U}{BOONDOX-calo}{\skewchar\font=45 }
\DeclareFontShape{U}{BOONDOX-calo}{m}{n}{
  <-> s*[1.05] BOONDOX-r-calo}{}
\DeclareFontShape{U}{BOONDOX-calo}{b}{n}{
  <-> s*[1.05] BOONDOX-b-calo}{}
\DeclareMathAlphabet{\mathcalb}{U}{BOONDOX-calo}{m}{n}
\SetMathAlphabet{\mathcalb}{bold}{U}{BOONDOX-calo}{b}{n}
\DeclareMathAlphabet{\mathbcalb}{U}{BOONDOX-calo}{b}{n}

%-----------------------------------------------------------------------------
% Algorithm Environment (formerly salgorithm.sty)
%-----------------------------------------------------------------------------
\newcommand{\ProgWordX}[1]{{\small\textbf{#1}}{\ }}
\newcommand{\Do}{\ProgWordX{do}\xspace}
\newcommand{\End}{\ProgWordX{end}\xspace}
\newcommand{\Return}{\ProgWordX{return}\xspace}%
\newcommand{\Procbegin}{\ProgWordX{begin}}
\newcommand{\If}{\ProgWordX{if}\xspace}
\newcommand{\Then}{\ProgWordX{then}\xspace}
\newcommand{\Else}{\ProgWordX{else}\xspace}
\newcommand{\For}{\ProgWordX{for}\xspace}
\newcommand{\While}{{\small\bfseries while}\ }

\newcommand{\xbeginlgox}{\begin{minipage}{1in}\begin{tabbing}
           \quad\=\qquad\=\qquad\=\qquad\=\qquad\=\qquad\=\qquad\=\kill}
\newcommand{\xendlgox}{\end{tabbing}\end{minipage}}

\newenvironment{algorithmEnv}{%
   \begin{tabular}{|l|}\hline\xbeginlgox\bigskip}%
   {\xendlgox\\\hline\end{tabular}}

\newenvironment{program}{%
   \begin{minipage}{4.0in}%
   \begin{tabbing}%
       \ \ \ \ \=\ \ \ \=\ \ \ \ \=\ \ \ \ \=\ \ \ \ \=%
      \ \ \ \ \=\ \ \ \ \=\ \ \ \ \=\ \ \ \ \=%
      \ \ \ \ \=\ \ \ \ \=\ \ \ \ \=\ \ \ \ \=\kill%
}{%
   \end{tabbing}%
   \end{minipage}%
}

\newcommand{\AlgorithmBox}[1]{%
   \fbox{\begin{program} #1 \end{program}}%
}

\definecolor{AlgorithmColor}{cmyk}{0.07,0.90,0,0.34}
\providecommand{\Algorithm}[1]{{\AlgorithmI{#1}\index{algorithm!#1@{\AlgorithmI{#1}}}}}
\newcommand{\CodeComment}[1]{\textcolor{blue}{\texttt{#1}}}

%-----------------------------------------------------------------------------
% Path & Directory Resolution
%-----------------------------------------------------------------------------
\newcommand{\RootDIR}{../}
\IfFileExists{../../../book.tex}{\renewcommand{\RootDIR}{../../../}}{}
\IfFileExists{../../book.tex}{\renewcommand{\RootDIR}{../../}}{}
\IfFileExists{../book.tex}{\renewcommand{\RootDIR}{../}}{}
\IfFileExists{./book.tex}{\renewcommand{\RootDIR}{}}{}

\providecommand{\File}[1]{#1}
\providecommand{\RootFile}[1]{\RootDIR #1}

%-----------------------------------------------------------------------------
% QED Duck Symbol
%-----------------------------------------------------------------------------
\newlength{\CharHeight}
\AtBeginDocument{\setlength{\CharHeight}{\fontcharht\font`X}}

\IfFileExists{styles/qed_duck.pdf}{%
   \providecommand{\myqedsymbol}{%
      \reflectbox{\includegraphics[height=1.4\CharHeight]{styles/qed_duck.pdf}}%
   }%
}{%
   \providecommand{\myqedsymbol}{$\blacksquare$}%
}

%-----------------------------------------------------------------------------
% BibLaTeX Setup (Pure Biber, No BibTeX)
%-----------------------------------------------------------------------------
\usepackage[bibencoding=utf8,style=alphabetic,refsegment=chapter,
indexing=cite,
backend=biber]{biblatex}

\ExecuteBibliographyOptions{%
   maxnames = 40,%
   maxalphanames = 4,%
   maxcitenames = 3,%
   minalphanames = 3,%
   minnames = 1,%
   giveninits = true,%
   doi = false,%
   url = false,%
   isbn = false%
}

\renewbibmacro{in:}{}

\AtEveryBibitem{%
   \ifentrytype{inproceedings}{%
      \clearname{editor}%
      \clearlist{editor}%
      \clearfield{editor}%
      \clearlist{publisher}%
      \clearfield{series}%
      \clearlist{location}%
      \clearfield{publisher}%
   }{}%
   \ifentrytype{book}{%
      \clearname{publisher}%
   }{}%
}
\AtEveryCitekey{%
   \ifentrytype{inproceedings}{%
      \clearlist{publisher}%
   }{}%
   \clearname{editor}%
   \clearlist{location}%
}
\DeclareFieldFormat[inproceedings]{date}{}
\DeclareFieldFormat[inproceedings]{pages}{#1, \printfield{year}}
\DeclareFieldFormat[inproceedings]{booktitle}{\mkbibemph{#1,\nopunct}}

\DeclareFieldFormat{journaltitle}{\mkbibemph{#1}\isdot}

\DeclareFieldFormat[article]{number}{\mkbibparens{#1}}
\DeclareFieldFormat[article]{date}{}
\DeclareFieldFormat[article]{pages}{#1, \printfield{year}}
\DeclareFieldFormat[article]{issue}{\nopunct#1}
\renewbibmacro*{issue+date}{%
   \iffieldundef{issue}{}%
   {%
      \printtext{\unspace}%
      (\printfield{issue})%
   }%
   \iffieldundef{pages}%
   {\addcomma\printfield{year}}%
   {}%
}

\renewcommand{\bibpagespunct}{%
   \ifentrytype{article}{\addcolon}{\addcomma}\space%
}

\DeclareFieldFormat[article]{titlecase}{\MakeSentenceCase{#1}}
\DeclareFieldFormat[inproceedings]{titlecase}{\MakeSentenceCase{#1}}
\DeclareFieldFormat[incollection]{titlecase}{\MakeSentenceCase{#1}}
\DeclareFieldFormat[book]{titlecase}{\MakeSentenceCase{#1}}

\renewbibmacro*{journal}{%
   \iffieldundef{journaltitle} {} {\printtext[journaltitle]{%
         \printfield[myplain]{journaltitle}%
         \setunit{\subtitlepunct}%
         \printfield[myplain]{journalsubtitle}}}}
\DeclareFieldFormat{myplain}{#1}

\renewbibmacro*{booktitle}{%
   \ifboolexpr{ test {\iffieldundef{booktitle}} and test
      {\iffieldundef{booksubtitle}} } {} {\printtext[booktitle]{%
         \printfield[myplain]{booktitle}%
         \setunit{\subtitlepunct}%
         \printfield[myplain]{booksubtitle}}%
      \newunit}%
   \printfield{booktitleaddon}}

\renewbibmacro*{volume+number+eid}{%
   \setunit*{\addcomma\addspace}%
   \printfield{volume}%
   \printfield{number}%
   \setunit{\addcomma\space}%
   \printfield{eid}%
}

\newbibmacro{string+doi}[1]{%
   \iffieldundef{doi}{%
      \iffieldundef{url}{%
         #1%
      }{%
         \href{\thefield{url}}{#1}%
      }%
   }{%
      \href{http://dx.doi.org/\thefield{doi}}{#1}}%
}
\DeclareFieldFormat*{title}{\usebibmacro{string+doi}{\mkbibemph{#1}}}
\renewcommand*{\multicitedelim}{\addcomma\space}
\DeclareFieldFormat[inproceedings]{title}{\usebibmacro{string+doi}{{#1}}}
\DeclareFieldFormat[article]{title}{\usebibmacro{string+doi}{{#1}}}
\DeclareFieldFormat[incollection]{title}{\usebibmacro{string+doi}{{#1}}}
\DeclareFieldFormat[book]{titlecase}{{#1}}

\DeclareCiteCommand{\tcite}%
{\usebibmacro{prenote}}%
{\usebibmacro{citeindex}\usebibmacro{cite}}%
{\multicitedelim}%
{\usebibmacro{postnote}}

%-----------------------------------------------------------------------------
% Hyperref & Visual Formatting
%-----------------------------------------------------------------------------
\usepackage{hyperref}
\hypersetup{%
   breaklinks,%
   colorlinks=true,%
   linkcolor=[rgb]{0.45,0.0,0.0},%
   citecolor=[rgb]{0,0,0.45}%
}

\definecolor{DeepNavy}{HTML}{002366}
\definecolor{RoyalBlue}{HTML}{004E92}
\definecolor{SteelBlue}{HTML}{4682B4}
\definecolor{TealBlue}{HTML}{008080}
\definecolor{DarkSlate}{HTML}{2F4F4F}
\definecolor{DustyGrey}{HTML}{696969}

\makeatletter
\AtBeginDocument{\let\@elt\relax}
\makeatother

\headstyles{bringhurst}
\chapterstyle{dash}

\renewcommand{\chapnamefont}{\normalfont\huge\bfseries\color{DeepNavy}}
\renewcommand{\chapnumfont}{\normalfont\huge\bfseries\color{DeepNavy}}
\renewcommand{\chaptitlefont}{\normalfont\Huge\bfseries\color{DeepNavy}}

\counterwithin{footnote}{chapter}

\setsubparaheadstyle{\normalfont\normalsize\itshape\color{DustyGrey}}
\setlength{\beforechapskip}{-20pt} 
\setlength{\afterchapskip}{30pt}
\setsecnumformat{\csname the#1\endcsname.\quad}

%-----------------------------------------------------------------------------
% Standalone Chapter vs Master Book Modes
%-----------------------------------------------------------------------------
\newcommand{\Copyright}{This work is licensed under the Creative Commons Attribution-Noncommercial 3.0 License. To view a copy of this license, visit \url{http://creativecommons.org/licenses/by-nc/3.0/} or send a letter to Creative Commons, 171 Second Street, Suite 300, San Francisco, California, 94105, USA.}

\newcommand{\RealChapterEndInner}{%
   {%
      \begingroup%
      \emergencystretch=1em%
      \printbibliography[heading=subbibliography,segment=\therefsegment]{}%
      \endgroup%
   }%
}

\providecommand{\BookMode}[1]{}
\providecommand{\StandAloneMode}[1]{#1}

\ifx\bookMode\undefined
   % STANDALONE CHAPTER MODE
   \renewcommand{\StandAloneMode}[1]{#1}
   \renewcommand{\BookMode}[1]{}

   \newcommand{\ChapterNumPage}[2]{%
      \if\relax\detokenize{#1}\relax
         \setcounter{chapter}{0}%
      \else\ifnum#1>0
         \setcounter{chapter}{\numexpr#1-1\relax}%
      \else
         \setcounter{chapter}{0}%
      \fi\fi
      \if\relax\detokenize{#2}\relax
         \setcounter{page}{1}%
      \else\ifnum#2>0
         \setcounter{page}{#2}%
      \else
         \setcounter{page}{1}%
      \fi\fi
   }

   \DeclareRobustCommand{\setheadingcolor}{\color{RoyalBlue}}
   \setsecheadstyle{\setheadingcolor\normalfont\Large\bfseries}
   \setsubsecheadstyle{\normalfont\large\bfseries\color{SteelBlue}}
   \setsubsubsecheadstyle{\normalfont\normalsize\bfseries\color{TealBlue}}
   \setparaheadstyle{\normalfont\normalsize\bfseries\color{DarkSlate}}

   \DeclareDocumentCommand{\Chapter}{s o m}{%
      \IfBooleanTF{#1}{%
         \IfValueTF{#2}{\chapter*[]{#2}}{\chapter*[]{#3}}%
      }{%
         \IfValueTF{#2}{\chapter*[]{#2}}{\chapter[]{#3}}%
      }%
      \setcounter{footnote}{0}%
   }
   \newcommand{\ChapterPrefix}{}
   \newcommand{\ChapterGen}{\pagestyle{plain}\vspace*{-0.7cm}}
   \newcommand{\ChapterEnd}{%
      \RealChapterEndInner{}%
      \end{document}%
   }

   \renewcommand{\File}[1]{#1}
   \renewcommand{\RootFile}[1]{../#1}
\else
   % MASTER BOOK MODE
   \renewcommand{\StandAloneMode}[1]{}
   \renewcommand{\BookMode}[1]{#1}

   \newcommand{\ChapterNumPage}[2]{}

   \DeclareDocumentCommand{\Chapter}{s o m}{%
      \IfBooleanTF{#1}{%
         \IfValueTF{#2}{\chapter*[#2]{#3}}{\chapter*{#3}}%
         \ifx\Dir\undefined\else
            \typeout{CHAPTER_MAP: \Dir/\CurrentChapFile => chap: 0, page: \thepage}%
         \fi
      }{%
         \IfValueTF{#2}{\chapter[#2]{#3}}{\chapter{#3}}%
         \ifx\Dir\undefined\else
            \typeout{CHAPTER_MAP: \Dir/\CurrentChapFile => chap: \thechapter, page: \thepage}%
         \fi
      }%
   }
   \newcommand{\ChapterPrefix}{%
      \vspace{-1cm}%
      \subsection*{\small
         {CS 574: Randomized Algorithms}\\
         {Sariel Har-Peled}\\
         {\today}}%
   } \newcommand{\ChapterGen}{}
   \newcommand{\ChapterEnd}{\RealChapterEndInner{}}

   \renewcommand{\File}[1]{\Dir/#1}
   \renewcommand{\RootFile}[1]{#1}

   \newcommand{\InputChap}[2]{%
      \def\Dir{#1}%
      \def\CurrentChapFile{#2}%
      \typeout{}%
      \typeout{------------------------------------------------------------}%
      \typeout{PROCESSING FILE: #1/#2}%
      \typeout{}%
      \input{#1/#2}%
   }
\fi

%-----------------------------------------------------------------------------
% Fragments & Quotes
%-----------------------------------------------------------------------------
\newcommand{\InputExt}[1]{%
   \IfFileExists{#1}%
   {\input{#1}}%
   {\IfFileExists{../#1}%
      {\input{../#1}}%
      {\IfFileExists{../../#1}%
         {\input{../../#1}}%
         {\typeout{ERROR: Unable to find #1}}%
      }%
   }%
}
\makeatletter
\def\FragNormKey#1{%
  \frag@norm@check#1def_\relax{#1}%
}
\def\frag@norm@check#1def_#2\relax#3{%
  \if\relax\detokenize{#1}\relax
    #3%
  \else
    def_#3%
  \fi
}
\makeatother

\newcommand{\FragmentName}[1]{\RootFile{fragment/\FragNormKey{#1}.tex}}

\newenvironment{fragment}[1]{}{}

\makeatletter
\AtBeginDocument{%
   \let\real@latex@label\label
   \def\label#1{%
      \global\@namedef{all@labels@seen@#1}{1}%
      \real@latex@label{#1}%
   }%
}

\DeclareDocumentCommand{\IncFragment}{m}{%
   \begingroup
      \def\label##1{%
         \@ifundefined{all@labels@seen@##1}{%
            \global\@namedef{all@labels@seen@##1}{1}%
            \real@latex@label{##1}%
         }{}%
      }%
      \IfFileExists{\FragmentName{#1}}%
         {\InputExt{\FragmentName{#1}}}%
         {\InputExt{\RootFile{fragment/#1.tex}}}%
   \endgroup
}
\makeatother

\newcommand{\Quote}[2]{%
   \SaveIndent
   \begin{center}%
      {\footnotesize
         \begin{quote}
             \RestoreIndent
             {#1}\\
             \hspace*{2cm}-- {#2.}
         \end{quote}%
      }%
      \index{quotation!#2}%
   \end{center}%
   \bigskip
}

\newcommand{\QuotePExt}[3]{%
   \bigskip%
   \hfill%
   \begin{minipage}{0.95\linewidth}%
       \footnotesize {#1}%

       \rule{\linewidth}{0.3mm}%
       \hfill #2, #3%
   \end{minipage}%
   \bigskip%
}

\epigraphfontsize{\small\itshape}
\setlength{\epigraphwidth}{0.6\textwidth}
\newcommand{\QuoteExt}[3]{%
   \vspace*{-1.6cm}%
   \epigraph{{\footnotesize #1}}{{#2}, {#3}}%
   \vspace*{-1.2cm}%
}

\providecommand{\delX}[1]{{\color{red}\sout{#1}}}
\providecommand{\newX}[1]{{\color{blue}#1}}
\providecommand{\chgY}[2]{{\color{red}\sout{#1}}{\color{blue}#2}}
\providecommand{\remX}[1]{{\color{purple}\textbf{[Remark: #1]}}}
\newcommand{\SpellIgnore}[1]{#1}
\newcommand{\EQ}{\mathrel{\;\;=\;\;}}
\newcommand{\LEQ}{\mathrel{\;\;\leq\;\;}}
\newcommand{\LT}{\mathrel{\;\;<\;\;}}
\newcommand{\GT}{\mathrel{\;\;<\;\;}}

\newcommand{\GEQ}{\mathrel{\;\;\geq\;\;}}

%-----------------------------------------------------------------------------
% Save a symbol that we know is going to get redefined.
\def\savesymbol#1{%
  \expandafter\let\expandafter\origsym\expandafter=\csname#1\endcsname
  \expandafter\let\csname orig#1\endcsname=\origsym
  \expandafter\let\csname#1\endcsname=\relax
}

% Restore a previously saved symbol, and rename the current one.
\def\restoresymbol#1#2{%
  \expandafter\let\expandafter\newsym\expandafter=\csname#2\endcsname
  \expandafter\global\expandafter\let\csname#1#2\endcsname=\newsym
  \expandafter\let\expandafter\origsym\expandafter=\csname orig#2\endcsname
  \expandafter\global\expandafter\let\csname#2\endcsname=\origsym
}

\usepackage{pictex}%
\savesymbol{triangle}%
\usepackage{tree}%
\restoresymbol{TreePKG}{triangle}%

% Definitions absorbed from sbook.sty
\providecommand{\MdPrc}{\nu}

\providecommand{\TwoFiguresBase}[4]{%
   \begin{tabular}{cc}
      \begin{minipage}{0.47\linewidth}
          \begin{center}
              {{{#1}}}
          \end{center}
      \end{minipage}
      &
      \begin{minipage}{0.47\linewidth}
          \begin{center}
              {{#3}}
          \end{center}
      \end{minipage}
      \\
      {#2} & {#4}
   \end{tabular}%
}

\providecommand{\TwoFigures}[6]{%
   \begin{tabular}{cc}
      \begin{minipage}{0.48\linewidth}
          \begin{center}
              {\IncludeGraphics[#1]{#2}}
          \end{center}
      \end{minipage}
      &
      \begin{minipage}{0.48\linewidth}
          \centerline{\IncludeGraphics[#4]{#5}}
      \end{minipage}
      \\
      {#3} & {#6}
   \end{tabular}%
}

\providecommand{\ThreeFigures}[6]{%
   \begin{tabular}{ccc}
        \begin{minipage}{0.30\linewidth}
            \IncludeGraphics[width=5.1cm]{#1}\\
        \end{minipage}
        &
        \begin{minipage}{0.30\linewidth}
            \IncludeGraphics[width=5.1cm]{#2}\\
        \end{minipage}
        &
        \begin{minipage}{0.30\linewidth}
            \IncludeGraphics[width=5.1cm]{#3}\\
        \end{minipage}
        \vspace{-1cm}
        \\
        {#4} & {#5}& {#6}
   \end{tabular}%
}

%-----------------------------------------------------------------------------
\providecommand{\remove}[1]{}
\providecommand{\NotAMSMode}[1]{#1}
\providecommand{\Solution}[1]{}
\providecommand{\Version}[1]{\hfill{Version: #1}}
\providecommand{\Distor}{\mathrm{dist}}
\providecommand{\Hardness}[1]{(h:#1)}
\providecommand{\QuoteOptExt}[3]{\QuoteExt{#1}{#2}{#3}}
\providecommand{\Hypercube}{\EuScript{H}}
\providecommand{\CoverTime}{\EuScript{C}}
\providecommand{\Nbr}{\Gamma}
\providecommand{\CommuteTime}[2]{\mathbf{CT}_{#1#2}}
\providecommand{\Resistance}{\mathbf{R}}
\providecommand{\EffectResist}[2]{\mathbf{R}_{#1#2}}
\providecommand{\NetworkE}{\EuScript{N}}
\providecommand{\trinom}[4]{\binom{#1}{#2 \;\; #3 \;\; #4}}
\providecommand{\Ranges}{\EuScript{R}}
\providecommand{\CClass}[1]{\ensuremath{\mathbf{#1}}}
\providecommand{\coCClass}[1]{\ensuremath{\mathrm{co-}}\CClass{#1}}
\providecommand{\CX}{\EuScript{C}}
\providecommand{\packet}{\mathbf{c}}
\providecommand{\packetA}{\mathbf{d}}
\providecommand{\mopt}{m_{\mathrm{opt}}}
\providecommand{\girth}{\mathop{\mathrm{girth}}}
\providecommand{\CR}{\mathop{\mathrm{cr}}}
\providecommand{\TRUE}{{\textcolor[named]{RedViolet}{\textsf{TRUE}}}\xspace}
\providecommand{\FALSE}{{\textcolor[named]{RedViolet}{\textsf{FALSE}}}\xspace}
\providecommand{\Part}[1]{\part{#1}}
\providecommand{\Section}[1]{\section{#1}}
\providecommand{\SubSection}[1]{\subsection{#1}}
\providecommand{\SubSubSection}[1]{\subsubsection{#1}}
\providecommand{\Exercise}[5]{%
\begin{exercise_h}[#1]%
   \exelab{#3}{}%
   #4%
\end{exercise_h}%
}
\providecommand{\si}[1]{#1}
\providecommand{\myParagraph}[1]{\noindent\textbf{#1}}
\providecommand{\DefName}[1]{#1}
\providecommand{\deflabel}[1]{#1}
\providecommand{\Description}[1]{}

%-----------------------------------------------------------------------------
% Mathematical Notations & Environments (from notations.sty)
%-----------------------------------------------------------------------------
%-----------------------------------------------------------------------------
% List environments (questions & compactenum)
%-----------------------------------------------------------------------------

\newlist{compactenumA}{enumerate}{5}%
\setlist[compactenumA]{topsep=0pt,itemsep=-1ex,partopsep=1ex,parsep=1ex,%
   label=(\Alph*)}%

\newlist{compactenumI}{enumerate}{5}%
\setlist[compactenumI]{topsep=0pt,itemsep=-1ex,partopsep=1ex,parsep=1ex,%
   label=(\Roman*)}%

\newlist{compactenuma}{enumerate}{5}%
\setlist[compactenuma]{topsep=0pt,itemsep=-1ex,partopsep=1ex,parsep=1ex,%
   label=(\alph*)}%

\newlist{compactitem}{enumerate}{5}%
\setlist[compactitem]{topsep=0pt,itemsep=-1ex,partopsep=1ex,parsep=1ex,%
   label=\textbullet}%

\newlist{compactenumi}{enumerate}{5}%
\setlist[compactenumi]{topsep=0pt,itemsep=-1ex,partopsep=1ex,parsep=1ex,%
   label=(\roman*)}%

\setcounter{secnumdepth}{5}

\newcommand{\HST}{\TermI{HST}\xspace}
\newcommand{\Term}[1]{\textsf{#1}}
\newcommand{\TermI}[1]{\Term{#1}\index{#1@\Term{#1}}}

\DeclareMathAlphabet{\mathpzc}{OT1}{pzc}{m}{it}

\newcommand{\seg}{\mathsf{s}}

\newcommand{\Diam}{\mathrm{diam}}
\newcommand{\diam}{\mathrm{diam}}

\newcommand{\Copt}{C_{\mathrm{opt}}}
\newcommand{\CenterPnt}{\mathrm{cen}}

\newcommand{\parent}{\overline{\mathrm{p}}}
\newcommand{\parentX}[2][\!]{\parent\pth[#1]{#2}}

\newcommand{\FMS}{\EuScript{X}\index{metric space}}  
\newcommand{\FMSB}{\EuScript{Y}\index{metric space}}  
\newcommand{\FastCut}{\Algorithm{FastCut}\xspace}
\newcommand{\MinCut}{\Algorithm{MinCut}\xspace}
\newcommand{\MinCutRep}{\Algorithm{MinCutRep}\xspace}
\newcommand{\Contract}{\Algorithm{Contract}\xspace}

\newcommand{\Grid}{\mathrm{G}}
\newcommand{\ClosestPair}{\mathcal{CP}}
\newcommand{\SphereChar}{\ensuremath{\mathbb{S}}}
\newcommand{\Sphere}[1]{{\SphereChar^{(#1)}}}

\newcommand{\PntSet}{{\mathsf{P}}}%
\providecommand{\PntSetA}{{\mathsf{Q}}}%

\newcommand{\query}{\mathcalb{q}}

\newcommand{\pnt}{\mathsf{p}}
\newcommand{\pntA}{\mathsf{u}}
\newcommand{\pntB}{\mathsf{v}}
\newcommand{\pntC}{{\mathsf{x}}}
\newcommand{\pntD}{{\mathsf{y}}}
\newcommand{\pntE}{{\mathsf{z}}}

\newcommand{\pntCoord}{p}
\newcommand{\pntCoordA}{q}

%%%%%%%%%%%%%%%%%%%%%%%%%%%%%%%%%%%%%%%%%%%%%%%%%%%%%%%%%%%%%%%%%%%%% 
%%%%%%%%%%%%%%%%%%%%%%%%%%%%%%%%%%%%%%%%%%%%%%%%%%%%%%%%%%%%%%%%%%%%%
% Markov chain notations.

\newcommand{\MCSpc}{\mathtt{S}}
\newcommand{\MCPrb}{\mathtt{P}}
\newcommand{\MCPrbY}[2]{\MCPrb_{#1#2}}
\newcommand{\MCPrbM}{\mathbf{P}}%
\newcommand{\MCPrbT}[3]{\MCPrb_{#1#2}^{(#3)}}
\newcommand{\MCFPrbT}[3]{\mathtt{r}_{#1#2}^{(#3)}}
\newcommand{\MCVstProb}[2]{\mathtt{f}_{#1#2}}
\newcommand{\MCExSteps}[2]{\mathtt{h}_{#1#2}}
\newcommand{\MCPVec}[1]{\mathbf{q}^{(#1)}}
%%%%%%%%%%%%%%%%%%%%%%%%%%%%%%%%%%%%%%%%%%%%%%%%%%%%%%%%%%%%%%%%%%%%%

\newcommand{\Spread}{\Phi}
\newcommand{\VC}{\ensuremath{\mathsf{VC}}\xspace}
\newcommand{\MC}{\ensuremath{\mathsf{MC}}\xspace}
\newcommand{\hyperplane}{h}

\newcommand{\QuickSort}{\Algorithm{QuickSort}\xspace}
\newcommand{\Alg}{\Algorithm{Alg}\xspace}

\newcommand{\Graph}{\mathsf{G}}%
\newcommand{\GraphB}{\mathsf{H}}
\newcommand{\GraphC}{\mathsf{K}}
\newcommand{\GraphD}{\mathsf{F}}

\providecommand{\G}{}
\renewcommand{\G}{\Graph}

\newcommand{\RLP}{\ComplexityClass{RLP}\xspace}
\newcommand{\USTCON}{\ComplexityClass{USTCON}}
\newcommand{\STCON}{\overrightarrow{\ComplexityClass{STCON}}}

\newcommand{\Mtrx}{\mathsf{M}}
\newcommand{\TMtrx}{\mathsf{Q}}
\newcommand{\IMtrx}{\mathcal{I}}

\newcommand{\Expander}[3]{$\left[#1,#2,#3\right]$-expander}

\newcommand{\ChartX}[1]{\ensuremath{\chi_{#1}}}
\newcommand{\Field}{\mathbb{F}}
\newcommand{\LSpace}{\mathbb{G}}
\newcommand{\LD}[2]{\mathsf{L{D}}\pth{#1,#2}}
\newcommand{\SEVN}{\widehat{\lambda_2}}
\newcommand{\NX}[2]{#1\!\left[#2\right]}
\newcommand{\Entry}[2]{#1\!\left[#2\right]}
\newcommand{\NXExt}[3]{#1_{#2}\!\left[#3\right]}
\newcommand{\RangeI}[1]{\llbracket #1 \rrbracket}
\newcommand{\DialX}[1]{\gamma\pth{#1}}
\newcommand{\BiDialX}[1]{{\gamma}_2\pth{#1}}

\newcommand{\EV}[1]{\widehat{\lambda_{#1}}}
\newcommand{\TEV}[1]{\widehat{\lambda_{#1}}}
\newcommand{\EVal}{\widehat{\lambda}}
\newcommand{\Basis}{\mathcal{B}}

\newcommand{\ZProd}[2]{\ensuremath{#1 \, \text{\textcircled{z}} #2}}
\newcommand{\RProd}[2]{\ensuremath{#1 \, \text{\textcircled{r}} #2}}

\newcommand{\OneVec}[1]{\mathbf{1}^{#1}}
\newcommand{\DstrbV}[1]{\mathbf{\mathsf{q}}^{(#1)}}
\newcommand{\DstrbE}[2]{\mathsf{q}^{(#1)}_{#2}}

\newcommand{\DEdges}{\widetilde{\Edges}}

\newcommand{\VV}{\mathsf{V}}%
\newcommand{\EE}{\mathsf{E}}%

\newcommand{\Vertices}{\VV}%
\newcommand{\Edges}{\EE}%

\newcommand{\EdgesX}[1]{\Edges\pth{#1}}
\newcommand{\VerticesX}[1]{V\pth{#1}}

\newcommand{\StringX}[1]{\mathsf{s}\pth{#1}}

\newcommand{\BMtrx}{\mathsf{B}}
\newcommand{\WMtrx}{\mathsf{W}}
\newcommand{\NWMtrx}{\overline{\mathsf{W}}}
\newcommand{\Sequence}{\EuScript{S}}

\newcommand{\med}{\mathrm{med}}
\newcommand{\medX}[1]{\mathrm{med}\pth{#1}}

\newcommand{\medianC}{\nu}%

\newlength{\savedparindent}
\newcommand{\SaveIndent}{\setlength{\savedparindent}{\parindent}}
\newcommand{\RestoreIndent}{\setlength{\parindent}{\savedparindent}}

\newcommand{\NNN}{\mathbb{N}}%
\providecommand{\NN}{\mathbb{N}}%
\AtBeginDocument{%
   \renewcommand{\Re}{\mathbb{R}}%
}
%\shogi
\newcommand{\reals}{\Re}%
\newcommand{\ZZ}{\mathbb{Z}}%

\newcommand{\level}{\mathop{\mathrm{level}}}

\providecommand{\TPDF}[2]{\texorpdfstring{#1}{#2}}

\newcommand{\Hint}[1]{\emph{[#1]}}
\newcommand{\Arr}{\mathop{\mathrm{\EuScript{A}}}}

\newcommand{\Array}{\mathcal{T}}%

\newcommand{\dirEdge}[2]{\left({{#1} \rightarrow {#2}} \right)}

%%%%%%%%%%%%%%%%%%%%%%%%%%%%%%%%%%%%%%%%%%%%%%%%%%%%%%%%%%%%%%%%%%%%%%
% Emphasize stuff...
\definecolor{blue25}{rgb}{0, 0, 11}
\newcommand{\emphic}[2]{%
%   \BADVERYBAD%
   \textcolor{blue25}{%
      \bfseries{\emph{#1}}}%
   \index{#2}}
\newcommand{\emphi}[1]{\emphic{#1}{#1}}
\newcommand{\emphOnly}[1]{\textcolor{blue25}{\bfseries{\emph{#1}}}}
\providecommand{\emphiC}[1]{\emphic{#1}{#1}}


\definecolor{almostblack}{rgb}{0, 0, 0.3}
\providecommand{\emphw}[1]{{\textcolor{almostblack}{\emph{#1}}}}%

%
%%%%%%%%%%%%%%%%%%%%%%%%%%%%%%%%%%%%%%%%%%%%%%%%%%%%%%%%%%%%%%%%%%%%%%

\newcommand{\volX}[1]{\mathrm{vol}\pth{#1}}

\newcommand{\Interval}{\EuScript{I}}
\newcommand{\IntervalA}{\EuScript{I}}

\newcommand{\Exp}{\mathrm{Exp}}

\let\Ball\undefined
\newcommand{\Ball}{\mathbf{b}}
\newcommand{\hdim}{d}
\newcommand{\Domain}{\mathcal{D}}
\newcommand{\points}[1]{[\bfseries{#1 points}]}
\newcommand{\Event}{\EuScript{E}}%
\newcommand{\EventB}{\EuScript{B}}%
\newcommand{\EventF}{\EuScript{F}}%
\newcommand{\EventG}{\EuScript{G}}%

\newcommand{\distCmd}[1]{\left\| {#1}  \right\|}
\newcommand{\distY}[2]{\distCmd{#1 - #2}}

\newcommand{\NormalC}{\mathbf{N}}
\newcommand{\NormalY}[2]{\NormalC\pth{#1, #2}}%
\newcommand{\NormalDistOneDim}{\NormalC}
\newcommand{\NormalDistX}[1]{\NormalC^{#1}}
\DeclareMathOperator{\lca}{lca}
\newcommand{\RangeSpace}{\mathsf{S}}%
\newcommand{\RangeSpaceA}{\mathsf{T}}

\newcommand{\range}{\mathbf{r}}
\newcommand{\rangeA}{\mathbf{r}'}

\newcommand{\BadProb}{\varphi}

\newcommand{\EntropyChar}{\mathbb{H}}
\newcommand{\Entropy}[1]{\EntropyChar\pth{#1}}
\newcommand{\EntropyX}[1]{\EntropyChar\pth[]{#1}}

\newcommand{\Matrix}{\mathsf{M}}
\newcommand{\MatrixB}{\mathsf{B}}
\newcommand{\MatrixC}{\mathsf{C}}
\newcommand{\MatrixD}{\mathsf{D}}

\renewcommand{\th}{th\xspace}

\newcommand{\Family}{\mathcal{F}}
\newcommand{\FamilyB}{\mathcal{G}}
\newcommand{\FamilyC}{\mathcal{H}}
\newcommand{\FMSC}{\EuScript{Y}\index{metric space}}

\newcommand{\median}{\overline{\mathsf{m}}}
\newcommand{\QuickSelect}{\Algorithm{Quick{}Select}\xspace}

\providecommand{\row}{\mathbf{r}}%
\newcommand{\rowc}{{r}}%

\newcommand{\vecb}{\mathbf{b}}

\newcommand{\vecA}{\mathbf{v}}
\newcommand{\vecB}{\mathbf{u}}
\newcommand{\vecRand}{\mathbf{r}}

\newcommand{\vecH}{\mathbf{h}}
\newcommand{\vecV}{\mathbf{v}}

\newcommand{\vecX}{\mathbf{x}}
\newcommand{\vecY}{\mathbf{y}}

\newcommand{\DistBinom}[2]{\mathrm{Bin}\pth{#1,#2}}

\newcommand{\DistGeom}[1]{\mathrm{Geom}\pth{#1}}

\newcommand{\Sample}{\mathsf{R}}

\newcommand{\lcm}{\mathop{\mathrm{lcm}}}
\newcommand{\bdiv}{\mathop{\mathrm{div}}}%
\newcommand{\divides}{\mid}%
\newcommand{\poly}{\mathrm{poly}}%

\newcommand{\Group}{\mathcal{G}} 
\newcommand{\GroupA}{\mathcal{H}} 

\newcommand{\algGCD}{\Algorithm{EuclidGCD}\xspace}%

\newcommand{\gOp}{\times}%
\newcommand{\gIdent}{\mathsf{i}}%

\newcommand{\ordX}[2][\!]{\mathop{\mathrm{o{r}d}}\pth[#1]{#2}}

\newcommand{\lSym}[2]{%
   \left(
       \ds {#1}\divides{#2}%
   \right) %
}
\newcommand{\jSym}[2]{\left\llbracket%
       \ds {#1}\divides{#2}%
       \right\rrbracket%
}
\newcommand{\Line}{\ell}

\newcommand{\LevelChar}{\mathbb{L}}
\newcommand{\AtMostKLevelChar}{\LevelChar_{\leq k}}
\newcommand{\AtMostKLevel}[1]{\LevelChar_{\leq k} \pth{ #1}}

\newcommand{\LineSet}{\mathsf{L}}%
\newcommand{\ArrX}[1]{\EuScript{A}\pth{#1}}%

\providecommand{\Matousek}{Matou{\v s}ek\xspace}
\providecommand{\Erdos}{Erd{\H o}s\xspace}
\providecommand{\Lovasz}{Lov\'asz\xspace}

\newcommand{\ProblemC}[1]{{\textsf{\textcolor{red}{#1}\index{problem!#1}}}}
\DefineNamedColor{named}{OliveGreen} {cmyk}{0.64,0,0.95,0.40}
\providecommand{\ComplexityClass}[1]{{{\textcolor[named]{OliveGreen}{%
            \textsc{#1}}}}}
\providecommand{\NPHard}{{\ComplexityClass{NP-Hard}}%
   \index{NP!hard}\xspace}
\providecommand{\NPComplete}{\ComplexityClass{NP-Complete}%
   \index{NP!complete}\xspace}
\providecommand{\APXHard}{\ComplexityClass{APX-Hard}%
   \index{APX!Hard}\xspace}
\providecommand{\NPC}{\ComplexityClass{NPC}%
   \index{NP!complete}\xspace}
\providecommand{\NP}{\ComplexityClass{NP}%
   \index{NP}\xspace}
\providecommand{\POLYT}{\ComplexityClass{P}\xspace}

\newcommand{\envFrame}[1]{\begin{center}\fbox{\begin{minipage}[c]{0.75\textwidth}{#1}\end{minipage}}\end{center}}

\newcommand{\mysqrt}[1]{\ts\!\!\sqrt{#1}}
\newcommand{\SegSet}{\mathsf{S}}
\newcommand{\ArrVD}[1]{\Arr^{|} \pth{#1}}
\newcommand{\RIC}{\textsf{RIC}\index{random incremental
      construction}\xspace}
\newcommand{\ArrB}{\EuScript{B}}

\newcommand{\TreeOpt}{\EuScript{T}_\mathrm{opt}}
\newcommand{\TreeOptA}{\EuScript{T}_\mathrm{opt}'}

\newcommand{\DepthExt}[2]{\mathrm{depth}_{#1}\pth{#2}}
\newcommand{\depthX}[1]{\mathrm{depth}\pth{#1}}

\newcommand{\TreeA}{\EuScript{U}}
\newcommand{\TreeB}{\EuScript{V}}

\newcommand{\Trap}{\sigma}
\newcommand{\TrapA}{\sigma'}
\newcommand{\TrapB}{\sigma''}

\newcommand{\SampleA}{\mathsf{R}'}

\newcommand{\AllRegions}{\mathcal{T}}
\newcommand{\RegionsSet}[1]{\mathcal{F}\pth{#1}}
\newcommand{\RegionsHeavy}[2]{\mathcal{F}_{\!\geq#1}\pth{#2}}
\newcommand{\NumOfRegionsZero}[1]{\ExChar \!f\pth{#1}}
\newcommand{\NumOfRegions}[2]{\ExChar\! f_{\!\geq #1}\pth[]{#2}}
\newcommand{\ObjSet}{\mathsf{S}}

\newcommand{\WeightX}[2][\!]{\omega \pth[#1]{#2}}
\newcommand{\Edge}{\mathsf{e}}

\newcommand{\Face}{\mathsf{f}}

\newcommand{\DefSet}[1]{D\pth{#1}}
\newcommand{\KillSet}[1]{K\pth{#1}}

\newcommand{\CSProbFull}[4]{
   \,\rho_{\!\!\stackrel{\rule[-.005cm]{0cm}{0.1cm}}{#1,
       #4}}\pth{\MakeSBig #2,#3}}

\newcommand{\disk}{D}
\newcommand{\diskA}{D'}

\newcommand{\IntervalSet}{\widehat{\mathcal{I}}}

\newcommand{\GroundSet}{\textsf{X}}
\newcommand{\FGroundSet}{\mathsf{x}}
\newcommand{\RangeSet}{{\mathcal{R}}}
\newcommand{\RangeSetA}{\mathcal{R}'}
\newcommand{\MeasureChar}{\overline{m}}
\newcommand{\Measure}[1]{\MeasureChar\pth{#1}}
\newcommand{\sMeasure}[1]{\overline{s}\pth{#1}}

\newcommand{\SetA}{S}%
\newcommand{\SetB}{T}%
\newcommand{\SetC}{N}%

\newcommand{\VCProj}[2]{#1_{|#2}}
\newcommand{\VCDim}{\mathrm{dim}_{\VC}}
\newcommand{\Dim}{\delta}
\newcommand{\DimA}{\delta'}
\newcommand{\PntSetB}{{{\mathsf{Q}}}}
\newcommand{\PntSetC}{{{\mathsf{R}}}}
\newcommand{\Growth}[2]{{\EuScript{G}}_{#1}\pth{#2}}
\newcommand{\ShatterFunc}[1]{\pi_{#1}}
\newcommand{\Distribution}{\EuScript{D}}
\newcommand{\Error}{\EuScript{E}}

\newcommand{\bd}{{\partial}}%
\newcommand{\eps}{{\varepsilon}}%

\newcommand{\clmlab}[1]{\label{claim:#1}}
\newcommand{\clmref}[1]{\HLink{Claim}{claim:#1}}

\newcommand{\CHX}[1]{{\mathcal{CH}}\pth{#1}}

\newcommand{\ts}{\hspace{0.6pt}}

\def\qedsymbol{\myqedsymbol}%

\newcommand{\qedsign}{\hfill{\hfill\myqedsymbol}}
\newcommand{\qedsignBig}{\hfill{\hfill\rule{4mm}{4mm}}}
\newcommand{\QED}{{\hfill{\hfill\qedsymbol}}}

\newcommand{\num}{\alpha}

\newcommand{\NEW}[1]{#1}

\newcommand{\ThreeFiguresExt}[9]{%
   \begin{tabular}{ccc}
     \begin{minipage}{0.31\linewidth}
         \begin{center}
             {\IncludeGraphics[#1]{#2}}
         \end{center}
     \end{minipage}
     &
       \begin{minipage}{0.31\linewidth}
           \centerline{\IncludeGraphics[#4]{#5}}
       \end{minipage}
     &
       \begin{minipage}{0.31\linewidth}
           \centerline{\IncludeGraphics[#7]{#8}}
       \end{minipage}
     \\
     ~\\[-0.2cm]
     {#3} & {#6} & {#9}
   \end{tabular}}

\newcommand{\IncludeGraphics}[2][]{%
   \typeout{}%
   \typeout{Graphics: #2}%
   \typeout{\ includegraphics[#1]{#2}}%
   \includegraphics[#1]{#2}
   \typeout{}%
}

%    -----------------------------------------------------------------
%    -----------------------------------------------------------------
%    -----------------------------------------------------------------
%    Handling references
%    -----------------------------------------------------------------
%    -----------------------------------------------------------------
%    -----------------------------------------------------------------
\newcommand{\HLinkShort}[2]{\hyperref[#2]{#1\ref*{#2}}}
\newcommand{\HLink}[2]{\hyperref[#2]{#1~\ref*{#2}}}
\newcommand{\HLinkPage}[2]{\hyperref[#2]{#1~\ref*{#2}%
      $_\text{p\pageref{#2}}$}}

\newcommand{\HLinkSuffix}[3]{\hyperref[#2]{#1\ref*{#2}{#3}}}
\newcommand{\HLinkPageSuffix}[3]{\hyperref[#2]{#1\ref*{#2}%
      #3$_\text{p\pageref{#2}}$}}

\newcommand{\figlab}[1]{\label{fig:#1}}
\newcommand{\figref}[1]{\HLink{Figure}{fig:#1}}

\newcommand{\exelab}[1]{\label{exer:#1}}
\newcommand{\exeref}[1]{\HLink{Exercise}{exer:#1}}

\providecommand{\deflab}[1]{\label{def:#1}}
\newcommand{\defref}[1]{\HLink{Definition}{def:#1}}
\newcommand{\defrefY}[2]{\hyperref[def:#2]{#1}}
\newcommand{\defrefpage}[1]{\HLinkPage{Definition}{def:#1}}

\newcommand{\etal}{\textit{et~al.}\xspace}

\newcommand{\obslab}[1]{\label{observation:#1}}
\newcommand{\obsref}[1]{\HLink{Observation}{observation:#1}}
\newcommand{\obsrefY}[2]{\hyperref[observation:#2]{#1}}

\newcommand{\exmlab}[1]{\label{example:#1}}
\newcommand{\exmref}[1]{\HLink{Example}{example:#1}}

\newcommand{\remlab}[1]{\label{rem:#1}}
\newcommand{\remref}[1]{\HLink{Remark}{rem:#1}}%

\newcommand{\chaplab}[1]{\label{chap:#1}}

\newcommand{\itemlab}[1]{\label{item:#1}}
\newcommand{\itemref}[1]{\HLinkSuffix{}{item:#1}{}}

\newcommand{\tedlab}[1]{\label{tedium:#1}}
\newcommand{\tedref}[1]{\HLink{Tedium}{tedium:#1}}%

\newcommand{\factlab}[1]{\label{fact:#1}}%

\newcommand{\alglab}[1]{\label{algorithm:#1}}

\newcommand{\net}{\mathcal{N}}%


\newcommand{\lemlab}[1]{\label{lemma:#1}}
\newcommand{\lemref}[1]{\HLink{Lemma}{lemma:#1}}%
\newcommand{\lemrefY}[2]{\hyperref[lemma:#2]{#1}}
\newcommand{\lemrefpage}[1]{\HLinkPage{Lemma}{lemma:#1}}
\newcommand{\lemrefshort}[1]{\HLinkShort{L}{lemma:#1}}

\newcommand{\problab}[1]{\label{prob:#1}}%
\newcommand{\probref}[1]{\HLink{Problem}{prob:#1}}%

\newcommand{\eqlab}[1]{\label{equation:#1}}
\newcommand{\Eqrefbare}[1]{\HLinkSuffix{(}{equation:#1}{)}}
\newcommand{\Eqref}[1]{\HLinkSuffix{Eq.~(}{equation:#1}{)}}
\newcommand{\Eqrefpage}[1]{\HLinkPageSuffix{Eq.~(}{equation:#1}{)}}

\newcommand{\seclab}[1]{\label{sec:#1}}
\newcommand{\secref}[1]{\HLink{Section}{sec:#1}}
\newcommand{\secrefY}[2]{\hyperref[sec:#2]{#1}}
\newcommand{\secrefpage}[1]{\HLinkPage{Section}{sec:#1}}

\newcommand{\thmlab}[1]{{\label{theo:#1}}}
\newcommand{\thmref}[1]{\HLink{Theorem}{theo:#1}}
\newcommand{\thmrefY}[2]{\hyperref[theo:#2]{#1}}
\newcommand{\thmrefpage}[1]{\HLinkPage{Theorem}{theo:#1}}

\newcommand{\corlab}[1]{\label{cor:#1}}
\newcommand{\corref}[1]{\HLink{Corollary}{cor:#1}}%
\newcommand{\correfY}[2]{\hyperref[cor:#2]{#1}}

\newcommand{\Yend}{Y_{\mathrm{end}}}%


%    -----------------------------------------------------------------
%    -----------------------------------------------------------------
%    -----------------------------------------------------------------
%    Handling references
%    -----------------------------------------------------------------
%    -----------------------------------------------------------------
%    -----------------------------------------------------------------

\newcommand{\ceil}[1]{\left\lceil {#1} \right\rceil}
\newcommand{\floor}[1]{\left\lfloor {#1} \right\rfloor}
\newcommand{\brc}[1]{\left\{ {#1} \right\}}
\newcommand{\pth}[2][\!]{\mleft({#2}\mright)}%
\newcommand{\brk}[1]{\mleft[{#1}\mright]}%
\newcommand{\pbrc}[2][\!\!]{#1\left[ {#2} \bigr. \right]}
\newcommand{\pbrcS}[1]{\left[ {#1} \right]}

\newcommand{\dotprod}[2]{{#1}\cdot{#2}}
\newcommand{\DotProd}[2]{\permut{{#1},{#2}}}
\newcommand{\permut}[1]{\left\langle {#1} \right\rangle}

\newcommand{\MakeBig}{}
\newcommand{\MakesBig}{}
\newcommand{\MakeSBig}{}
\newcommand{\MakeVBig}{}

\newcommand{\DistChar}{\mathsf{d}}%
\newcommand{\DistSetY}[2]{\DistChar\pth{#1,#2}}
\newcommand{\distG}[2]{\DistChar_\Graph\pth{#1, #2}}
\newcommand{\DistMtrCharX}[1]{\DistChar_{#1}}

\newcommand{\dY}[2]{\DistChar\pth{#1, #2}}%
\newcommand{\dGY}[2]{\DistChar_\Graph\pth{#1, #2}}%

\newcommand{\cset}{\mathcal{C}}%

\newcommand{\MTR}{\mathcal{M}}
\newcommand{\DistMtr}[2]{\DistChar\pth{#1,#2}}

\newcommand{\dist}[1]{\left\| {#1} \right\|}
\newcommand{\norm}[1]{\left\| {#1} \right\|}

\providecommand{\ds}{\displaystyle}
\newcommand{\sep}[1]{\,\left|\, {#1} \bigr.\right.}%
\newcommand{\sepw}[1]{\left|\, {#1} \right.}%

\newcommand{\ProbExt}[2]{\mathop{\mathbf{Pr}}_{#1}\pbrcx{#2}}
\newcommand{\ProbLTR}{\mathbb{P}}%

\newcommand{\Prob}[1]{\mathop{\ProbLTR} \mleft[ #1 \mright]}%
\newcommand{\ProbCond}[2]{\mathop{\ProbLTR}\!\left[%
       #1 \;\middle\vert\; #2 \right]}
\newcommand{\ProbY}[2]{\ProbLTR^{}_{\!#1}\pbrc{#2}}

\newcommand{\ExChar}{\mathbb{E}}%
\newcommand{\ExSym}{\mathop{\ExChar}}%
\newcommand{\Ex}[2][\!]{\ExSym#1\pbrcx{#2}}
\newcommand{\ExExt}[2]{\ExSym_{#1}\!\pbrcx{#2}}
\newcommand{\ExExtS}[2]{\ExSym_{\substack{#1}}\!\pbrcx{#2}}
\newcommand{\ExCond}[2]{\ExSym\!\left[%
       #1 \;\middle\vert\; #2 \right]}

\newcommand{\VarSymb}{\mathop{\mathbb{V}}}
\newcommand{\Var}[1]{\mathop{\mathbb{V}}\!\pbrcx{#1}}

\newcommand{\pbrcx}[1]{\left[ {#1} \right]}%
\newcommand{\nfrac}[2]{{#1}/{#2}}
\newcommand{\cardin}[1]{\left| {#1} \right|}%

\newcommand\CCancel[2][black]{{\renewcommand\CancelColor{\color{#1}}\cancel{#2}}}

\newcommand{\IncGraphPage}[4][]{%
   \IncludeGraphics[page=#4,#1]{\File{#2/#3}}%
}

\newenvironment{enumEx}{
   \begin{compactenumA}
    }{\end{compactenumA}}

\newcommand{\cl}{{\mathrm{cl}}}

\usepackage[amsmath,thmmarks]{ntheorem}
\theoremseparator{.}%

\newtheorem{theorem:unnumbered}{Theorem}[section]
\newtheorem{theorem}{Theorem}[section]

\newtheorem{observation}[theorem]{Observation}
\newtheorem{conjecture}[theorem]{Conjecture}
\newtheorem{lemma}[theorem]{Lemma}
\newtheorem{corollary}[theorem]{Corollary}
\newtheorem{claim}[theorem]{Claim}
\newtheorem{assumption}[theorem]{Assumption}
\newtheorem{tedium}[theorem]{Tedium}

\theoremstyle{plain}%
\theoremheaderfont{\sffamily}
\theorembodyfont{\upshape}%

\newtheorem{defn}[theorem]{Definition}
\newtheorem{definition}[theorem]{Definition}

\newtheorem*{remark_unnumbered}[theorem:unnumbered]{Remark}%
\newtheorem*{remarks}[theorem]{Remarks}%
\newtheorem{remark}[theorem]{Remark}%
\newtheorem{fact}[theorem]{Fact}%
\newtheorem{example}[theorem]{Example}
\newtheorem{exercise}[theorem]{Exercise}
\newtheorem{problem}[theorem]{Problem}
\newtheorem{algorithm}[theorem]{Algorithm}%
\newtheorem{xca}[theorem]{Exercise}
\newtheorem{openproblem}[theorem]{Open Problem}
\newtheorem{exercise_h}[theorem]{Exercise}

% Proof environments
\theoremheaderfont{\em}%
\theorembodyfont{\upshape}%
\theoremstyle{nonumberplain}
\theoremseparator{}
\theoremsymbol{\myqedsymbol}
\newtheorem{proof}{Proof:}
\newtheorem{proofextin}{Proof}
\newtheorem{altproof}{Alternative proof:}

\theoremsymbol{}
\newtheorem{proofnoqed}{Proof:}

\theoremsymbol{\myqedsymbol}
\theorempreskip{0pt}%
\theorempostskip{0pt}%
\newtheorem{proofof}{Proof of\!}%

\newenvironment{proofExt}[1]{{\emph{#1:}}}{\hfill{\hfill\myqedsymbol}}
\newenvironment{proofnoproof}{{}}{{\hfill{\hfill\myqedsymbol}}}

%%%%%%%%%%%%%%%%%%%%%%%%%%%%%%%%%%%%%%%%%%%%%%%%%%%%%%%%%%%%%%%%%%%%
%%%%%%%%%%%%%%%%%%%%%%%%%%%%%%%%%%%%%%%%%%%%%%%%%%%%%%%%%%%%%%%%%%%%
% Macros to embed figures in text

\newcommand{\EBRY}[2]{#1_{\langle#2\rangle}}%

\newcommand{\NiceFrac}[2]{%
   {\raisebox{0.2em}{\ensuremath{#1}}\!\!\!}%
   {\rotatebox[origin=c]{-20}{\Big/}}% Rotated /
   {\! \raisebox{-0.2em}{\ensuremath{#2}} }}%

% end - Macros to embed figures in text
%%%%%%%%%%%%%%%%%%%%%%%%%%%%%%%%%%%%%%%%%%%%%%%%%%%%%%%%%%%%%%%%%%%% 
%%%%%%%%%%%%%%%%%%%%%%%%%%%%%%%%%%%%%%%%%%%%%%%%%%%%%%%%%%%%%%%%%%%%

%%%%%%%%%%%%%%%%%%%%%%%%%%%%%%%%%%%%%%%%%%%%%%%%%%%%%%%%%%%%%%%%%%
% Color in math
%\providecommand{\Mh}[1]{{#1}}%
\providecommand{\IfColorMath}[2]{#2}%

\IfFileExists{.latex_printer_friendly}{\def\GenPrinterVer{1}}{}%

\ifx\GenPrinterVer\undefined
   \IfFileExists{.latex_color}{\def\GenColorMath{1}}{}
\else
   \renewcommand{\IfPrinterVer}[2]{#1}%
\fi

\ifx\GenColorMath\undefined
\else
\renewcommand{\IfColorMath}[2]{#1}%
%\renewcommand{\Mh}[1]{{\textcolor{red}{#1}}}%
\fi
% Color in math - end
%%%%%%%%%%%%%%%%%%%%%%%%%%%%%%%%%%%%%%%%%%%%%%%%%%%%%%%%%%%%%%%%%%

\newcommand{\HSetY}[2]{H\pth{#1, #2}}%
\newcommand{\CHSetY}[2]{C\pth{#1, #2}}%
\newcommand{\Holder}{H\"older\xspace}
\newcommand{\OldOm}{{\mathcal{O}}}
\newcommand{\NewOm}{{\mathcal{N}}}

%%%%%%%%%%%%%%%%%%%%%%%%%%%%%%%%%%%%%%%%%%%%%%%%%%%%%%%%%%%%%%%%%%%%%%%
%---------------------------------------------------------------
\newcommand{\IntRange}[1]{\left\llbracket #1 \right\rrbracket}
\newcommand{\IRX}[1]{\left\llbracket #1 \right\rrbracket}%
\newcommand{\IRY}[2]{\left\llbracket #1:#2 \right\rrbracket}

%---------------------------------------------------------------
%%%%%%%%%%%%%%%%%%%%%%%%%%%%%%%%%%%%%%%%%%%%%%%%%%%%%%%%%%%%%%%%%%%%%%%

\newcommand{\obj}{f}%
\newcommand{\Wg}{\omega}%
\newcommand{\TotalW}{W}%

\newcommand{\rankY}[2]{{#1}_{ \permut{#2}}}%%
\newcommand{\areaX}[1]{\mathrm{area}\pth{#1}}%%

\newcommand{\mybinom}[2]%
{\Bigl(\begin{array}{@{}c@{}}#1\\#2\end{array}\Bigr)}

\providecommand{\Set}[2]{\left\{ #1 \;\middle\vert\; #2 \right\}}
\newcommand*{\biga}[1]{\vcenter{\hbox{\scalebox{1.5}{\ensuremath#1}}}}
\newcommand{\CNF}{\TermI{CNF}\xspace}%
\newcommand{\Tree}{T}%

\DeclareSymbolFont{sfletters}{OML}{cmbrm}{m}{it}

\DeclareMathSymbol{\salpha}{\mathord}{sfletters}{"0B}

%%%%%%%%%%%%%%%%%%%%%%%%%%%%%%%%%%%%%%%%%%%%%%%%%%%%%%%%%%%%%%%%%%%%

\providecommand{\ProblemInd}[2]{\textsf{\textcolor{red}{#1}\index{Problem!#2}}}
\newcommand{\NPProblemInd}[4]{
   \begin{center}
      \vspace{0.6cm}
       \fbox{
          \begin{minipage}{0.9\linewidth}
              \vspace{-0.6cm} \hspace{-0.35cm}
              {\large \bfseries \ProblemInd{#1}{#2}}
              \vspace{+0.2cm}%

              \noindent{\textcolor{blue}{\bfseries \textsf{Instance}}:} {#3}
              
              \noindent{\textcolor{blue}{\bfseries \textsf{Question}:}} {#4}
          \end{minipage}
       }
   \end{center}
   \smallskip
   }

\newcommand{\NPProblemOpt}[4]{\NPProblemInd{#1}{#2}{#3}{#4}}

%%%%%%%%%%%%%%%%%%%%%%%%%%%%%%%%%%%%%%%%%%%%%%%%%%%%%%%%%%%%%%%%%%%%

\newcommand{\Hastad}{H{\aa}stad\xspace}

\newcommand{\ANNbr}{\Term{ANN}\index{approximate nearest neighbor}\xspace}
\newcommand{\LSH}{\TermI{L{S}H}\xspace}

\newcommand{\DS}{\EuScript{D}} % DataStructure
\newcommand{\DSC}{\mathsf{D}} % DataStructure

\newcommand{\distH}[2]{\DistChar_{H}\pth{#1, #2}}
\newcommand{\distHY}[2]{\DistChar_{H}\pth{#1, #2}}

\newcommand{\DSAprxNearNbr}{\ensuremath{\DataStructure{{\approx%
         \text{Near{}}}}}%
   \index{near neighbor!data-structure!approximate}%
}
\newcommand{\DataStructure}[1]{\ensuremath{\DS_{#1}}}
\newcommand{\DistNN}[2]{\DistChar\pth{#1,#2}}
\newcommand{\pntc}{{p}}%

\newcommand{\seq}{{M}}%
\newcommand{\cTimes}{\beta}%

\newcommand{\DSTimes}{L}%
\newcommand{\rad}{r}%

\newcommand{\HypercubeX}[1]{\Hypercube^{#1}}
\newcommand{\kFunctions}[2]{\mathrm{combine}\pth{#1,#2}}
\newcommand{\WEqualChar}{\Bumpeq}
\newcommand{\WEqual}[2]{{#1}\WEqualChar{#2}}
\newcommand{\FamilyAmp}{\FamilyC_{\WEqualChar}}
\newcommand{\PSPACE}{\ComplexityClass{PSPACE}\xspace}%

\newcommand{\PS}{P}%

\newcommand{\PSetFar}{\ensuremath{\PntSet^{}_{\geq}\xspace}}%
\newcommand{\DA}{D}%
\newcommand{\uniqX}[1]{\mathrm{uniq}\pth{#1}}%
\newcommand{\diffX}[1]{\mathop{}\!\mathrm{d}#1}

\newcommand{\Frq}{F}%
\newcommand{\frq}{f}%
\newcommand{\STRM}{\mathcal{S}}%
\newcommand{\strm}{{s}}%

\providecommand{\Turan}{Tur\'an\xspace}
\newcommand{\fGrowX}[1]{G^{}_{\Dim}\pth{#1}}
\newcommand{\CompressFunc}{\mathtt{Compress}}
\newcommand{\RSA}{\Term{RSA}\xspace}
\newcommand{\MST}{\Term{MST}\xspace}
\providecommand{\CClassPP}{\CClass{P{}P}\xspace}
\providecommand{\ID}{\mathrm{id}}

\newcommand{\EncodeFunc}{\text{\ttfamily{Enc}}}
\newcommand{\DecodeFunc}{\text{\ttfamily{Dec}}}
\newcommand{\discX}[1]{\mathrm{disc}\pth{#1}}%

\newcommand{\propertyC}{\psi}%
\newcommand{\propertyX}[1]{\propertyC\pth{#1}}%
\newcommand{\numC}{\Gamma}%

\newcommand{\degreeX}[1]{\mathrm{degree}\pth{#1}}

\newcommand{\ind}{m}

\providecommand{\nequiv}{\not\equiv}
\providecommand{\sqcupplus}{\sqcup^+}

\newcommand{\AlgorithmI}[1]{{\protect\color{RedViolet}{#1}}}

\usepackage{marginnote}
\usepackage{ragged2e} 
% Short, subtle snark macro
\setlength{\marginparwidth}{1.5in}
\setlength{\marginparsep}{10pt}
\newcommand{\snark}[1]{%
   \marginnote{\raggedright\scriptsize\sffamily\color{gray}#1}%
}

\newcommand{\binsY}[2]{\sqcup_{\geq #1}\pth{#2}}
\newcommand{\ballsY}[2]{\text{\Football}_{\geq #1}\pth{#2}}

\newcommand{\BLX}[1]{{{\sqcupplus{}}_{\geq #1}}}

\MakePerPage{footnote} % Resets footnote counter on every page

\allowdisplaybreaks % Allows page breaks inside multiline math environments

\let\Mh\undefined


% 1. Create custom page style for chapter start pages
\makepagestyle{chapfirstpage}

\makefootrule{chapfirstpage}{\textwidth}{0.4pt}{1ex}

% 2. Define footers: {Left}{Center}{Right}
% For two-sided documents (recto/verso):
\makeoddfoot{chapfirstpage}{\footnotesize Sariel Har-Peled}{\thepage}{\today}
\makeevenfoot{chapfirstpage}{\thepage}{Sariel Har-Peled}{}

% (Leave headers empty on chapter start pages)
\makeoddhead{chapfirstpage}{}{}{}
\makeevenhead{chapfirstpage}{}{}{}

% 3. Alias 'chapter' style to use your new style
\aliaspagestyle{chapter}{chapfirstpage}

\newcommand{\pdf}{\textsf{pdf}\xspace}
\newcommand{\cdf}{\textsf{cdf}\xspace}

% -----------------------------------------------------------------------------
% Bibliography & Document Begin
%-----------------------------------------------------------------------------
\IfFileExists{\jobname.bib}{%
   \addbibresource{\jobname.bib}%
}{%
   \addbibresource{shortcuts.bib}%
   \addbibresource{geometry.bib}%
}
\makeindex
\begin{document}
}{\input{prefix_latex}}%
  \fi
\fi
\makeatother
%

       
\BookMode{\newpage}%

\noindent%
\ChapterNumPage{29}{230}
\Chapter{Approximating the Number of Distinct Elements in a Stream}

\QuotePExt{%
   ``See? Genuine-sounding indignation. I programmed that myself. It's
   the first thing you need in a university environment: the ability
   to take offense at any slight, real or imagined.''%
}{Robert Sawyer}{Factoring Humanity}


\noindent%
\ChapterPrefix{}%
\ChapterGen{}

~

% \bigskip%

\bigskip%



%%%%%%%%%%%%%%%%%%%%%%%%%%%%%%%%%%%%%%%%%%%%%%%%%%%%%%%%%%%%%%%%%%%%%%%%%%
%%%%%%%%%%%%%%%%%%%%%%%%%%%%%%%%%%%%%%%%%%%%%%%%%%%%%%%%%%%%%%%%%%%%%%%%%% 

\section{Counting number of distinct elements}

We first analyze an idealized estimator based on fully random hash
values. We then turn to pairwise independent sampling.

\subsection{First order statistic}

Let $X_1, \ldots, X_u$ be $u$ independent random variables uniformly
distributed in $[0,1]$. Let $Y = \min (X_1, \ldots, X_u)$. The value $Y$ is
the \emphOnly{first order statistic}\index{first order statistic} of
$X_1,\ldots, X_u$.


%XXX
For a continuous variable $X$, the \emphOnly{probability density
   function}\index{probability density function} (i.e.,
\emphOnly{\pdf}\index{\pdf}) is the ``probability'' of $X$ having this
value. Since this is not well defined, one looks on the
\emphOnly{cumulative distribution function}\index{cumulative distribution
   function} $F(x) = \Prob{X \leq x}$. The \pdf is then the derivative of
the \cdf.  Discrete distribution do not have ordinary \pdf{}s. This can
be simulated by Dirac's $\delta$-functions, but fortunately, this pile of
amusement and headache is outside our scope. In particular, the
probability $\Prob{X=\alpha}$, for any specific $\alpha$, is zero.

\begin{lemma}
    The \emph{probability density function} of $Y$ is
    \begin{equation*}
        f(x) = \binom{u}{1} 1 (1-x)^{u-1}.
    \end{equation*}
\end{lemma}
\begin{proof}
    The minimum exceed $x$ only if all the $u$ numbers are larger than
    $x$. Thus,
    $ F(x)=\Prob{Y \leq x } =1 - (1-x)^u$. As such
    $f(x) = F'(x) = u(1-x)^{u-1}$.
\end{proof}
\begin{proof}[Second version: no need to read it.]
    Considering the \pdf of $X_1$ being $x$, and all other $X_i$s
    being bigger. We have that this \pdf is
    \begin{align*}
        g(x)
        &=%
        \Prob{ (X_1 = x) \cap \smash{\bigcap_{i=2}^u} (X_i>X_1)\Bigr.}
        \\&=
        \ProbCond{\Bigl. \smash{\bigcap_{i=2}^u} (X_i>X_1)}{X_1=x} \Prob{X_1=x}
        \\&=%
        (1-x)^{u-1}.
    \end{align*}
    Since every one of the $X_i$ has equal probability to realize $Y$,
    we have $f(x) = u g(x)$.
\end{proof}

\begin{lemma}
    \lemlab{order:stat}%
    %
    We have $\Ex{Y} = \frac{1}{u+1}$,
    $\Ex{Y^2} = \frac{2}{(u+1)(u+2)}$, and
    $\Var{Y}=\frac{u}{(u+1)^2 (u+2)}$.
\end{lemma}
\begin{proof}
    Let $A,B$ be two independent uniform samples from $[0,1]$,
    independent of $X_1,\ldots,X_u$. Conditioned on $Y$, the probability
    that $A<Y$ is $Y$. On the other hand, $A<Y$ exactly when $A$ is
    the smallest of $u+1$ samples. Thus, by symmetry,
    \[
        \Ex{Y} = \Prob{A<Y} = \frac{1}{u+1}.
    \]
    Similarly, conditioned on $Y$, the probability that both $A$ and
    $B$ are smaller than $Y$ is $Y^2$. This happens exactly when
    $A,B$ are the two smallest of $u+2$ samples. All pairs are equally
    likely to be the two smallest, so
    \[
        \Ex{Y^2} = \Prob{A<Y,\ B<Y}
        = \frac{1}{\binom{u+2}{2}}
        = \frac{2}{(u+1)(u+2)}.
    \]
    Consequently,
    % BEGIN EDIT W5: Split the variance calculation to fit the line.
    \begingroup\color{blue}
    \begin{align*}
        \Var{Y} &= \Ex{Y^2} - \pth{\Ex{Y}}^2\\
        &= \frac{2}{(u+1)(u+2)} - \frac{1}{(u+1)^2}
        = \frac{u}{(u+1)^2(u+2)}.
    \end{align*}
    \endgroup
    % END EDIT W5
\end{proof}

\subsection{The algorithm}

\paragraph*{A single estimator.}



Assume that we have a perfectly random hash function $h$ that randomly
maps $N = \brc{1,\ldots, n}$ to $[0,1]$. Assume that the stream
$\STRM=\strm_1,\ldots,\strm_m$, has $u$ unique numbers in $N$. Then the set
$\{h(\strm_1),\ldots, h(\strm_m)\}$ contains $u$ random numbers uniformly
distributed in $[0,1]$.  The algorithm computes $Y = \min_i h(\strm_i)$.

\paragraph*{Explanation.}
Note that $Y$ is \emph{not} an estimator for $u$ -- instead, as
$\Ex{Y} = 1/(u+1)$, we are estimating $1/(u+1)$. The key observation is
that an $1 \pm \eps$ estimator for $1/(u+1)$, is $1 \pm O(\eps)$ estimator for
$u+1$, which is in turn an $1 \pm O(\eps)$ estimator for $u$.

\begin{lemma}
    Let $\eps, \BadProb \in (0,1)$ be parameters. Given a stream
    $\STRM$ of items from $\{1,\ldots, n\}$ one can return an estimate
    $X$, such that
    \begin{equation}
        \Prob{ (1-\eps/4)\frac{1}{u+1} \leq X \leq
           (1+\eps/4)\frac{1}{u+1}} \geq 1-\BadProb,
        \eqlab{range-est}
    \end{equation}
    where $u$ is the
    number of unique elements in $\STRM$. This requires
    $O\pth{ \frac{1 }{\eps^2 } \log \frac{1}{\BadProb}}$ space.
\end{lemma}
\begin{proof}
    The basic estimator $Y$ has $\mu = \Ex{Y} = \frac{1}{u+1}$ and
    $\nu = \Var{Y}=\frac{u}{(u+1)^2 (u+2)}$. We now plug this
    estimator into the mean/median framework. By \lemref{order:stat},
    for $c$ some absolute constant, this requires maintaining $M$
    estimators, where $M$ is larger than
    \begin{equation*}
        c \frac{4\cdot 16 \nu}{\eps^2 \mu^2 } \log \frac{1}{\BadProb}
        = 
        O\pth{  \frac{u^2 }{\eps^2 u^2 } \log \frac{1}{\BadProb}}
        = 
        O\pth{  \frac{1 }{\eps^2  } \log \frac{1}{\BadProb}}.
    \end{equation*}
\end{proof}

Assume $X$ is in the desired range, and write it in the form
$X = (1+\eta)/(u+1)$, where $|\eta|\leq\eps/4$. For a nonempty stream, $u\geq1$, so
\begin{align*}
    \frac{|(1/X-1)-u|}{u}
    &= \frac{u+1}{u}\frac{|\eta|}{|1+\eta|}
    \leq \frac{2|\eta|}{1-|\eta|}
    \leq \frac{2\eps}{4-\eps}<\eps.
\end{align*}
Thus, $1/X-1$ is a $(1\pm\eps)$-estimate of the number of distinct
elements.

\paragraph*{The algorithm revisited.}
Compute $X$ as above, and output the quantity $1/X-1$.

\bigskip

This immediately implies the following.

\begin{lemma}
    Let $\eps,\BadProb\in(0,1)$. Under the unreasonable\footnote{Criminal
       even!} assumption that we can sample perfectly random functions
    from $\{1,\ldots,n\}$ to $[0,1]$, and store each function and each hash
    value in $O(1)$ words, one can process a nonempty stream with $u$
    distinct elements and return an estimate $\rho$ satisfying
    \[
        \Prob{(1-\eps)u\leq\rho\leq(1+\eps)u}\geq1-\BadProb,
    \]
    using $O(\eps^{-2}\log(2/\BadProb))$ words.
\end{lemma}


%%%%%%%%%%%%%%%%%%%%%%%%%%%%%%%%%%%%%%%%%%%%%%%%%%%%%%%%%%%%%%%%%%%%%%%%%%%%%% 
%%%%%%%%%%%%%%%%%%%%%%%%%%%%%%%%%%%%%%%%%%%%%%%%%%%%%%%%%%%%%%%%%%%%%%%%%%%%%% 
%%%%%%%%%%%%%%%%%%%%%%%%%%%%%%%%%%%%%%%%%%%%%%%%%%%%%%%%%%%%%%%%%%%%%%%%%%%%%%
\newcommand{\distinctX}[1]{\mathrm{set}\pth{#1}}

% BEGIN EDIT W5: Shorten the overfull heading and name the assumption.
\begingroup\color{blue}
\section{Sampling with pairwise independent randomness}
\endgroup
% END EDIT W5

Assume that we have a stream of elements
$\STRM = \strm_1, \ldots, \strm_m$, all taken from the set
$\{1,\ldots, n\}$.  Let $\distinctX{\STRM}$ denote the set of values that
appear in $\STRM$, and set $S=\distinctX{\STRM}$. Thus,
\begin{equation*}
    \Frq_0 = \Frq_0(\STRM) = \cardin{S}
\end{equation*}
is the number of distinct values in the stream $\STRM$.


Assume that we have a random sequence of bits
$\mathcal{B} \equiv B_1, \ldots, B_n$, such that $\Prob{B_i=1} =p$, for some
$p$. Furthermore, we can compute $B_i$ efficiently. Assume that the bits
of $\mathcal{B}$ are pairwise independent, and can be computed using
constant space.


\paragraph*{The sampling algorithm.}
When the $i$\th item $\strm_i$ arrives, we compute $B_{\strm_i}$. If this
bit is $1$, then we insert $\strm_i$ into the random sample $R$ (if it is
already in $R$, there is no need to store a second copy, naturally).


This defines a natural random sample
\begin{equation*}
    R = \Set{i}{B_i =1 \text{ and } i \in S} \subseteq S.
\end{equation*}

\begin{lemma}
    \lemlab{f:e:2:ev}%
    %
    For the above random sample $R$, let $X = \cardin{R}$. We have
    that $\Ex{X} = p \nu$ and $\Var{X} = p \nu - p^2 \nu$, where
    $\nu = \Frq_0(\STRM)$ is the number of distinct elements in $S$.
\end{lemma}
\begin{proof}
    Since $X=\sum_{i\in S}B_i$, linearity of expectation gives
    \[
        \Ex{X}=\sum_{i\in S}\Ex{B_i}=p\nu.
    \]
    Pairwise independence makes the covariances zero, so
    \[
        \Var{X}=\sum_{i\in S}\Var{B_i}
        =\nu p(1-p)=p\nu-p^2\nu.
    \]
\end{proof}
  
\begin{lemma}
    \lemlab{estimator}%
    %
    Let $\eps \in (0,1/4)$.  Given $O(1/\eps^2)$ space, and a
    parameter $N$. Consider the task of estimating the size of
    $\Frq_0 = \cardin{\distinctX{\STRM}}$, where $\Frq_0 > N/4$.
    Then, the algorithm described below outputs one of the following:
    \begin{compactenumA}
        \smallskip%
        \item $\Frq_0 > 2N$.
        \smallskip%
        \item Output a number $\rho$ such that
        $(1-\eps)\Frq_0 \leq \rho \leq (1+\eps)\Frq_0$.
    \end{compactenumA}
    \smallskip%
    (Note, that the two options are not disjoint.)  The output of this
    algorithm is correct, with probability $\geq 7/8$.
\end{lemma}
\begin{proof}
    We set $p= \frac{c}{N\eps^2}$, where $c$ is a constant to be determined
    shortly\footnote{If $p>1$ call the probability police and set its
       value to $1$, which makes what follows uninteresting.}. Let
    $T = p N = O(1/\eps^2)$. We sample a random sample $R$ from $S$, by
    scanning the elements of $S$, and adding $i\in S$ to $R$ if $B_i=1$, If
    the random sample is larger than $8T$, at any point, then the
    algorithm outputs that $\cardin{S} > 2N$.

    In all other cases, the algorithm outputs $\cardin{R} /p$ as the
    estimate for the size of $S$, together with $R$.

    To bound the failure probability, consider first the case that
    $N/4 < \cardin{\distinctX{\STRM}}$. In this case, we have by the
    above, that
    \begin{align*}
        \Prob{\cardin{X - \Ex{X}} > \eps \Ex{X}}%
        &\leq%
        \Prob{\cardin{X - \Ex{X}} > \eps
           \frac{\Ex{X}}{\sqrt{\Var{X}}}
           \sqrt{\Var{X}}}%
        \\&\leq%
        \frac{\Var{X}}{\eps^2(\Ex{X})^2}
        \leq
        \frac{1}{8},
    \end{align*}

    if $\frac{\Var{X}}{\eps^2 (\Ex{X})^2} \leq \frac{1}{8}$, For
    $\nu = \Frq_0 \geq N/4$, this happens if
    $\frac{p \nu}{\eps^2 p^2 \nu^2} \leq \frac{1}{8}$. This in turn is
    equivalent to $8/\eps^2 \leq p \nu $. This in turn happens if
    \begin{equation*}
        \frac{c}{N\eps^2} \cdot \frac{N}{4} \geq \frac{8}{\eps^2},
    \end{equation*}
    which implies that this holds for $c=32$.  Namely, the algorithm
    in this case would output a $(1\pm\eps)$-estimate for
    $\cardin{S}$.

    Thus, assume that $X$ is within the ``good'' range, which happens
    with probability $\geq 7/8$. If the sample gets bigger than $8T$, then
    the above readily implies that the size of $S$ is at least
    $(1-\eps)8T/p > 2N$. Namely, the output of the algorithm is correct in
    this case.
\end{proof}
  
\begin{lemma}
    \lemlab{f:e:0:resolution}%
    %
    Let $\eps \in (0,1/4)$ and $\BadProb \in (0,1)$.  Given
    $O(\eps^{-2} \log \BadProb^{-1})$ space, and a parameter $N$, and
    the task is to estimate $\Frq_0$ of $\STRM$, given that
    $\Frq_0 > N/4$.  Then, there is an algorithm that would output one
    of the following:
    \begin{compactenumA}
        \smallskip%
        \item $\Frq_0 > 2N$.
        \smallskip%
        \item Output a number $\rho$ such that
        $(1-\eps)\Frq_0 \leq \rho \leq (1+\eps)\Frq_0$.
    \end{compactenumA}
    \smallskip%
    (Note, that the two options are not disjoint.)  The output of this
    algorithm is correct, with probability $\geq 1- \BadProb$.
\end{lemma}

\begin{proof}
    We run $O( \log \BadProb^{-1})$ independent copies of
    \lemref{estimator}. If half of them returns that $\Frq_0 > 2N$, then
    the algorithm returns that $\Frq_0 > 2N$. Otherwise, the algorithm
    returns the median of the estimates returned, and return it as the
    desired estimate. The correctness readily follows by a repeated
    application of Chernoff's inequality. Indeed, for the result to be
    wrong, at least a quarter of the estimators have to fail, but the
    probability of failure per estimator is at most $1/8$.
\end{proof}


\begin{lemma}
    Let $\eps \in (0,1/4)$.  Given $O(\eps^{-2} \log^2 n)$ space, one
    can read the stream $\STRM$ once, and output a number $\rho$, such
    that $(1-\eps)\Frq_0 \leq \rho \leq (1+\eps)\Frq_0$. The estimate
    is correct with high probability (i.e., $\geq 1-1/n^{O(1)}$).
\end{lemma}
\begin{proof}
    Let $N_i = 2^i$, for $i=1, \ldots, M$, where $M = \ceil{\lg n}$. Run
    $M$ copies of \lemref{f:e:0:resolution}, for each value of $N_i$,
    with $\BadProb = 1/n^{O(1)}$. Let $Y_1, \ldots, Y_M$ be the outputs of
    these algorithms for the stream. A prefix of these outputs, are going
    to be ``$\Frq_0 > 2N_i$'', Let $j$ be the first $Y_j$ that is a
    number.

  
    % Return this number as the desired estimate. The correctness
    % is easy---the first estimate that is a number, is a correct estimate
    % with high probability. Since $N_M \geq n$, it also follows that $Y_M$
    % must be a number. As such, there is a first number in the sequence,
    % and the algorithm would output an estimate.

    More precisely, there is a minimal index $i$, such that
    $N_i/4 \leq \Frq_0 \leq 2N_i$, and $Y_i$ is a good estimate, with high
    probability. The algorithm returns $Y_i$ as the desired estimate.
    Importantly, all earlier estimates are correct with high probability
    as well.
    % If any of the $Y_j$, for $j < i$, is an estimate, then
    % it is correct (again) with high probability.
\end{proof}

\remX{G1: Choose the first scale with $N_i\geq\Frq_0$ and union-bound only the levels up to that scale. The displayed $\Frq_0\leq2\Frq_0$ does not establish this. Handle the empty stream and $n=1$ separately; see review.md.}

%%%%%%%%%%%%%%%%%%%%%%%%%%%%%%%%%%%%%%%%%%%%%%%%%%%%%%%%%%%%%%%%%%%%%%%%%%%%%% 
%%%%%%%%%%%%%%%%%%%%%%%%%%%%%%%%%%%%%%%%%%%%%%%%%%%%%%%%%%%%%%%%%%%%%%%%%%%%%% 
%%%%%%%%%%%%%%%%%%%%%%%%%%%%%%%%%%%%%%%%%%%%%%%%%%%%%%%%%%%%%%%%%%%%%%%%%%%%%%

\section{Bibliographical notes}


%%%%%%%%%%%%%%%%%%%%%%%%%%%%%%%%%%%%%%%%%%%%%%%%%%%%%%%%%%%%%%%%%%%%%%%%%%%%%% 
%%%%%%%%%%%%%%%%%%%%%%%%%%%%%%%%%%%%%%%%%%%%%%%%%%%%%%%%%%%%%%%%%%%%%%%%%%%%%% 
%%%%%%%%%%%%%%%%%%%%%%%%%%%%%%%%%%%%%%%%%%%%%%%%%%%%%%%%%%%%%%%%%%%%%%%%%%%%%%


\StandAloneMode{
   \section{From previous lectures}

   \IncFragment{median:estimator}

   \IncFragment{Chebyshev}

\remX{E4: The imported standardized Chebyshev statement fails when the variance is zero. Require finite positive variance, or use a positive absolute threshold; see review.md. The shared fragment is unchanged.}
   \IncFragment{chernoff_small_delta}
   \IncFragment{k_wise_independent_mod_p}

   
%   \IncFragment{chernoff:small:delta:1}
   % \IncFragment{normal:distribution}
   % \IncFragment{A:N:N:PLEB:Euclidean}
   % \IncFragment{A:N:N:PLEB:slow}
   % \IncFragment{H:S:T:Euclidean}
}


%%%%%%%%%%%%%%%%%%%%%%%%%%%%%%%%%%%%%%%%%%%%%%%%%%%%%%%%%%%%%%%%%%
\nocite{mr-ra-95}

\ChapterEnd{}

% -----------
